\documentclass[11pt,a4paper]{article}

% ============================================================================
% PACKAGES (matching the GAUSE method deep-dive)
% ============================================================================
\usepackage[utf8]{inputenc}
\usepackage[T1]{fontenc}
\usepackage{lmodern}
\usepackage{amsmath,amssymb,amsthm}
\usepackage{mathtools}
\usepackage{bm}
\usepackage{algorithm}
\usepackage{algorithmic}
\usepackage{booktabs}
\usepackage{multirow}
\usepackage{array}
\usepackage{float}
\usepackage{graphicx}
\usepackage{xcolor}
\usepackage[numbers]{natbib}
\usepackage{hyperref}
\usepackage{tikz}
\usetikzlibrary{shapes,arrows,positioning,calc,patterns,decorations.pathreplacing,shapes.geometric,arrows.meta,fit,backgrounds,intersections}
\usepackage{colortbl}
\usepackage{pgfplots}
\pgfplotsset{compat=1.17}
\usepgfplotslibrary{fillbetween}
\usepackage{listings}
\usepackage{tcolorbox}
\usepackage{enumitem}
\usepackage{geometry}
\geometry{margin=1in}

% ============================================================================
% THEOREM ENVIRONMENTS
% ============================================================================
\theoremstyle{definition}
\newtheorem{definition}{Definition}[section]
\newtheorem{example}{Example}[section]
\newtheorem{remark}{Remark}[section]

\theoremstyle{plain}
\newtheorem{theorem}{Theorem}[section]
\newtheorem{proposition}{Proposition}[section]
\newtheorem{lemma}{Lemma}[section]
\newtheorem{corollary}{Corollary}[section]

% ============================================================================
% CUSTOM COMMANDS
% ============================================================================
\newcommand{\R}{\mathbb{R}}
\newcommand{\E}{\mathbb{E}}
\newcommand{\Var}{\mathrm{Var}}
\newcommand{\Cov}{\mathrm{Cov}}
\newcommand{\N}{\mathcal{N}}
\newcommand{\KL}{\mathrm{KL}}
\newcommand{\TV}{\mathrm{TV}}
\newcommand{\nm}{\textsc{RISP}}
\newcommand{\Ginv}{G_{\mathrm{inv}}}
\newcommand{\Gforget}{G_{\mathrm{forget}}}
\newcommand{\argmax}{\operatornamewithlimits{arg\,max}}
\newcommand{\argmin}{\operatornamewithlimits{arg\,min}}

% Code listing style
\lstset{
    language=Python,
    basicstyle=\ttfamily\small,
    keywordstyle=\color{blue},
    commentstyle=\color{gray},
    stringstyle=\color{red},
    numbers=left,
    numberstyle=\tiny\color{gray},
    frame=single,
    breaklines=true,
    captionpos=b
}

% Colored boxes for intuition
\newtcolorbox{intuition}[1][]{
    colback=blue!5!white,
    colframe=blue!75!black,
    fonttitle=\bfseries,
    title=Intuition,
    #1
}

\newtcolorbox{keypoint}[1][]{
    colback=green!5!white,
    colframe=green!50!black,
    fonttitle=\bfseries,
    title=Key Point,
    #1
}

\newtcolorbox{warning}[1][]{
    colback=red!5!white,
    colframe=red!75!black,
    fonttitle=\bfseries,
    title=Important,
    #1
}

% ============================================================================
% DOCUMENT
% ============================================================================
\title{\Huge\textbf{\nm{}: A Mathematical Deep Dive}\\[0.5cm]
\Large Retention and Invariance in Decision-Focused\\
Strategy Pools for Non-Stationary Markets}

\author{Yuhao Li\\
University of Pennsylvania\\
\texttt{li88@sas.upenn.edu}}

\date{June 2026}

\begin{document}

\maketitle

\begin{abstract}
This document is the exhaustive, ground-up mathematical companion to the
\nm{} research paper (\emph{When the Regime Returns}) and its explanatory
report. It develops, from first principles, every tool the framework's
theory uses---decision-focused learning and the SPO+ surrogate,
concentration of measure, distances between distributions and Pinsker's
inequality, exponentiated-gradient dynamics, and invariant risk across
environments---then assembles them into the four results that organize the
contribution: the episode-transfer lemma (the invariance gap), the
eviction-rate lemma (the forgetting deficit), the post-reactivation regret
decomposition with its predicted-and-confirmed super-additive interaction,
and the break-even-dormancy pricing of retention against carrying costs.
The framework's central claim is that financial non-stationarity is
two-layered---regimes \emph{recur} (so dormant expertise is worth
retaining, and any reward-driven capacity allocator provably fails to
retain it) and recurrences \emph{differ} (so retained expertise must be
trained for across-episode invariance, or it faithfully preserves the last
episode's accidents)---and that the two layers are governed by independent
design choices whose product, but not either factor alone, recovers a
hand-designed oracle assignment without hand design. Every theoretical
statement is given with its full proof or an explicitly labeled sketch;
every assumption carries an audit of its idealizations; every experimental
number in the evidence part traces to a released result file; and the
document reports, at full prominence, the three places where our
pre-registered expectations were refuted by our own experiments (the
real-data structure gate, the competition-batching story, and the
high-signal regime where the entire mechanism inverts). \textbf{Prerequisites:}
probability at the level of a first graduate course, linear algebra, and
convexity basics; everything else is developed here.
\end{abstract}

\tableofcontents
\newpage

% ============================================================================
% PART I: MATHEMATICAL FOUNDATIONS (fragment; preamble lives in master file)
% ============================================================================
\part{Mathematical Foundations}

Before describing the \nm{} mechanism itself, we assemble the mathematical
toolkit its theory rests on. Four ingredients recur throughout. First,
\emph{decision-focused learning}: \nm{} never scores a forecast by how close it
is to the truth, but by the portfolio decision it induces, so we need the exact
geometry of the decision layer, the induced regret, and a trainable surrogate
(SPO+). Second, \emph{concentration of measure}: every retention and transfer
claim ultimately reduces to ``a quantity estimated from $5$--$7$ episodes is
close to its mean,'' and we must be honest about what such small samples can
and cannot certify. Third, \emph{distances between distributions}: the
regime-conditional guarantees we import are stated in divergence
(KL/total-variation) terms, and a small but load-bearing ``Pinsker bridge''
converts them into statements about bounded decision regret. Fourth,
\emph{online learning on the simplex}: the gate that routes episodes to niches
is an exponentiated-gradient update, and we must understand both what its
classical regret theory gives us and---just as importantly---what it does
\emph{not} give us once losses become endogenous. The final section formalizes
why pooled empirical risk minimization fails across non-stationary episodes and
why an across-episode variance penalty detects exactly the failure mode,
positioning the \nm{} objective inside the invariant-risk literature.

% ============================================================================
\section{Decision-Focused Learning from First Principles}
\label{sec:dd1-dfl}

\subsection{Prediction Versus Decision: the Predict-then-Optimize Pipeline}
\label{sec:dd1-pto}

Classical supervised learning evaluates a forecast $\hat{y}$ of an unknown
quantity $y$ by a \emph{prediction loss} such as the squared error
$\|\hat{y}-y\|_2^2$. But in \nm{}---as in a large class of operational
problems---the forecast is not the product. The forecast is an
\emph{intermediate input} to an optimization problem whose solution is the
actual decision, and the decision is what earns or loses money. This is the
\textbf{predict-then-optimize} pipeline \cite{elmachtoub2022smart,
mandi2020smart, mandi2024decision}:
\begin{equation}
\label{eq:dd1-pipeline}
x \;\xrightarrow{\ \text{model}\ }\; \hat{y} = f_\theta(x)
\;\xrightarrow{\ \text{optimizer}\ }\; z^*(\hat{y}) \in \argmax_{z \in Z}
F(z, \hat{y})
\;\xrightarrow{\ \text{nature reveals } y\ }\; \text{realized value } F(z^*(\hat{y}), y).
\end{equation}
Here $x$ are observable features (market state), $y \in \R^n$ is the unknown
parameter vector (next-period asset returns), $Z \subseteq \R^n$ is a known
feasible set of decisions (portfolios), and $F : Z \times \R^n \to \R$ is the
objective.

\begin{definition}[Linear-objective decision layer]
\label{def:dd1-linear-layer}
A decision layer is \textbf{linear-objective} if $F(z, y) = y^\top z$, i.e.\
the unknown parameter enters the objective linearly, while all problem
structure (combinatorics, budgets, bounds) is carried by the feasible set $Z$,
which does not depend on $y$. We write
\begin{equation}
\label{eq:dd1-oracle-def}
z^*(y) \in \argmax_{z \in Z} \, y^\top z, \qquad
V(y) := \max_{z \in Z} \, y^\top z ,
\end{equation}
and call the map $y \mapsto z^*(y)$ the \textbf{decision oracle} and $V$ the
\textbf{support function} of $Z$.
\end{definition}

\begin{remark}
Linearity of the objective in $y$ is a genuine restriction, but it covers
shortest paths, assignment, knapsack relaxations, and---our case---linear
portfolio construction. It buys us two things: the optimal value $V(y)$ is a
maximum of linear functions of $y$ and hence \emph{convex} in $y$, and the SPO+
surrogate of Section~\ref{sec:dd1-spoplus} exists in closed form.
\end{remark}

\nm{}'s decision layer is a cardinality-constrained, long-only allocation.

\begin{definition}[Cardinality-constrained long-only portfolio]
\label{def:dd1-portfolio}
Fix the universe size $n \in \N$, a per-asset cap $w_{\max} > 0$, a gross
budget $W > 0$, and a cardinality $k \in \{1, \dots, n\}$. The feasible set is
\begin{equation}
\label{eq:dd1-Z}
Z \;=\; \Bigl\{\, z \in [0, w_{\max}]^n \;:\; \textstyle\sum_{i=1}^n z_i \le W,
\;\; \|z\|_0 \le k \,\Bigr\},
\end{equation}
where $\|z\|_0 = |\{i : z_i \neq 0\}|$ counts nonzero positions. The objective
is the portfolio return $F(z, y) = y^\top z$ for a return vector $y \in \R^n$.
\end{definition}

Note that $Z$ is \emph{nonconvex} (the $\ell_0$ constraint carves the box into
$\binom{n}{k}$-many faces), so generic convex-optimization machinery does not
apply. Fortunately the problem admits an exact, closed-form, sort-based
solution.

\begin{theorem}[Exact decision oracle]
\label{thm:dd1-oracle}
Assume $W \ge k\, w_{\max}$ (the budget never binds before the cap and the
cardinality do). For $y \in \R^n$ let
\[
P(y) = \{ i : y_i > 0 \}, \qquad
S^*(y) = \text{the } \min\{k, |P(y)|\} \text{ indices of } P(y) \text{ with
largest } y_i,
\]
with ties broken toward the smaller index. Then
\begin{equation}
\label{eq:dd1-oracle-form}
z^*(y)_i \;=\; w_{\max} \,\mathbf{1}\{ i \in S^*(y) \}
\end{equation}
is an optimal solution of $\max_{z \in Z} y^\top z$, with optimal value
$V(y) = w_{\max} \sum_{i \in S^*(y)} y_i \ge 0$.
\end{theorem}

\begin{proof}
\emph{Feasibility.} The vector in \eqref{eq:dd1-oracle-form} has entries in
$\{0, w_{\max}\} \subseteq [0, w_{\max}]$, support size $|S^*(y)| \le k$, and
gross exposure $|S^*(y)|\, w_{\max} \le k\, w_{\max} \le W$. So
$z^*(y) \in Z$.

\emph{Optimality.} Let $z \in Z$ be arbitrary with support
$T = \{i : z_i > 0\}$, so $|T| \le k$. Using $0 \le z_i \le w_{\max}$ and
writing $(u)_+ = \max\{u, 0\}$,
\[
y^\top z \;=\; \sum_{i \in T} y_i z_i
\;\le\; \sum_{i \in T} (y_i)_+ \, z_i
\;\le\; w_{\max} \sum_{i \in T} (y_i)_+ .
\]
The first inequality holds termwise: if $y_i \le 0$ then $y_i z_i \le 0 =
(y_i)_+ z_i$; if $y_i > 0$ the two sides agree. Now
$\sum_{i \in T} (y_i)_+$ is a sum of at most $k$ of the nonnegative numbers
$\{(y_i)_+\}_{i=1}^n$, hence is at most the sum of the $k$ largest of them,
which equals $\sum_{i \in S^*(y)} y_i$ (indices outside $P(y)$ contribute
$(y_i)_+ = 0$ and are never needed). Therefore
$y^\top z \le w_{\max} \sum_{i \in S^*(y)} y_i = y^\top z^*(y)$. Finally
$V(y) \ge 0$ because $z = 0 \in Z$.
\end{proof}

\begin{keypoint}
The oracle is a \textbf{sort}: rank the forecast coordinates, keep the top $k$
strictly positive ones, allocate the cap $w_{\max}$ to each. It costs
$O(n \log n)$, is exact (no relaxation gap despite the nonconvex $\ell_0$
constraint), and will be called \emph{twice per gradient step} by the SPO+
subgradient of Section~\ref{sec:dd1-spoplus}. The standing assumption
$W \ge k\, w_{\max}$ is maintained throughout the document; when it fails the
oracle becomes a fractional-knapsack rule and only the boundary case changes.
\end{keypoint}

Throughout, we fix the tie-breaking rule above so that $z^*(\cdot)$ is a
well-defined \emph{function} (not a set-valued map); all statements below are
for this selection.

\subsection{Decision Regret and Its Geometry}
\label{sec:dd1-regret}

\begin{definition}[Decision regret]
\label{def:dd1-regret}
For a forecast $\hat{y} \in \R^n$ and realized returns $y \in \R^n$, the
\textbf{decision regret} (also: SPO loss \cite{elmachtoub2022smart}) is
\begin{equation}
\label{eq:dd1-regret}
\rho(\hat{y}; y) \;=\; F\bigl(z^*(y), y\bigr) - F\bigl(z^*(\hat{y}), y\bigr)
\;=\; y^\top z^*(y) - y^\top z^*(\hat{y}).
\end{equation}
It is the return given up by acting on the forecast instead of acting on the
truth.
\end{definition}

\begin{proposition}[Range of the regret]
\label{prop:dd1-regret-range}
Suppose $\|y\|_\infty \le y_{\max}$. Then for every $\hat{y} \in \R^n$,
\begin{equation}
\label{eq:dd1-B}
0 \;\le\; \rho(\hat{y}; y) \;\le\; B, \qquad B := 2\, k\, w_{\max}\, y_{\max}.
\end{equation}
\end{proposition}

\begin{proof}
\emph{Lower bound.} $z^*(\hat{y}) \in Z$, and $z^*(y)$ maximizes $y^\top z$
over $Z$ (Theorem~\ref{thm:dd1-oracle}), so
$y^\top z^*(y) \ge y^\top z^*(\hat{y})$.

\emph{Upper bound.} For any $z \in Z$ with support $T$, $|T| \le k$ and
$0 \le z_i \le w_{\max}$, so
\[
\bigl| y^\top z \bigr| \;\le\; \sum_{i \in T} |y_i|\, z_i
\;\le\; k\, w_{\max}\, y_{\max}.
\]
Hence $y^\top z^*(y) \le k\, w_{\max}\, y_{\max}$ and
$y^\top z^*(\hat{y}) \ge -\,k\, w_{\max}\, y_{\max}$; subtracting gives
$\rho \le 2 k\, w_{\max}\, y_{\max}$.
\end{proof}

\begin{remark}
The constant $B$ is attained only in a worst case where the forecast selects
$k$ assets that each lose $y_{\max}$ while the best portfolio gains $y_{\max}$
on each of $k$ assets. With daily equity returns this is astronomically
pessimistic; we return to this honestly in
Remark~\ref{rem:dd1-B-loose} when $B$ enters the Pinsker bridge.
\end{remark}

The next observation is small but conceptually central: the decision layer
destroys all information about the \emph{scale} of the forecast.

\begin{proposition}[Positive scale invariance; regret is a ranking loss]
\label{prop:dd1-scale-invariance}
For every $c > 0$ and $\hat{y} \in \R^n$:
\begin{enumerate}
\item $z^*(c\,\hat{y}) = z^*(\hat{y})$, and consequently
$\rho(c\,\hat{y}; y) = \rho(\hat{y}; y)$.
\item More generally, $\rho(\hat{y}; y)$ depends on $\hat{y}$ only through
the sign pattern $(\mathrm{sign}(\hat{y}_i))_{i}$ and the relative order of
the positive coordinates of $\hat{y}$.
\end{enumerate}
\end{proposition}

\begin{proof}
(1) Multiplication by $c > 0$ preserves the positivity set,
$P(c\hat{y}) = P(\hat{y})$, and preserves all pairwise comparisons
$c\hat{y}_i > c\hat{y}_j \iff \hat{y}_i > \hat{y}_j$ (including ties, which
are broken by the same index rule). Hence $S^*(c\hat{y}) = S^*(\hat{y})$ and,
by \eqref{eq:dd1-oracle-form}, $z^*(c\hat{y}) = z^*(\hat{y})$. The regret
\eqref{eq:dd1-regret} depends on $\hat{y}$ only through $z^*(\hat{y})$.

(2) Inspecting Theorem~\ref{thm:dd1-oracle}: the selected set $S^*(\hat{y})$
is determined by which coordinates are strictly positive and by the ordering
among those, and nothing else. Since $\rho(\cdot\,; y)$ factors through
$S^*(\hat{y})$, the claim follows.
\end{proof}

\begin{intuition}
The optimizer asks the forecast exactly one question: \emph{``which $k$ assets
do you like, and do you like them at all?''} Calibration, magnitude, even the
gap between the best and second-best asset are invisible to the decision. A
loss that trains the forecaster should therefore reward correct
\emph{rankings}, not small residuals---this is the entire raison d'\^etre of
decision-focused learning \cite{elmachtoub2022smart, shah2022lodl}.
\end{intuition}

\begin{example}[Zero MSE correlation with regret: a 4-asset case]
\label{ex:dd1-mse-vs-regret}
Let $n = 4$, $k = 2$, $w_{\max} = 0.5$, $W = 1$, and realized returns
\[
y = (0.05,\; 0.03,\; -0.02,\; 0.01).
\]
The oracle selects the two largest positive coordinates: $S^*(y) = \{1, 2\}$,
$z^*(y) = (0.5, 0.5, 0, 0)$, and the optimal value is
$V(y) = 0.5\,(0.05 + 0.03) = 0.04$.

\textbf{Forecast A (high MSE, zero regret).} Take $\hat{y}^A = 10\,y =
(0.5, 0.3, -0.2, 0.1)$, a wildly miscalibrated forecast. Its squared error is
\[
\mathrm{MSE}(\hat{y}^A)
= \tfrac{1}{4}\bigl(0.45^2 + 0.27^2 + 0.18^2 + 0.09^2\bigr)
= \tfrac{1}{4}(0.2025 + 0.0729 + 0.0324 + 0.0081)
\approx 0.0790 .
\]
But by Proposition~\ref{prop:dd1-scale-invariance},
$z^*(\hat{y}^A) = z^*(y)$, so $\rho(\hat{y}^A; y) = 0$: the decision is
\emph{perfect}.

\textbf{Forecast B (tiny MSE, large regret).} Take
$\hat{y}^B = (0.025,\; 0.030,\; -0.020,\; 0.031)$. Its squared error is
\[
\mathrm{MSE}(\hat{y}^B)
= \tfrac{1}{4}\bigl(0.025^2 + 0^2 + 0^2 + 0.021^2\bigr)
= \tfrac{1}{4}(0.000625 + 0.000441)
\approx 0.000267,
\]
about $296$ times smaller than forecast A's. Yet the induced ranking puts
assets $4$ and $2$ on top: $S^*(\hat{y}^B) = \{4, 2\}$,
$z^*(\hat{y}^B) = (0, 0.5, 0, 0.5)$, realized value
$y^\top z^*(\hat{y}^B) = 0.5\,(0.03 + 0.01) = 0.02$, and
\[
\rho(\hat{y}^B; y) = 0.04 - 0.02 = 0.02,
\]
i.e.\ \emph{half} the achievable return is forfeited. MSE ranks B far above
A; the decision ranks A strictly above B.
\end{example}

\begin{warning}
Example~\ref{ex:dd1-mse-vs-regret} is not a pathological corner case. In a
cross-section of assets the return differences that determine the top-$k$ set
are routinely smaller than the irreducible noise level, so an MSE-optimal
forecast spends its capacity matching magnitudes that the decision layer
ignores, while leaving the decision-relevant orderings to chance. This is the
failure that the SPO family of losses is designed to repair.
\end{warning}

\subsection{The SPO Loss, Its Pathologies, and the SPO+ Surrogate}
\label{sec:dd1-spoplus}

The natural training loss is the regret itself, viewed as a function of the
prediction: $\ell_{\mathrm{SPO}}(\hat{y}; y) := \rho(\hat{y}; y)$. It is,
unfortunately, almost unusable directly.

\begin{example}[Discontinuity and nonconvexity of the SPO loss]
\label{ex:dd1-spo-discont}
Let $n = 2$, $k = 1$, $w_{\max} = 1$, $W \ge 1$, $y = (0.04, 0.02)$, and
consider the forecast family $\hat{y}(t) = (t, 0.02)$ for $t \in \R$. The
oracle picks asset $1$ iff $t > 0.02$ (at $t = 0.02$ the tie-break also picks
asset $1$), and asset $2$ for $0 < t < 0.02$. Hence
\[
\ell_{\mathrm{SPO}}\bigl(\hat{y}(t); y\bigr) =
\begin{cases}
0, & t \ge 0.02 \quad (\text{decision } z = e_1,\ \text{value } 0.04),\\[2pt]
0.02, & t < 0.02 \quad (\text{decision } z = e_2,\ \text{value } 0.02),
\end{cases}
\]
a step function: \textbf{discontinuous} at $t = 0.02$, with derivative $0$
everywhere else. Any nonconstant function that is piecewise constant is
nonconvex (a finite convex function on $\R$ is continuous, and a continuous
piecewise-constant function is constant), so $\ell_{\mathrm{SPO}}$ is also
\textbf{nonconvex} in $\hat{y}$. Gradient-based training receives either no
signal (zero gradient) or an undefined one (at the jump).
\end{example}

The remedy of \cite{elmachtoub2022smart} is a convex surrogate, SPO+. Their
derivation is stated for cost \emph{minimization}; \nm{} maximizes return, so
we derive the maximization form explicitly via the substitution $c = -y$.

\paragraph{Derivation.} In the minimization convention one has cost vectors
$c$, decisions $w^*(c) \in \argmin_{w \in Z} c^\top w$, and the SPO+ loss
\begin{equation}
\label{eq:dd1-spoplus-min}
\ell_{\mathrm{SPO+}}^{\min}(\hat{c}; c)
\;=\; \max_{w \in Z}\,\bigl\{ (c - 2\hat{c})^\top w \bigr\}
\;+\; 2\,\hat{c}^\top w^*(c) \;-\; c^\top w^*(c),
\end{equation}
obtained from the SPO loss by a duality relaxation and the (provably
tightest) scaling $\hat{c} \mapsto 2\hat{c}$ \cite{elmachtoub2022smart}.
Substituting $c = -y$ and $\hat{c} = -\hat{y}$, and noting that
$\argmin_{w} (-y)^\top w = \argmax_{z} y^\top z$ so $w^*(c) = z^*(y)$, the
three terms of \eqref{eq:dd1-spoplus-min} become
$\max_z (2\hat{y} - y)^\top z$, $-2\hat{y}^\top z^*(y)$, and
$+\,y^\top z^*(y)$. We adopt the result as our definition.

\begin{definition}[SPO+ loss, maximization form]
\label{def:dd1-spoplus}
For forecast $\hat{y}$ and realized returns $y$,
\begin{equation}
\label{eq:dd1-spoplus}
\ell_{\mathrm{SPO+}}(\hat{y}; y)
\;=\; \underbrace{\max_{z \in Z}\, (2\hat{y} - y)^\top z}_{\textstyle
= \,V(2\hat{y} - y)}
\;-\; 2\,\hat{y}^\top z^*(y)
\;+\; y^\top z^*(y).
\end{equation}
Both extremal problems are instances of the same oracle
(Theorem~\ref{thm:dd1-oracle}), applied to the \emph{perturbed} vector
$2\hat{y} - y$ and to the true vector $y$.
\end{definition}

A reassuring sanity check: at the truth, the surrogate vanishes,
\[
\ell_{\mathrm{SPO+}}(y; y) = V(y) - 2\,y^\top z^*(y) + y^\top z^*(y)
= V(y) - y^\top z^*(y) = 0 .
\]

\begin{proposition}[Convexity in the prediction]
\label{prop:dd1-spoplus-convex}
For every fixed $y$, the map $\hat{y} \mapsto \ell_{\mathrm{SPO+}}(\hat{y}; y)$
is convex on $\R^n$ (indeed, piecewise affine).
\end{proposition}

\begin{proof}
The first term is $\hat{y} \mapsto V(2\hat{y} - y) =
\max_{z \in Z} \{ 2 z^\top \hat{y} - z^\top y \}$, a pointwise maximum of
affine functions of $\hat{y}$, hence convex (a supremum of convex functions is
convex, and affine functions are convex). The second term
$-2\,\hat{y}^\top z^*(y)$ is affine in $\hat{y}$ ($z^*(y)$ is a constant once
$y$ is fixed), and the third term is a constant. A sum of convex functions is
convex. Since $Z$ has finitely many extreme points (each coordinate of an
optimal $z$ is $0$ or $w_{\max}$ on one of finitely many supports), the
maximum is over finitely many affine functions, so the loss is piecewise
affine.
\end{proof}

\begin{proposition}[Majorization: SPO+ upper-bounds the decision regret]
\label{prop:dd1-spoplus-majorize}
For all $\hat{y}, y \in \R^n$,
\begin{equation}
\label{eq:dd1-majorize}
\rho(\hat{y}; y) \;\le\; \ell_{\mathrm{SPO+}}(\hat{y}; y).
\end{equation}
\end{proposition}

\begin{proof}
Lower-bound the maximum in \eqref{eq:dd1-spoplus} by the particular feasible
point $z = z^*(\hat{y})$:
\[
\ell_{\mathrm{SPO+}}(\hat{y}; y)
\;\ge\; (2\hat{y} - y)^\top z^*(\hat{y}) - 2\,\hat{y}^\top z^*(y)
+ y^\top z^*(y)
\;=\; 2\,\hat{y}^\top z^*(\hat{y}) - y^\top z^*(\hat{y})
- 2\,\hat{y}^\top z^*(y) + y^\top z^*(y).
\]
Now $z^*(\hat{y})$ maximizes $\hat{y}^\top z$ over $Z$, so
$\hat{y}^\top z^*(\hat{y}) \ge \hat{y}^\top z^*(y)$, i.e.\
$2\,\hat{y}^\top z^*(\hat{y}) - 2\,\hat{y}^\top z^*(y) \ge 0$. Dropping this
nonnegative quantity,
\[
\ell_{\mathrm{SPO+}}(\hat{y}; y)
\;\ge\; y^\top z^*(y) - y^\top z^*(\hat{y})
\;=\; \rho(\hat{y}; y). \qedhere
\]
\end{proof}

Combining Propositions~\ref{prop:dd1-regret-range} and
\ref{prop:dd1-spoplus-majorize}: minimizing the convex surrogate pushes down
the true (discontinuous) regret. The surrogate is moreover statistically
faithful in the following sense.

\begin{theorem}[Fisher consistency of SPO+; \cite{elmachtoub2022smart}, stated
without proof]
\label{thm:dd1-fisher}
Suppose the conditional distribution of $y$ given $x$ is continuous, centrally
symmetric about its mean $\bar{y}(x) = \E[y \mid x]$, and assigns the optimal
decision uniquely for almost every realization. Then the minimizer of the
conditional SPO+ risk $\E\bigl[\ell_{\mathrm{SPO+}}(\hat{y}; y) \mid x\bigr]$
over $\hat{y} \in \R^n$ is $\hat{y} = \bar{y}(x)$; in particular, minimizing
SPO+ risk also minimizes the true SPO (regret) risk.
\end{theorem}

\paragraph{Subgradient.} Convexity (Proposition~\ref{prop:dd1-spoplus-convex})
gives the loss a nonempty subdifferential everywhere; the support-function
structure makes one element of it available in closed form. We first record
the elementary lemma usually attributed to Danskin in this finite setting.

\begin{lemma}[Subgradient of a support function]
\label{lem:dd1-danskin}
Let $h(u) = \max_{z \in Z} u^\top z$ with $Z$ compact, and let
$\bar{z} \in \argmax_{z \in Z} u^\top z$. Then $\bar{z} \in \partial h(u)$,
i.e.\ $h(u') \ge h(u) + \bar{z}^\top (u' - u)$ for all $u'$.
\end{lemma}

\begin{proof}
$h(u') \ge u'^\top \bar{z} = u^\top \bar{z} + (u' - u)^\top \bar{z}
= h(u) + \bar{z}^\top (u' - u)$, where the first step holds because $\bar{z}$
is feasible for the maximization defining $h(u')$ and the last because
$\bar{z}$ attains $h(u)$.
\end{proof}

\begin{proposition}[SPO+ subgradient and the linear-model chain rule]
\label{prop:dd1-spoplus-grad}
Fix $y$ and let $\hat{y} \in \R^n$. Then
\begin{equation}
\label{eq:dd1-subgrad}
g(\hat{y}; y) \;:=\; 2\,\Bigl( z^*\bigl(2\hat{y} - y\bigr) - z^*(y) \Bigr)
\;\in\; \partial_{\hat{y}}\, \ell_{\mathrm{SPO+}}(\hat{y}; y),
\end{equation}
so a subgradient-descent step moves the forecast in the direction
\begin{equation}
\label{eq:dd1-descent-dir}
-\,g(\hat{y}; y) \;=\; 2\,\Bigl( z^*(y) - z^*\bigl(2\hat{y} - y\bigr) \Bigr),
\end{equation}
i.e.\ \emph{toward} the assets the truth selects and \emph{away from} the
assets the perturbed forecast selects. If the forecast is produced by a linear
model $\hat{y} = \Theta x$ with $\Theta \in \R^{n \times d}$ and features
$x \in \R^d$, then by the affine chain rule for subgradients,
\begin{equation}
\label{eq:dd1-chain}
g(\Theta x; y)\, x^\top \;\in\; \partial_{\Theta}\,
\ell_{\mathrm{SPO+}}(\Theta x;\, y) \;\subseteq\; \R^{n \times d},
\end{equation}
and for a batch $\{(x_s, y_s)\}_{s=1}^m$ the subgradient of the average loss
is $\frac{1}{m} \sum_s g(\Theta x_s; y_s)\, x_s^\top$.
\end{proposition}

\begin{proof}
Write $\ell_{\mathrm{SPO+}}(\hat{y}; y) = h(2\hat{y} - y) - 2\,\hat{y}^\top
z^*(y) + \text{const}$, with $h$ the support function of $Z$. By
Lemma~\ref{lem:dd1-danskin}, $z^*(2\hat{y} - y) \in \partial h(2\hat{y} - y)$
(our oracle applied to the vector $2\hat{y} - y$ is exactly an argmax for
$h$). The map $\hat{y} \mapsto 2\hat{y} - y$ is affine with Jacobian $2 I$, so
by the affine chain rule $2\,z^*(2\hat{y} - y)$ is a subgradient of
$\hat{y} \mapsto h(2\hat{y} - y)$. The term $-2\,\hat{y}^\top z^*(y)$ is
differentiable with gradient $-2\,z^*(y)$. Subdifferentials of sums of convex
functions (one of them differentiable) add, giving \eqref{eq:dd1-subgrad}.
The chain rule \eqref{eq:dd1-chain} is the same affine-composition rule
applied to $\Theta \mapsto \Theta x$: for any $\Theta'$,
$\ell(\Theta' x) \ge \ell(\Theta x) + g^\top(\Theta' - \Theta)x
= \ell(\Theta x) + \langle g x^\top, \Theta' - \Theta\rangle_F$.
\end{proof}

\begin{remark}[Sign conventions]
\label{rem:dd1-sign}
In the cost-minimization convention of \cite{elmachtoub2022smart} the same
object reads $2\bigl(w^*(c) - w^*(2\hat{c} - c)\bigr)$. Under $c = -y$ the two
expressions map onto each other; the expression
$2\bigl(z^*(y) - z^*(2\hat{y} - y)\bigr)$ that appears in informal
descriptions of \nm{} is the \emph{descent direction}
\eqref{eq:dd1-descent-dir} of the maximization form, not the subgradient
itself. Both conventions produce the identical parameter update; only the
bookkeeping differs.
\end{remark}

\begin{intuition}
The update \eqref{eq:dd1-descent-dir} has a ``contrastive'' reading. The
oracle is run on the truth $y$ and on the reflection $2\hat{y} - y$ of the
truth through the forecast (the point such that $\hat{y}$ is the midpoint of
$y$ and $2\hat{y} - y$; the reflection exaggerates the forecast's deviation
from the truth). Assets selected by the truth but not by the exaggerated
forecast get their predicted returns pushed \emph{up} by $2 w_{\max} \eta$ per
step; assets selected only by the exaggerated forecast get pushed \emph{down}
by the same amount; assets on which both agree (or both pass) receive no
gradient. The loss is exactly zero gradient when the exaggerated forecast
already makes the right selection---training stops adjusting magnitudes the
moment the ranking is robustly correct, which is precisely the
scale-invariance of Proposition~\ref{prop:dd1-scale-invariance} made
operational.
\end{intuition}

\begin{example}[One SPO+ gradient step on four assets]
\label{ex:dd1-spoplus-step}
Continue with $n = 4$, $k = 2$, $w_{\max} = 0.5$, $W = 1$, true returns
$y = (0.05, 0.03, -0.02, 0.01)$, and current forecast
\[
\hat{y} = (0.02,\; 0.01,\; 0.04,\; 0.00).
\]

\emph{Step 1: the two oracle calls.} The truth selects
$S^*(y) = \{1, 2\}$, so $z^*(y) = (0.5, 0.5, 0, 0)$. The perturbed vector is
\[
2\hat{y} - y = (0.04 - 0.05,\; 0.02 - 0.03,\; 0.08 + 0.02,\; 0 - 0.01)
= (-0.01,\; -0.01,\; 0.10,\; -0.01),
\]
whose only positive coordinate is the third, so
$S^*(2\hat{y} - y) = \{3\}$ and $z^*(2\hat{y} - y) = (0, 0, 0.5, 0)$.

\emph{Step 2: loss value.} By \eqref{eq:dd1-spoplus},
\[
\ell_{\mathrm{SPO+}} = \underbrace{0.10 \times 0.5}_{V(2\hat y - y) = 0.05}
\;-\; \underbrace{2\,(0.02 + 0.01)(0.5)}_{2\hat{y}^\top z^*(y) = 0.03}
\;+\; \underbrace{0.04}_{y^\top z^*(y)} \;=\; 0.06 .
\]
For comparison, the true regret of $\hat{y}$: the forecast selects its top
two positive coordinates $S^*(\hat{y}) = \{3, 1\}$, realized value
$0.5\,(-0.02 + 0.05) = 0.015$, so
$\rho(\hat{y}; y) = 0.04 - 0.015 = 0.025 \le 0.06$, consistent with
Proposition~\ref{prop:dd1-spoplus-majorize}.

\emph{Step 3: subgradient and update.} By \eqref{eq:dd1-subgrad},
\[
g = 2\bigl( (0, 0, 0.5, 0) - (0.5, 0.5, 0, 0) \bigr)
= (-1,\; -1,\; 1,\; 0).
\]
A descent step with rate $\eta = 0.02$ gives
\[
\hat{y}' = \hat{y} - \eta\, g
= (0.02 + 0.02,\; 0.01 + 0.02,\; 0.04 - 0.02,\; 0.00)
= (0.04,\; 0.03,\; 0.02,\; 0.00).
\]
The new forecast ranks the assets $1 \succ 2 \succ 3 \succ 4$, selects
$S^*(\hat{y}') = \{1, 2\}$, and achieves $\rho(\hat{y}'; y) = 0$: one step
flipped the decision from losing $0.025$ to optimal. Note what the step did
\emph{not} do: it made the forecast magnitudes \emph{less} accurate for asset
$3$ in MSE terms en route to fixing the ranking.
\end{example}

\begin{keypoint}
Summary of the decision-focused toolkit used by \nm{}:
\begin{itemize}
\item the decision layer is an $O(n \log n)$ sort oracle
      (Theorem~\ref{thm:dd1-oracle});
\item the quantity that matters is the regret
      $\rho \in [0, B]$, a bounded \emph{ranking} loss
      (Propositions~\ref{prop:dd1-regret-range},
      \ref{prop:dd1-scale-invariance});
\item $\rho$ itself is untrainable
      (Example~\ref{ex:dd1-spo-discont}), so each niche head is trained on
      the convex majorant $\ell_{\mathrm{SPO+}}$
      (Definition~\ref{def:dd1-spoplus}), whose subgradient costs two oracle
      calls (Proposition~\ref{prop:dd1-spoplus-grad}).
\end{itemize}
Every ``episode risk'' appearing later in the document is an average of
$\ell_{\mathrm{SPO+}}$ (for training) or of $\rho$ (for evaluation) over the
days of an episode.
\end{keypoint}

% ============================================================================
\section{Concentration of Measure}
\label{sec:dd1-concentration}

\nm{}'s retention and transfer arguments repeatedly take the form: ``the
average regret of niche head $r$ over its $E_r$ stored episodes is small;
therefore its regret on a fresh episode of the same regime is small with high
probability.'' Since $E_r$ is tiny ($5$--$7$ in the companion experiments),
we need the sharpest available finite-sample inequalities---and an honest
accounting of how weak even those are at such sample sizes.

\subsection{Markov, Chebyshev, and Cantelli}

\begin{theorem}[Markov's inequality]
\label{thm:dd1-markov}
Let $X \ge 0$ be a random variable with $\E[X] < \infty$. For every $t > 0$,
\begin{equation}
\label{eq:dd1-markov}
\Pr\bigl( X \ge t \bigr) \;\le\; \frac{\E[X]}{t}.
\end{equation}
\end{theorem}

\begin{proof}
Since $X \ge 0$, the pointwise inequality
$t\,\mathbf{1}\{X \ge t\} \le X$ holds on every outcome: if $X \ge t$ the left
side is $t \le X$; otherwise the left side is $0 \le X$. Taking expectations
and using monotonicity and linearity,
$t\,\Pr(X \ge t) = \E\bigl[t\,\mathbf{1}\{X \ge t\}\bigr] \le \E[X]$.
Divide by $t$.
\end{proof}

\begin{corollary}[Chebyshev's inequality]
\label{cor:dd1-chebyshev}
Let $X$ have mean $\mu$ and variance $\sigma^2 < \infty$. For every $t > 0$,
\begin{equation}
\label{eq:dd1-chebyshev}
\Pr\bigl( |X - \mu| \ge t \bigr) \;\le\; \frac{\sigma^2}{t^2}.
\end{equation}
\end{corollary}

\begin{proof}
Apply Theorem~\ref{thm:dd1-markov} to the nonnegative variable
$(X - \mu)^2$ at level $t^2$:
$\Pr(|X - \mu| \ge t) = \Pr\bigl((X - \mu)^2 \ge t^2\bigr)
\le \E[(X-\mu)^2] / t^2 = \sigma^2 / t^2$.
\end{proof}

Chebyshev is two-sided. When only one direction of deviation is harmful---a
niche head performing \emph{worse} on a fresh episode than on its stored
ones---we can do strictly better.

\begin{theorem}[Cantelli's one-sided inequality]
\label{thm:dd1-cantelli}
Let $X$ have mean $\mu$ and variance $\sigma^2 < \infty$. For every $t > 0$,
\begin{equation}
\label{eq:dd1-cantelli}
\Pr\bigl( X - \mu \ge t \bigr) \;\le\; \frac{\sigma^2}{\sigma^2 + t^2}.
\end{equation}
\end{theorem}

\begin{proof}
Let $Y = X - \mu$, so $\E[Y] = 0$ and $\E[Y^2] = \sigma^2$. Fix an arbitrary
auxiliary constant $u > 0$. On the event $\{Y \ge t\}$ we have
$Y + u \ge t + u > 0$, hence $(Y + u)^2 \ge (t + u)^2$ there, and
$(Y+u)^2 \ge 0$ everywhere. Therefore
\[
\Pr( Y \ge t )
\;=\; \Pr\bigl( Y + u \ge t + u \bigr)
\;\le\; \Pr\bigl( (Y + u)^2 \ge (t + u)^2 \bigr)
\;\le\; \frac{\E\bigl[(Y + u)^2\bigr]}{(t + u)^2}
\;=\; \frac{\sigma^2 + u^2}{(t + u)^2},
\]
where the middle step is Markov (Theorem~\ref{thm:dd1-markov}) applied to the
nonnegative variable $(Y + u)^2$, and the last step expands
$\E[(Y+u)^2] = \E[Y^2] + 2u\,\E[Y] + u^2 = \sigma^2 + u^2$.

The bound holds for every $u > 0$, so we minimize the right-hand side. Let
$\varphi(u) = (\sigma^2 + u^2)/(t + u)^2$. Then
\[
\varphi'(u)
= \frac{2u\,(t+u)^2 - 2(t+u)(\sigma^2 + u^2)}{(t+u)^4}
= \frac{2\bigl[u(t + u) - (\sigma^2 + u^2)\bigr]}{(t+u)^3}
= \frac{2\,(ut - \sigma^2)}{(t+u)^3},
\]
which vanishes at $u^* = \sigma^2 / t$, is negative for $u < u^*$ and positive
for $u > u^*$; so $u^*$ is the global minimizer on $(0, \infty)$.
Substituting,
\[
\varphi(u^*)
= \frac{\sigma^2 + \sigma^4/t^2}{\bigl(t + \sigma^2/t\bigr)^2}
= \frac{\sigma^2\,(t^2 + \sigma^2)/t^2}{(t^2 + \sigma^2)^2 / t^2}
= \frac{\sigma^2}{\sigma^2 + t^2}. \qedhere
\]
\end{proof}

\begin{remark}[Why one-sidedness matters for the episode-transfer lemma]
\label{rem:dd1-onesided}
In the transfer argument of Part II, the certified quantity is the average
regret of a niche head over its stored episodes, and the failure event is
one-directional: the fresh-episode regret exceeding that average by more than
a tolerance $t$. A fresh episode on which the head does \emph{better} than
its certificate is a windfall, not a failure. Cantelli prices exactly the
harmful tail: at $t = \sigma$, Chebyshev gives the vacuous bound $1$ split
across two tails (effectively $1/2$ per side by symmetry heuristics, but with
no valid one-sided guarantee better than $1$), while Cantelli gives
$\sigma^2/(2\sigma^2) = 1/2$ \emph{as a theorem}, and at $t = 2\sigma$ gives
$1/5$ versus Chebyshev's two-sided $1/4$. Moreover Cantelli requires no
boundedness and no symmetry---only a variance---which suits regret
distributions that are skewed by construction (they live in $[0, B]$ and pile
up near $0$ for a well-adapted head).
\end{remark}

\subsection{Hoeffding's Inequality}

Variance-based bounds decay polynomially in $t$. For \emph{bounded} variables
---and the regret is bounded in $[0, B]$ by
Proposition~\ref{prop:dd1-regret-range}---exponential tails are available.
The engine is a bound on the moment generating function.

\begin{lemma}[Hoeffding's lemma; proof sketch]
\label{lem:dd1-hoeffding-lemma}
Let $X \in [a, b]$ almost surely, with $\E[X] = 0$. Then for all
$\lambda \in \R$,
\begin{equation}
\label{eq:dd1-hoeffding-lemma}
\E\bigl[ e^{\lambda X} \bigr] \;\le\;
\exp\!\Bigl( \frac{\lambda^2 (b - a)^2}{8} \Bigr).
\end{equation}
\end{lemma}

\begin{proof}[Proof sketch]
Define the cumulant generating function $\psi(\lambda) = \log \E[e^{\lambda
X}]$, which is finite for all $\lambda$ (boundedness) and twice
differentiable. Direct computation gives $\psi(0) = 0$ and
$\psi'(0) = \E[X] = 0$. The second derivative is
\[
\psi''(\lambda) = \E_\lambda[X^2] - \bigl(\E_\lambda[X]\bigr)^2
= \Var_\lambda(X),
\]
the variance of $X$ under the exponentially tilted law
$d\Pr_\lambda \propto e^{\lambda X}\, d\Pr$ (this is the step we do not carry
out in full detail; it is a one-line differentiation under the expectation,
legitimate by dominated convergence on a bounded variable). The tilted law is
still supported in $[a, b]$, and \emph{any} random variable $V \in [a, b]$
satisfies
\[
\Var(V) \;\le\; \E\Bigl[\Bigl(V - \tfrac{a + b}{2}\Bigr)^2\Bigr]
\;\le\; \Bigl(\tfrac{b - a}{2}\Bigr)^2,
\]
where the first inequality holds because the variance minimizes the expected
squared deviation over all centering constants, and the second because
$|V - (a+b)/2| \le (b-a)/2$ pointwise. Hence
$\psi''(\lambda) \le (b - a)^2 / 4$ for all $\lambda$, and Taylor's theorem
with the exact remainder gives, for some $\xi$ between $0$ and $\lambda$,
\[
\psi(\lambda) = \psi(0) + \lambda\,\psi'(0)
+ \tfrac{\lambda^2}{2}\,\psi''(\xi)
\;\le\; \frac{\lambda^2 (b-a)^2}{8}.
\]
Exponentiate.
\end{proof}

\begin{theorem}[Hoeffding's inequality]
\label{thm:dd1-hoeffding}
Let $X_1, \dots, X_n$ be independent with $X_i \in [a_i, b_i]$ almost surely,
and let $\bar{X} = \frac{1}{n}\sum_{i=1}^n X_i$, $\mu = \E[\bar{X}]$. Then for
every $t > 0$,
\begin{equation}
\label{eq:dd1-hoeffding}
\Pr\bigl( \bar{X} - \mu \ge t \bigr)
\;\le\; \exp\!\Bigl( - \frac{2 n^2 t^2}{\sum_{i=1}^n (b_i - a_i)^2} \Bigr),
\end{equation}
and the same bound holds for $\Pr(\mu - \bar{X} \ge t)$; a union bound gives
the two-sided version with an extra factor $2$. In particular, if
$X_i \in [0, B]$ for all $i$, then with probability at least $1 - \delta$,
\begin{equation}
\label{eq:dd1-hoeffding-radius}
\bigl| \bar{X} - \mu \bigr| \;\le\; B \sqrt{ \frac{\ln(2/\delta)}{2n} }.
\end{equation}
\end{theorem}

\begin{proof}
Let $S = \sum_i (X_i - \E X_i)$, so $\{\bar{X} - \mu \ge t\} =
\{S \ge nt\}$. For any $\lambda > 0$, by Markov's inequality applied to the
nonnegative variable $e^{\lambda S}$ (the Chernoff step) and independence,
\[
\Pr( S \ge nt ) \;\le\; e^{-\lambda n t}\, \E\bigl[ e^{\lambda S} \bigr]
\;=\; e^{-\lambda n t} \prod_{i=1}^n \E\bigl[ e^{\lambda (X_i - \E X_i)}
\bigr]
\;\le\; \exp\!\Bigl( -\lambda n t + \frac{\lambda^2}{8} \sum_{i=1}^n
(b_i - a_i)^2 \Bigr),
\]
using Lemma~\ref{lem:dd1-hoeffding-lemma} for each centered $X_i - \E X_i \in
[a_i - \E X_i,\, b_i - \E X_i]$, an interval of the same width $b_i - a_i$.
The exponent $-\lambda n t + \frac{\lambda^2}{8}\sum_i (b_i - a_i)^2$ is a
convex parabola in $\lambda$, minimized at
$\lambda^* = 4nt / \sum_i (b_i - a_i)^2 > 0$, with minimum value
$-2 n^2 t^2 / \sum_i (b_i - a_i)^2$. This proves \eqref{eq:dd1-hoeffding};
the lower tail follows by applying the result to $-X_i$. For
\eqref{eq:dd1-hoeffding-radius}, set $a_i = 0$, $b_i = B$, equate the
two-sided bound $2\exp(-2nt^2/B^2)$ to $\delta$, and solve for $t$.
\end{proof}

\subsection{Empirical Bernstein: Paying the Observed Variance, Not the Range}

Hoeffding's radius \eqref{eq:dd1-hoeffding-radius} scales with the worst-case
range $B$, no matter how concentrated the data actually are. When the sample
is tiny this is painful, because $B$ for the decision regret is a gross
worst-case constant (Proposition~\ref{prop:dd1-regret-range}). Bernstein-type
inequalities replace $B$ by the standard deviation, but the classical
Bernstein bound needs the \emph{true} variance, which is unknown. The
empirical-Bernstein inequality of Maurer and Pontil
\cite{maurer2009empirical} closes the loop by using the \emph{sample}
variance, at the price of a lower-order $B/n$ term.

\begin{theorem}[Empirical Bernstein inequality \cite{maurer2009empirical};
stated without proof]
\label{thm:dd1-empbernstein}
Let $X_1, \dots, X_n$ ($n \ge 2$) be i.i.d.\ with values in $[0, B]$, sample
mean $\bar{X}$, and sample variance
$V_n = \frac{1}{n - 1}\sum_{i=1}^n (X_i - \bar{X})^2$. Then with probability
at least $1 - \delta$,
\begin{equation}
\label{eq:dd1-empbernstein}
\mu \;\le\; \bar{X} \;+\; \sqrt{ \frac{2\, V_n \ln(2/\delta)}{n} }
\;+\; \frac{7\, B \ln(2/\delta)}{3\,(n - 1)} .
\end{equation}
\end{theorem}

\begin{intuition}
Read \eqref{eq:dd1-empbernstein} as: \emph{you pay the spread you actually
observed} (the $\sqrt{V_n/n}$ term), plus a \emph{fixed entrance fee}
proportional to $B/n$ for the privilege of estimating the variance from the
same data. When the observations are tightly clustered ($\sqrt{V_n} \ll B$)
and $n$ is moderate, the first term is far below Hoeffding's
$B\sqrt{\ln(2/\delta)/2n}$ and the entrance fee is negligible. When $n$ is
very small, the entrance fee dominates everything.
\end{intuition}

\begin{remark}[Why this matters at $E_r = 5$--$7$ episodes]
\label{rem:dd1-small-n}
In \nm{}, the per-niche certificate averages the episode risks of the
$E_r \in \{5, \dots, 7\}$ episodes stored in niche $r$'s memory. The entire
design of the niche structure---episodes are routed to a niche precisely
because the niche head's decision behavior on them is similar---aims to make
the \emph{within-niche} spread of episode risks small, i.e.\ $V_n \ll B^2$.
Empirical Bernstein is the inequality that can, in principle, convert that
design property into a tighter certificate, because it is the only one in
this family whose radius adapts to the realized within-niche concentration.
Hoeffding cannot see the clustering at all. The sobering quantitative reality
at $n \in \{5, \dots, 7\}$ is worked out in the next example.
\end{remark}

\begin{example}[Estimation radii with $6$ episodes]
\label{ex:dd1-radius}
Take $n = 6$ episodes, regret range $B = 0.4$, confidence $\delta = 0.1$, so
$\ln(2/\delta) = \ln 20 \approx 2.9957$.

\emph{Hoeffding.} The two-sided radius \eqref{eq:dd1-hoeffding-radius} is
\[
B \sqrt{\frac{\ln(2/\delta)}{2n}}
= 0.4 \times \sqrt{\frac{2.9957}{12}}
= 0.4 \times \sqrt{0.24964}
= 0.4 \times 0.49964
\approx 0.1999 \approx 0.20 .
\]
The certificate is ``mean regret within $\pm 0.20$''---\emph{half the entire
range} $[0, 0.4]$. As a quantitative guarantee this is nearly vacuous.

\emph{Empirical Bernstein.} Suppose the six stored episode regrets are
tightly clustered with sample standard deviation $\sqrt{V_n} = 0.05$. The
variance term of \eqref{eq:dd1-empbernstein} is
\[
\sqrt{\frac{2 \times 0.0025 \times 2.9957}{6}}
= \sqrt{0.0024964} \approx 0.050,
\]
a four-fold improvement---but the entrance fee is
\[
\frac{7 \times 0.4 \times 2.9957}{3 \times 5}
= \frac{8.388}{15} \approx 0.559,
\]
which alone exceeds the whole range $B$. At $n = 6$ the empirical-Bernstein
certificate is \emph{worse} than Hoeffding's; the variance adaptivity only
pays once $n$ is large enough that the $B/(n-1)$ term recedes (here, roughly
$n \gtrsim 40$ for these constants).
\end{example}

\begin{warning}
Example~\ref{ex:dd1-radius} is the honest small-sample accounting that frames
every finite-sample statement in this document: at $E_r = 5$--$7$ episodes,
\emph{no} distribution-free concentration inequality certifies the per-niche
mean regret to a useful tolerance. The role of these bounds in the \nm{}
theory is therefore \emph{structural}, not numerical: they identify
\emph{which} quantities control transfer (the within-niche mean and
variance of episode risk, and only those), and they predict the correct
\emph{qualitative} ordering of mechanisms (lower within-niche variance
$\Rightarrow$ tighter transfer). That prediction is exactly what the
variance-penalized objective of Section~\ref{sec:dd1-invariance} acts on, and
it is tested empirically rather than certified theoretically.
\end{warning}

% ============================================================================
\section{Distances Between Distributions and the Pinsker Bridge}
\label{sec:dd1-distances}

\nm{} inherits from the invariant predict-and-optimize line of work
\cite{yan2026invpnco} guarantees of the form ``within a regime, the
divergence between the training mixture of episodes and a deployment episode
is at most $\epsilon$.'' Such statements are about \emph{distributions}; what
we need are statements about \emph{regret}. This section builds the standard
bridge: total variation converts distributional closeness into closeness of
bounded expectations, and Pinsker's inequality converts KL divergence into
total variation.

Throughout, $P$ and $Q$ are probability measures on a common measurable space
$(\Omega, \mathcal{F})$, and we fix a dominating $\sigma$-finite measure
$\mu$ (e.g.\ $\mu = P + Q$) with densities $p = dP/d\mu$, $q = dQ/d\mu$. All
statements are independent of the choice of $\mu$.

\subsection{Total Variation Distance}

\begin{definition}[Total variation]
\label{def:dd1-tv}
The \textbf{total variation distance} between $P$ and $Q$ is
\begin{equation}
\label{eq:dd1-tv-def}
\TV(P, Q) \;=\; \sup_{A \in \mathcal{F}} \bigl| P(A) - Q(A) \bigr|.
\end{equation}
\end{definition}

\begin{proposition}[Variational and $L_1$ characterizations]
\label{prop:dd1-tv-variational}
With $p, q$ as above,
\begin{equation}
\label{eq:dd1-tv-three}
\TV(P, Q)
\;=\; \frac{1}{2} \int_\Omega |p - q| \, d\mu
\;=\; \sup_{g : \Omega \to [0, 1]} \Bigl| \E_P[g] - \E_Q[g] \Bigr|,
\end{equation}
where the supremum is over all measurable functions $g$ with values in
$[0,1]$, and is attained at $g = \mathbf{1}_{A^*}$ with
$A^* = \{p \ge q\}$.
\end{proposition}

\begin{proof}
Write $\Delta = p - q$ and note the key cancellation
\begin{equation}
\label{eq:dd1-cancellation}
\int_\Omega \Delta \, d\mu = \int p \, d\mu - \int q \, d\mu = 1 - 1 = 0
\quad\Longrightarrow\quad
\int_{A^*} \Delta \, d\mu \;=\; \int_{(A^*)^c} (-\Delta) \, d\mu ,
\end{equation}
where $A^* = \{\Delta \ge 0\}$; both sides of the implication equal
$\frac{1}{2}\int |\Delta| \, d\mu$ since
$\int |\Delta| = \int_{A^*} \Delta + \int_{(A^*)^c} (-\Delta)$.

\emph{Step 1: every $[0,1]$-valued $g$ is dominated.} For any such $g$,
\[
\E_P[g] - \E_Q[g] = \int g\, \Delta \, d\mu
= \int_{A^*} g\, \Delta \, d\mu + \int_{(A^*)^c} g\, \Delta \, d\mu
\;\le\; \int_{A^*} 1 \cdot \Delta \, d\mu + \int_{(A^*)^c} 0 \cdot \Delta \,
d\mu
= \tfrac{1}{2} \textstyle\int |\Delta| \, d\mu,
\]
using $0 \le g \le 1$, $\Delta \ge 0$ on $A^*$, and $\Delta < 0$ on
$(A^*)^c$. Applying the same bound to $1 - g$ (also $[0,1]$-valued) gives
$\E_Q[g] - \E_P[g] \le \frac{1}{2}\int|\Delta|$ as well (note
$\E_Q[g]-\E_P[g] = \E_P[1-g]-\E_Q[1-g]$). Hence
$\sup_g |\E_P g - \E_Q g| \le \frac{1}{2}\int |\Delta|\, d\mu$.

\emph{Step 2: attainment.} The choice $g = \mathbf{1}_{A^*}$ achieves
$\E_P[g] - \E_Q[g] = \int_{A^*} \Delta\, d\mu = \frac{1}{2}\int|\Delta|\,
d\mu$ by \eqref{eq:dd1-cancellation}. So the supremum over $[0,1]$-valued $g$
\emph{equals} $\frac{1}{2}\int |\Delta|$.

\emph{Step 3: events.} Indicators $\mathbf{1}_A$ are particular
$[0,1]$-valued functions, so
$\sup_A |P(A) - Q(A)| \le \sup_g |\E_P g - \E_Q g|$; and the optimizer found
in Step 2 \emph{is} an indicator, so the two suprema coincide, and both equal
$\frac{1}{2}\int|p - q|\,d\mu$.
\end{proof}

\begin{intuition}
$\TV$ answers: \emph{over all bounded tests normalized to $[0,1]$, how
differently can the two distributions score?} The optimal test is the
likelihood-ratio region $\{p \ge q\}$. The $[0,1]$-function form---rather
than the events-only form---is the one we will actually use, because a
bounded loss rescaled into $[0,1]$ is such a function but is not an
indicator.
\end{intuition}

\subsection{KL Divergence and Gibbs' Inequality}

\begin{definition}[Kullback--Leibler divergence]
\label{def:dd1-kl}
If $P \ll Q$ (i.e.\ $q = 0 \Rightarrow p = 0$, $\mu$-a.e.),
\begin{equation}
\label{eq:dd1-kl-def}
\KL(P \,\|\, Q) \;=\; \int_{\{p > 0\}} p \,\log \frac{p}{q} \, d\mu
\;=\; \E_P\Bigl[ \log \frac{p}{q} \Bigr],
\end{equation}
with the conventions $0 \log 0 = 0$; if $P \not\ll Q$ we set
$\KL(P\|Q) = +\infty$.
\end{definition}

\begin{proposition}[Gibbs' inequality / non-negativity]
\label{prop:dd1-gibbs}
$\KL(P \| Q) \ge 0$, with equality if and only if $P = Q$.
\end{proposition}

\begin{proof}
Assume $P \ll Q$ (otherwise $\KL = \infty > 0$). Then
\[
-\KL(P\|Q) = \E_P\Bigl[ \log \frac{q}{p} \Bigr]
\;\le\; \log \E_P\Bigl[ \frac{q}{p} \Bigr]
= \log \int_{\{p > 0\}} q \, d\mu
\;\le\; \log 1 = 0,
\]
where the first inequality is Jensen's inequality for the concave function
$\log$, and the second uses $\int_{\{p>0\}} q\, d\mu \le \int q \, d\mu = 1$.
For the equality case: $\log$ is \emph{strictly} concave, so Jensen is an
equality iff $q/p$ is $P$-a.s.\ constant; combined with
$\int_{\{p>0\}} q = 1$ (forced by the second inequality being tight) the
constant is $1$, i.e.\ $p = q$ $\mu$-a.e., i.e.\ $P = Q$.
\end{proof}

\subsection{Pinsker's Inequality}

We give the classical proof \cite{tsybakov2009introduction} in two moves: a
data-processing (two-point) reduction via the log-sum inequality, then the
scalar inequality for Bernoulli pairs.

\begin{lemma}[Log-sum inequality]
\label{lem:dd1-logsum}
For nonnegative reals $a_1, \dots, a_m$ and positive reals
$b_1, \dots, b_m$,
\begin{equation}
\label{eq:dd1-logsum}
\sum_{i=1}^m a_i \log \frac{a_i}{b_i}
\;\ge\; \Bigl( \sum_{i=1}^m a_i \Bigr)
\log \frac{\sum_{i=1}^m a_i}{\sum_{i=1}^m b_i},
\end{equation}
with the convention $0 \log 0 = 0$. The same holds with sums replaced by
integrals of nonnegative functions.
\end{lemma}

\begin{proof}
Let $\phi(t) = t \log t$ for $t > 0$, $\phi(0) = 0$; $\phi$ is convex on
$[0, \infty)$ since $\phi''(t) = 1/t > 0$. Set $b = \sum_i b_i$ and apply
Jensen's inequality with weights $b_i / b$ to the points $t_i = a_i / b_i$:
\[
\sum_{i=1}^m \frac{b_i}{b}\, \phi\Bigl( \frac{a_i}{b_i} \Bigr)
\;\ge\; \phi\Bigl( \sum_{i=1}^m \frac{b_i}{b} \cdot \frac{a_i}{b_i} \Bigr)
= \phi\Bigl( \frac{\sum_i a_i}{b} \Bigr).
\]
Multiplying both sides by $b$ and expanding $\phi$:
\[
\sum_i b_i \cdot \frac{a_i}{b_i} \log\frac{a_i}{b_i}
= \sum_i a_i \log \frac{a_i}{b_i}
\;\ge\; b \cdot \frac{\sum_i a_i}{b} \log \frac{\sum_i a_i}{b}
= \Bigl(\sum_i a_i\Bigr) \log \frac{\sum_i a_i}{\sum_i b_i}.
\]
The integral version is Jensen for the same $\phi$ with the normalized
measure $b(\omega)\,d\mu / \int b \, d\mu$.
\end{proof}

\begin{lemma}[Two-point (data-processing) reduction]
\label{lem:dd1-dataproc}
For any event $A \in \mathcal{F}$, writing $a = P(A)$, $b = Q(A)$ and
$\mathrm{kl}(a \,\|\, b) = a \log\frac{a}{b} + (1-a)\log\frac{1-a}{1-b}$ for
the Bernoulli KL,
\begin{equation}
\label{eq:dd1-dataproc}
\KL(P \,\|\, Q) \;\ge\; \mathrm{kl}\bigl( P(A) \,\|\, Q(A) \bigr).
\end{equation}
\end{lemma}

\begin{proof}
Split the defining integral over the partition $\{A, A^c\}$ and apply the
integral log-sum inequality (Lemma~\ref{lem:dd1-logsum}) on each cell:
\[
\KL(P\|Q)
= \int_A p \log \frac{p}{q}\, d\mu + \int_{A^c} p \log \frac{p}{q}\, d\mu
\;\ge\; \Bigl(\int_A p\Bigr) \log \frac{\int_A p}{\int_A q}
+ \Bigl(\int_{A^c} p\Bigr) \log \frac{\int_{A^c} p}{\int_{A^c} q}
= \mathrm{kl}\bigl(P(A) \| Q(A)\bigr). \qedhere
\]
\end{proof}

\begin{lemma}[Scalar Pinsker]
\label{lem:dd1-scalar-pinsker}
For all $a, b \in [0, 1]$ (with $\mathrm{kl}(a\|b) = \infty$ when undefined),
\begin{equation}
\label{eq:dd1-scalar-pinsker}
\mathrm{kl}(a \,\|\, b) \;\ge\; 2\,(a - b)^2 .
\end{equation}
\end{lemma}

\begin{proof}
If $b \in \{0, 1\}$ and $a \neq b$ the left side is $+\infty$ and the claim
is trivial; if $a = b$ both sides vanish. So fix $b \in (0,1)$ and define on
$(0,1)$ (extending continuously to the endpoints where finite)
\[
f(a) \;=\; \mathrm{kl}(a \| b) - 2 (a - b)^2
\;=\; a \log\frac{a}{b} + (1 - a)\log\frac{1-a}{1-b} - 2(a-b)^2 .
\]
Then $f(b) = 0$, and differentiating in $a$:
\[
f'(a) = \log \frac{a}{b} - \log \frac{1-a}{1-b} - 4(a - b),
\qquad f'(b) = 0,
\]
\[
f''(a) = \frac{1}{a} + \frac{1}{1 - a} - 4 = \frac{1}{a(1-a)} - 4 \;\ge\; 0,
\]
since $a(1-a) \le 1/4$ for $a \in (0,1)$ (AM--GM: $a(1-a) \le
\bigl(\frac{a + (1-a)}{2}\bigr)^2 = \frac14$). Thus $f'$ is nondecreasing
and vanishes at $a = b$: $f' \le 0$ on $(0, b]$ and $f' \ge 0$ on $[b, 1)$,
so $f$ attains its minimum over $(0,1)$ at $a = b$, where it is $0$. Hence
$f(a) \ge 0$ for all $a$, which is \eqref{eq:dd1-scalar-pinsker}.
\end{proof}

\begin{theorem}[Pinsker's inequality]
\label{thm:dd1-pinsker}
For all probability measures $P, Q$,
\begin{equation}
\label{eq:dd1-pinsker}
\TV(P, Q) \;\le\; \sqrt{ \tfrac{1}{2}\, \KL(P \,\|\, Q) } .
\end{equation}
\end{theorem}

\begin{proof}
If $\KL(P\|Q) = \infty$ there is nothing to prove. Otherwise let
$A^* = \{p \ge q\}$, so that by
Proposition~\ref{prop:dd1-tv-variational}, $\TV(P, Q) = P(A^*) - Q(A^*)$.
Chaining Lemmas~\ref{lem:dd1-dataproc} and \ref{lem:dd1-scalar-pinsker},
\[
\KL(P \| Q)
\;\ge\; \mathrm{kl}\bigl(P(A^*) \,\|\, Q(A^*)\bigr)
\;\ge\; 2\,\bigl(P(A^*) - Q(A^*)\bigr)^2
\;=\; 2\, \TV(P, Q)^2 .
\]
Rearranging gives \eqref{eq:dd1-pinsker}.
\end{proof}

\subsection{The Pinsker Bridge: from Divergence to Regret}

\begin{theorem}[Bounded-loss transfer / Pinsker bridge]
\label{thm:dd1-bridge}
Let $\ell : \Omega \to [0, B]$ be any measurable loss, and write
$L_P = \E_P[\ell]$, $L_Q = \E_Q[\ell]$. Then
\begin{equation}
\label{eq:dd1-bridge}
\bigl| L_P - L_Q \bigr|
\;\le\; B \cdot \TV(P, Q)
\;\le\; B \, \sqrt{ \tfrac{1}{2} \,\KL(P \,\|\, Q) } .
\end{equation}
\end{theorem}

\begin{proof}
The function $g = \ell / B$ takes values in $[0, 1]$, so by the variational
characterization (Proposition~\ref{prop:dd1-tv-variational}),
\[
\bigl| L_P - L_Q \bigr| = B \,\bigl| \E_P[g] - \E_Q[g] \bigr|
\;\le\; B \sup_{g' : 0 \le g' \le 1} \bigl| \E_P[g'] - \E_Q[g'] \bigr|
= B \cdot \TV(P, Q).
\]
The second inequality in \eqref{eq:dd1-bridge} is Pinsker
(Theorem~\ref{thm:dd1-pinsker}).
\end{proof}

\begin{keypoint}
\textbf{Role in the \nm{} theory.} Fix a regime $r$ and apply
Theorem~\ref{thm:dd1-bridge} with: $\Omega$ the space of (feature, return)
days; $\ell(\omega) = \rho(\hat{y}_\theta(x_\omega); y_\omega)$ the decision
regret of niche head $\theta$, bounded by $B = 2 k\, w_{\max}\, y_{\max}$
(Proposition~\ref{prop:dd1-regret-range}); $P$ the mixture of the niche's
stored training episodes; $Q$ a fresh deployment episode of the same regime.
An Inv-PnCO-style assumption \cite{yan2026invpnco}---that episodes of one
regime are divergence-close, $\KL(Q \| P) \le \epsilon_r$ (or
$\TV \le \tau_r$)---then yields a \emph{regret} transfer guarantee:
\[
\text{fresh-episode regret}
\;\le\; \text{stored-episode regret} + B\sqrt{\epsilon_r / 2},
\]
which composes with the concentration bounds of
Section~\ref{sec:dd1-concentration} (controlling the stored-episode average)
to give the episode-transfer lemma of Part II. The bridge is what lets a
statement about \emph{data distributions} (checkable by regime statistics)
control a statement about \emph{money} (regret of induced portfolios), for
\emph{any} head simultaneously---no union bound over heads is needed because
\eqref{eq:dd1-bridge} is deterministic given the divergence.
\end{keypoint}

\begin{remark}[Honest looseness accounting]
\label{rem:dd1-B-loose}
Every appearance of $B = 2 k\, w_{\max}\, y_{\max}$ in
\eqref{eq:dd1-bridge} imports the worst case of
Proposition~\ref{prop:dd1-regret-range}: the forecast shorting itself into
the $k$ worst assets while the oracle rides the $k$ best, all at the extreme
return $y_{\max}$. Empirically observed decision regrets in the companion
experiments range roughly \emph{fifty times smaller} than $B$. The bridge's
\emph{form}---transfer error scaling as
$(\text{regret scale}) \times \sqrt{\text{divergence}}$---is the content we
rely on; its \emph{constant} should be read as a certificate of structure,
not of magnitude. Tightening $B$ to a high-probability regret envelope (e.g.\
a quantile of the realized regret distribution) is possible but converts the
deterministic bridge into a probabilistic one; we keep the loose constant and
flag it.
\end{remark}

% ============================================================================
\section{Online Learning: Hedge, Exponentiated Gradient, and Their Limits}
\label{sec:dd1-online}

The \nm{} gate maintains a probability vector over niches and updates it
multiplicatively in observed decision quality. This update has three classical
pedigrees---Hedge \cite{freund1997decision}, exponentiated gradient
\cite{kivinen1997exponentiated}, and multiplicative weights
\cite{arora2012multiplicative}---and one classical guarantee, the
$O(\sqrt{T \log N})$ adversarial regret bound \cite{cesa2006prediction}. This
section develops all of it, and then states precisely why the classical
guarantee does \emph{not} cover \nm{}'s use of the update.

\subsection{The Experts Problem and the Hedge Update}

\begin{definition}[Prediction with expert advice]
\label{def:dd1-experts}
There are $N$ \textbf{experts}. At each round $t = 1, \dots, T$: the learner
chooses a distribution $p_t \in \Delta_N = \{p \in [0,1]^N : \sum_i p_i =
1\}$; an adversary simultaneously chooses a loss vector
$\ell_t \in [0, 1]^N$; the learner suffers $\langle p_t, \ell_t \rangle$ and
observes all of $\ell_t$. The learner's \textbf{regret} is
\begin{equation}
\label{eq:dd1-experts-regret}
R_T \;=\; \sum_{t=1}^T \langle p_t, \ell_t \rangle
\;-\; \min_{i \in [N]} \sum_{t=1}^T \ell_{t,i},
\end{equation}
the gap to the best \emph{single} expert in hindsight.
\end{definition}

\begin{definition}[Hedge / multiplicative weights]
\label{def:dd1-hedge}
Fix a learning rate $\eta > 0$. Initialize $w_{1,i} = 1$ for all $i$ and
$p_1 = (1/N, \dots, 1/N)$. After observing $\ell_t$,
\begin{equation}
\label{eq:dd1-hedge}
w_{t+1, i} = w_{t, i}\, e^{-\eta\, \ell_{t,i}},
\qquad
p_{t+1, i} = \frac{w_{t+1, i}}{\sum_{j=1}^N w_{t+1, j}}
= \frac{p_{t,i}\, e^{-\eta \ell_{t,i}}}{\sum_j p_{t,j}\, e^{-\eta
\ell_{t,j}}}.
\end{equation}
\end{definition}

This lineage runs through universal portfolio selection
\cite{cover1991universal} and online portfolio optimization
\cite{li2014online}: there, the ``experts'' are constant-rebalanced
portfolios or trading strategies and the weights track wealth shares. In
\nm{} the experts are niche heads and the weights drive both selection and
training-resource allocation---a difference we return to in
Section~\ref{sec:dd1-online-limits}.

\subsection{Exponentiated Gradient as Mirror Descent with Entropy}

The exponential form \eqref{eq:dd1-hedge} is not an ad hoc choice: it is the
\emph{exact} solution of a regularized linear step where the regularizer is
the KL divergence on the simplex \cite{kivinen1997exponentiated}.

\begin{proposition}[EG update solves the KL-regularized linear step]
\label{prop:dd1-eg-mirror}
Fix $p_t$ in the relative interior of $\Delta_N$ and $\eta > 0$. The unique
solution of
\begin{equation}
\label{eq:dd1-mirror-step}
p_{t+1} \;=\; \argmin_{p \in \Delta_N}
\Bigl\{ \eta\, \langle \ell_t, p \rangle + \KL(p \,\|\, p_t) \Bigr\},
\qquad
\KL(p \| p_t) = \sum_{i=1}^N p_i \log \frac{p_i}{p_{t,i}},
\end{equation}
is exactly the Hedge update \eqref{eq:dd1-hedge}.
\end{proposition}

\begin{proof}
The objective $J(p) = \eta \langle \ell_t, p\rangle + \KL(p\|p_t)$ is
strictly convex on $\Delta_N$ ($\langle \ell_t, \cdot\rangle$ is linear and
$\KL(\cdot \| p_t)$ is strictly convex, being a sum of strictly convex scalar
functions $p_i \mapsto p_i \log(p_i / p_{t,i})$), so the minimizer is unique.
Moreover the entropy term's derivative blows up at the boundary
($\partial_{p_i} p_i \log p_i = \log p_i + 1 \to -\infty$ as $p_i \to 0^+$),
which forces the minimizer into the relative interior; hence the
nonnegativity constraints are inactive and we may use a Lagrangian with only
the equality constraint. Set
\[
\mathcal{L}(p, \lambda)
= \sum_{i=1}^N \Bigl[ \eta\, \ell_{t,i}\, p_i
+ p_i \log \frac{p_i}{p_{t,i}} \Bigr]
+ \lambda \Bigl( \sum_{i=1}^N p_i - 1 \Bigr).
\]
Stationarity in each coordinate:
\[
\frac{\partial \mathcal{L}}{\partial p_i}
= \eta\, \ell_{t,i} + \log \frac{p_i}{p_{t,i}} + 1 + \lambda = 0
\quad\Longleftrightarrow\quad
p_i = p_{t,i}\, e^{-\eta \ell_{t,i}}\, e^{-(1 + \lambda)} .
\]
The factor $e^{-(1+\lambda)}$ is a positive constant independent of $i$;
the constraint $\sum_i p_i = 1$ pins it to
$\bigl(\sum_j p_{t,j} e^{-\eta \ell_{t,j}}\bigr)^{-1}$. This is precisely
\eqref{eq:dd1-hedge}.
\end{proof}

\begin{intuition}
The step \eqref{eq:dd1-mirror-step} balances two pulls: ``move mass toward
low-loss coordinates'' (the linear term) against ``do not move far from where
you are, measured \emph{multiplicatively}'' (the KL term). Measuring movement
by KL rather than Euclidean distance is what makes the update respect the
simplex geometry: probabilities can shrink toward zero but never jump across
it, weights never go negative, and a coordinate's \emph{relative} change
$p_{t+1,i}/p_{t,i}$ depends only on its own loss and the normalizer. This is
the precise sense in which the \nm{} gate is a ``well-behaved simplex
ascent'': bounded relative updates, invariance to common loss shifts, and no
projection step that could zero out a niche in one round.
\end{intuition}

\subsection{The Classical Regret Guarantee}

\begin{theorem}[Hedge regret bound \cite{freund1997decision,
cesa2006prediction}]
\label{thm:dd1-hedge-regret}
For losses $\ell_t \in [0,1]^N$ and any $\eta > 0$, Hedge guarantees
\begin{equation}
\label{eq:dd1-hedge-bound}
R_T \;\le\; \frac{\ln N}{\eta} + \frac{\eta\, T}{8},
\end{equation}
and with the tuned rate $\eta = \sqrt{8 \ln N / T}$,
$R_T \le \sqrt{(T/2) \ln N} = O(\sqrt{T \log N})$.
\end{theorem}

\begin{proof}[Proof (potential-function argument)]
Let $W_t = \sum_{i=1}^N w_{t,i}$ with the weights of
Definition~\ref{def:dd1-hedge}, and consider the potential $\ln W_t$.

\emph{Per-round drop.} Since $p_{t,i} = w_{t,i} / W_t$,
\[
\ln \frac{W_{t+1}}{W_t}
= \ln \sum_{i=1}^N \frac{w_{t,i}}{W_t}\, e^{-\eta \ell_{t,i}}
= \ln \E_{i \sim p_t}\bigl[ e^{-\eta \ell_{t,i}} \bigr]
\;\le\; -\eta\, \E_{i \sim p_t}[\ell_{t,i}] + \frac{\eta^2}{8}
= -\eta\, \langle p_t, \ell_t \rangle + \frac{\eta^2}{8},
\]
where the inequality is Hoeffding's lemma
(Lemma~\ref{lem:dd1-hoeffding-lemma}) applied to the random variable
$-\eta\,(\ell_{t,I} - \E\ell_{t,I})$, $I \sim p_t$, supported on an interval
of width $\eta \cdot 1$. Summing over $t$ and telescoping,
\begin{equation}
\label{eq:dd1-potential-upper}
\ln W_{T+1} - \ln W_1
\;\le\; -\eta \sum_{t=1}^T \langle p_t, \ell_t \rangle + \frac{\eta^2 T}{8}.
\end{equation}

\emph{Comparator floor.} For any fixed expert $i^*$, since all weights are
positive,
\begin{equation}
\label{eq:dd1-potential-lower}
\ln W_{T+1} \;\ge\; \ln w_{T+1, i^*}
= \ln\Bigl( 1 \cdot \prod_{t=1}^T e^{-\eta \ell_{t, i^*}} \Bigr)
= -\eta \sum_{t=1}^T \ell_{t, i^*},
\end{equation}
while $\ln W_1 = \ln N$.

\emph{Combine.} Chaining \eqref{eq:dd1-potential-lower} $\le \ln W_{T+1}$
with \eqref{eq:dd1-potential-upper}:
\[
-\eta \sum_t \ell_{t, i^*} - \ln N
\;\le\; -\eta \sum_t \langle p_t, \ell_t \rangle + \frac{\eta^2 T}{8}
\;\;\Longrightarrow\;\;
\sum_t \langle p_t, \ell_t \rangle - \sum_t \ell_{t, i^*}
\;\le\; \frac{\ln N}{\eta} + \frac{\eta T}{8}.
\]
Taking $i^*$ to be the best expert gives \eqref{eq:dd1-hedge-bound};
substituting the tuned $\eta$ gives
$2\sqrt{(\ln N)(T)/8} = \sqrt{(T/2)\ln N}$.
\end{proof}

Remarkably, the bound holds against an adversary who sees the learner's
distribution before choosing losses---but, crucially, an adversary whose
losses are \emph{handed to fixed experts}, a point we scrutinize below.

\subsection{The Replicator Limit}

For small $\eta$, the discrete multiplicative update becomes the replicator
dynamics of evolutionary game theory---the lens through which the
winner-take-all niche competition of \nm{} (and of GAUSE before it
\cite{li2026gause}) is analyzed.

\begin{proposition}[Small-$\eta$ replicator limit]
\label{prop:dd1-replicator}
For the Hedge update \eqref{eq:dd1-hedge} with bounded losses,
\begin{equation}
\label{eq:dd1-replicator}
p_{t+1, i} - p_{t, i}
\;=\; \eta\, p_{t,i} \bigl( \langle p_t, \ell_t \rangle - \ell_{t,i} \bigr)
\;+\; O(\eta^2),
\end{equation}
so as $\eta \to 0$ the rescaled-time trajectories solve the replicator ODE
$\dot{p}_i = p_i (\bar{f} - f_i)$ with ``fitness'' $f_i = \ell_{t,i}$ (lower
loss = higher fitness, $\bar{f} = \langle p, \ell\rangle$ the population
mean).
\end{proposition}

\begin{proof}
Expand the exponentials to first order: $e^{-\eta \ell_{t,i}} = 1 - \eta
\ell_{t,i} + O(\eta^2)$, uniformly in $i$ by boundedness. The normalizer is
\[
\sum_j p_{t,j} e^{-\eta \ell_{t,j}}
= 1 - \eta \langle p_t, \ell_t \rangle + O(\eta^2),
\qquad\text{so}\qquad
\Bigl( \sum_j p_{t,j} e^{-\eta \ell_{t,j}} \Bigr)^{-1}
= 1 + \eta \langle p_t, \ell_t \rangle + O(\eta^2).
\]
Multiplying,
\[
p_{t+1,i} = p_{t,i} \bigl(1 - \eta \ell_{t,i}\bigr)
\bigl(1 + \eta \langle p_t, \ell_t\rangle\bigr) + O(\eta^2)
= p_{t,i}\Bigl[ 1 + \eta\bigl( \langle p_t, \ell_t \rangle - \ell_{t,i}
\bigr) \Bigr] + O(\eta^2). \qedhere
\]
\end{proof}

\begin{intuition}
A niche grows in gate weight iff it beats the \emph{current
weighted average} of all niches; the growth rate is proportional to both its
advantage and its current weight. This is exactly the dynamics under which
competitive exclusion and niche differentiation are usually derived---and it
is a \emph{dynamical-systems} object (fixed points, basins, exclusion), not a
regret object.
\end{intuition}

\subsection{What the Classical Guarantee Does Not Cover}
\label{sec:dd1-online-limits}

\begin{warning}
\textbf{The adversarial regret bound does not transfer to \nm{}.}
Theorem~\ref{thm:dd1-hedge-regret} compares the learner to the best
\emph{fixed} expert under loss sequences that are \emph{exogenous}: expert
$i$'s loss at round $t$ does not depend on how much weight the learner has
been giving expert $i$. In \nm{} both premises fail, for the same reason they
fail in GAUSE \cite{li2026gause}:
\begin{enumerate}
\item \emph{Losses are endogenous.} Niche $i$'s decision quality at round $t$
depends on its head parameters, which were trained with data and gradient
budget allocated \emph{in proportion to past gate weights}. The loss sequence
$\ell_{t,i}$ is a functional of the learner's own trajectory
$p_{1:t-1}$---outside the oblivious- and even the standard adaptive-adversary
model, where ``expert $i$'s loss'' must remain well-defined independently of
the comparison.
\item \emph{The comparator is meaningless.} ``The best fixed niche in
hindsight'' is not a coherent benchmark when the niche's identity (its
trained parameters) is itself an artifact of the weight history; under a
different weight history the ``same'' niche would be a different predictor.
A regret bound against an incoherent comparator is vacuous even if formally
true.
\item \emph{Winner-take-all coupling.} The \nm{} assignment concentrates
training on high-weight niches, which is the very mechanism that creates
specialization; this positive feedback is what the replicator analysis
(Proposition~\ref{prop:dd1-replicator}) studies and what a worst-case regret
analysis is designed to neutralize.
\end{enumerate}
Accordingly, \nm{} uses the EG/Hedge update for its \emph{structural}
virtues established in Proposition~\ref{prop:dd1-eg-mirror} and the
intuition following it---a strictly interior, multiplicatively bounded,
geometry-respecting ascent on the simplex---and \emph{not} for the bound
\eqref{eq:dd1-hedge-bound}. Every stability claim about the coupled dynamics
is proved separately in Part II; none is inherited from
Theorem~\ref{thm:dd1-hedge-regret}.
\end{warning}

\subsection{Sleeping Experts and Adaptive Regret: the Gap \nm{} Addresses}
\label{sec:dd1-sleeping}

Non-stationarity has its own classical formalisms, and it is important to
locate precisely where they stop short of the \nm{} setting.

\begin{definition}[Sleeping experts \cite{freund1997sleeping}]
\label{def:dd1-sleeping}
In the sleeping-experts (specialists) problem, at each round $t$ an
adversary additionally reveals an \textbf{awake set} $A_t \subseteq [N]$;
only awake experts make predictions, the learner must put its mass on $A_t$,
and asleep experts are unevaluated. The standard guarantee is per-expert: for
every expert $i$,
\begin{equation}
\label{eq:dd1-sleeping-regret}
\sum_{t \,:\, i \in A_t} \langle p_t, \ell_t \rangle
\;-\; \sum_{t \,:\, i \in A_t} \ell_{t,i}
\;\le\; O\bigl( \sqrt{ T_i \,\log N } \bigr),
\qquad T_i = |\{ t : i \in A_t \}|,
\end{equation}
i.e.\ the learner competes with each specialist \emph{on the rounds where
that specialist is awake}.
\end{definition}

\begin{definition}[Adaptive and strongly adaptive regret
\cite{hazan2009adaptive, daniely2015strongly}]
\label{def:dd1-adaptive}
The \textbf{adaptive regret} of an algorithm is its worst regret over any
contiguous interval,
\begin{equation}
\label{eq:dd1-adaptive-regret}
\mathrm{AR}_T \;=\; \max_{1 \le s \le e \le T}
\Bigl[ \sum_{t=s}^{e} \langle p_t, \ell_t \rangle
- \min_{i} \sum_{t=s}^{e} \ell_{t,i} \Bigr];
\end{equation}
an algorithm is \textbf{strongly adaptive} \cite{daniely2015strongly} if for
every interval length $\tau$ its regret on every length-$\tau$ interval is
$O(\mathrm{poly}(\log T) \cdot \sqrt{\tau})$. Such algorithms track the best
expert \emph{per regime segment} rather than globally.
\end{definition}

\begin{keypoint}
\textbf{Why dormant \emph{fixed} experts retain trivially.} In both
frameworks an expert is a fixed function (or a fixed external advisor). When
expert $i$ sleeps for a thousand rounds, nothing happens to it: its
parameters do not exist or do not change, and when it wakes its predictions
are exactly what they would have been. Retention of dormant competence is
\emph{free}---it is an assumption, not an achievement. The guarantees
\eqref{eq:dd1-sleeping-regret} and \eqref{eq:dd1-adaptive-regret} are
therefore statements purely about the \emph{weighting} layer: how fast the
learner re-concentrates on the right expert after a regime switch.
\end{keypoint}

\begin{remark}[The precise gap]
\label{rem:dd1-gap}
\nm{} operates in a setting these formalisms do not model:
\textbf{capacity-bounded \emph{learning} experts under reward-driven
training allocation}. Formally: each expert $i$ carries parameters
$\theta_{t,i}$ in a bounded-capacity class; at each round a shared training
budget is divided according to (a function of) the gate weights $p_t$, and
$\theta_{t+1, i}$ is produced by gradient steps on data weighted by expert
$i$'s allocation; the loss $\ell_{t,i} = \ell(\theta_{t,i};\,
\text{episode}_t)$ is thereby a functional of the entire weight history. In
this setting dormancy is no longer safe by assumption: an expert that
receives weight (hence budget, hence gradient updates on
\emph{out-of-niche} data, or no updates while the shared feature extractor
drifts) can have its dormant competence \emph{overwritten}---the
continual-learning failure mode studied for mixture-of-experts learners in
\cite{li2024moecl} and in the GAUSE retention analysis \cite{li2026gause}.
The sleeping-experts guarantee says nothing here because the premise ``the
specialist's awake-round losses are fixed targets'' is false; strongly
adaptive algorithms recover the right \emph{weights} quickly but have no
mechanism to ensure there remains a competent expert to weight. \nm{}'s
contribution lives exactly in this gap: the niche memory and the
variance-penalized per-niche objective are mechanisms for making the
``expert exists and is still competent when its regime returns'' premise
\emph{true}, after which adaptive-weighting logic can do its classical job.
\end{remark}

% ============================================================================
\section{Invariant Risk and the Failure of ERM Across Environments}
\label{sec:dd1-invariance}

The last foundation explains \emph{what is being learned inside each niche}.
A niche head sees a handful of episodes of ``its'' regime; those episodes
share the regime's stable structure but differ in incidental, episode-specific
correlations. Pooled empirical risk minimization (ERM) over the episodes is
the default---and it fails in a precise, derivable way: it loads on the
incidental correlations exactly in proportion to how unbalanced the small
sample of episodes happens to be. This section sets up the standard
two-feature construction, derives the failure quantitatively, shows that the
\emph{across-episode variance of per-episode loss} is a statistic that detects
exactly the failure, and positions the \nm{} per-niche objective in the
invariant-risk literature \cite{arjovsky2019irm, krueger2021rex,
peters2016causal}.

\subsection{Environments, Episodes, and Pooled ERM}

\begin{definition}[Multi-environment learning]
\label{def:dd1-environments}
Data arrive grouped into \textbf{episodes} (environments)
$e \in \mathcal{E}$; episode $e$ carries a distribution $P_e$ over
observation--target pairs. A training set consists of samples from $E$
episodes $e_1, \dots, e_E$ drawn from an environment distribution
$\mathcal{Q}$; at deployment a \emph{fresh} episode $e' \sim \mathcal{Q}$ is
encountered. For a predictor $h$ and loss $\ell$, write
$L_e(h) = \E_{P_e}[\ell(h(x), y)]$ for the per-episode risk. \textbf{Pooled
ERM} minimizes $\frac{1}{E}\sum_{j=1}^E L_{e_j}(h)$ (empirically); the
\textbf{invariant} ideal is a predictor whose risk---and whose optimality---is
stable across all $e$ in the support of $\mathcal{Q}$.
\end{definition}

In \nm{}, $\mathcal{E}$ is the set of episodes routed to one niche, $E = E_r
\in \{5, \dots, 7\}$, and $\ell$ is the SPO+ surrogate of
Section~\ref{sec:dd1-spoplus}; for the analysis of this section we use the
squared loss, where everything is exactly computable, and return to the SPO+
case at the end.

\subsection{The Two-Feature Construction}

\begin{definition}[Invariant vs.\ spurious feature model]
\label{def:dd1-twofeature}
Within episode $e$, independent draws are generated as
\begin{equation}
\label{eq:dd1-sem}
x_{\mathrm{inv}} \sim \N(0, 1), \qquad
\varepsilon \sim \N(0, \sigma_\varepsilon^2), \qquad
y = \theta\, x_{\mathrm{inv}} + \varepsilon, \qquad
x_{\mathrm{sp}} = \gamma_e\, y + \nu, \quad \nu \sim \N(0, \sigma_\nu^2),
\end{equation}
with $x_{\mathrm{inv}}, \varepsilon, \nu$ mutually independent. The
\textbf{invariant} coefficient $\theta$ is the same in every episode; the
\textbf{spurious} (anti-causal) coefficient $\gamma_e$ is drawn once per
episode, i.i.d.\ across episodes with $\E[\gamma] = 0$ and
$\Var[\gamma] = \sigma_\gamma^2 > 0$. Write $s^2 := \Var(y) = \theta^2 +
\sigma_\varepsilon^2$. A linear predictor is
$h_{a,b}(x) = a\, x_{\mathrm{inv}} + b\, x_{\mathrm{sp}}$, and the
\textbf{invariant predictor} is $h_{\theta, 0}$.
\end{definition}

\begin{intuition}
$x_{\mathrm{inv}}$ is a genuine driver of $y$ (a stable regime
characteristic); $x_{\mathrm{sp}}$ is a symptom---it is generated \emph{from}
$y$, with an episode-specific gain $\gamma_e$ (an incidental correlation:
e.g.\ a factor that happened to co-move with returns during those particular
weeks). Because $x_{\mathrm{sp}}$ is correlated with $y$ within each episode,
regression happily uses it; because the correlation's sign and size are
episode-specific accidents, doing so buys in-sample fit with out-of-episode
fragility. Averaging the $\gamma_e$'s over infinitely many episodes would
cancel them ($\E\gamma = 0$); the problem is that a niche has $5$--$7$
episodes, and the empirical mean $\bar{\gamma} = \frac{1}{E}\sum_j
\gamma_{e_j}$ is a $\Theta(\sigma_\gamma E^{-1/2})$-sized accident that does
\emph{not} cancel.
\end{intuition}

\begin{proposition}[Pooled-OLS limit loads on the spurious feature]
\label{prop:dd1-pooled-ols}
Consider pooled least squares over $E$ training episodes with equal weights
and infinite within-episode samples (so all second moments are exact). Write
$\bar{\gamma} = \frac{1}{E}\sum_{j} \gamma_{e_j}$,
$\widehat{m}_2 = \frac{1}{E}\sum_j \gamma_{e_j}^2$,
$\widehat{\sigma}_\gamma^2 = \widehat{m}_2 - \bar{\gamma}^2$. The pooled OLS
coefficients $(\widehat{a}, \widehat{b}) = \argmin_{a, b} \frac{1}{E}
\sum_j L_{e_j}(h_{a,b})$ are
\begin{equation}
\label{eq:dd1-pooled-limit}
\widehat{b}
= \frac{\bar{\gamma}\, \sigma_\varepsilon^2}{D},
\qquad
\widehat{a}
= \theta \cdot \frac{\sigma_\nu^2 + s^2\, \widehat{\sigma}_\gamma^2}{D},
\qquad
D := \sigma_\nu^2 + s^2\, \widehat{\sigma}_\gamma^2
+ \bar{\gamma}^2 \sigma_\varepsilon^2 \;>\; 0.
\end{equation}
In particular $\widehat{b} \neq 0$ whenever $\bar{\gamma} \neq 0$: the pooled
solution loads on $x_{\mathrm{sp}}$ \emph{in proportion to the empirical mean
of the episode coefficients}, which for small $E$ is nonzero with
probability one (for continuous $\gamma$) and of typical size
$\sigma_\gamma / \sqrt{E}$.
\end{proposition}

\begin{proof}
Within episode $e$, using \eqref{eq:dd1-sem} and independence:
\[
\E[x_{\mathrm{inv}}^2] = 1, \quad
\E[x_{\mathrm{inv}} x_{\mathrm{sp}}] = \gamma_e \E[x_{\mathrm{inv}} y]
= \gamma_e \theta, \quad
\E[x_{\mathrm{sp}}^2] = \gamma_e^2 s^2 + \sigma_\nu^2, \quad
\E[x_{\mathrm{inv}} y] = \theta, \quad
\E[x_{\mathrm{sp}} y] = \gamma_e s^2 .
\]
Pooling with equal weights averages these moments over the $E$ training
episodes, so the pooled normal equations
$\bar{\Sigma} (\widehat a, \widehat b)^\top = \bar v$ have
\[
\bar{\Sigma} = \begin{pmatrix} 1 & \bar{\gamma}\theta \\
\bar{\gamma}\theta & \widehat{m}_2 s^2 + \sigma_\nu^2 \end{pmatrix},
\qquad
\bar v = \begin{pmatrix} \theta \\ \bar{\gamma} s^2 \end{pmatrix}.
\]
The determinant is
\[
\det \bar{\Sigma}
= \widehat{m}_2 s^2 + \sigma_\nu^2 - \bar{\gamma}^2 \theta^2
= \sigma_\nu^2 + s^2(\widehat{m}_2 - \bar{\gamma}^2)
+ \bar{\gamma}^2 (s^2 - \theta^2)
= \sigma_\nu^2 + s^2 \widehat{\sigma}_\gamma^2
+ \bar{\gamma}^2 \sigma_\varepsilon^2 = D > 0 .
\]
By Cramer's rule,
\[
\widehat{b} = \frac{1 \cdot \bar{\gamma} s^2 - \bar{\gamma}\theta \cdot
\theta}{D} = \frac{\bar{\gamma}(s^2 - \theta^2)}{D}
= \frac{\bar{\gamma}\, \sigma_\varepsilon^2}{D},
\qquad
\widehat{a} = \frac{(\widehat{m}_2 s^2 + \sigma_\nu^2)\theta -
\bar{\gamma}\theta \cdot \bar{\gamma} s^2}{D}
= \frac{\theta\bigl( \sigma_\nu^2 + s^2 \widehat{\sigma}_\gamma^2
\bigr)}{D}. \qedhere
\]
\end{proof}

Note the mechanism: $\widehat{b} \propto \bar{\gamma}$, an average of
mean-zero accidents that a law of large numbers would kill---\emph{if $E$
were large}. ERM is not wrong asymptotically; it is wrong at $E = 5$--$7$,
which is the only regime \nm{} lives in. Pooled ERM also genuinely
\emph{prefers} this loading in-sample:

\begin{lemma}[ERM strictly descends into the spurious direction]
\label{lem:dd1-erm-descends}
The per-episode risk of $h_{a,b}$ is
\begin{equation}
\label{eq:dd1-Le}
L_e(h_{a,b})
= \bigl( \theta - a - b\gamma_e\theta \bigr)^2
+ \bigl( 1 - b\gamma_e \bigr)^2 \sigma_\varepsilon^2
+ b^2 \sigma_\nu^2,
\end{equation}
and at the invariant predictor $(a, b) = (\theta, 0)$ the pooled risk has
$\frac{\partial}{\partial b}\, \frac{1}{E}\sum_j L_{e_j} =
-2\bar{\gamma}\sigma_\varepsilon^2$. Hence whenever $\bar{\gamma} \neq 0$,
moving $b$ slightly in the direction of $\operatorname{sign}(\bar\gamma)$
strictly reduces pooled training risk below the invariant predictor's.
\end{lemma}

\begin{proof}
The residual is $y - a x_{\mathrm{inv}} - b x_{\mathrm{sp}}
= (\theta - a - b\gamma_e \theta) x_{\mathrm{inv}} + (1 - b\gamma_e)
\varepsilon - b \nu$ (substitute $x_{\mathrm{sp}} = \gamma_e(\theta
x_{\mathrm{inv}} + \varepsilon) + \nu$); squaring and using independence and
unit/known variances gives \eqref{eq:dd1-Le}. Differentiate in $b$:
\[
\frac{\partial L_e}{\partial b}
= -2\gamma_e\theta\,(\theta - a - b\gamma_e\theta)
- 2\gamma_e \sigma_\varepsilon^2 (1 - b\gamma_e) + 2 b \sigma_\nu^2
\;\xrightarrow{\;(a,b) = (\theta, 0)\;}\;
0 - 2\gamma_e \sigma_\varepsilon^2 + 0 .
\]
Averaging over the training episodes gives $-2\bar\gamma
\sigma_\varepsilon^2$, nonzero iff $\bar\gamma \neq 0$.
\end{proof}

\begin{proposition}[Expected excess risk on a fresh episode]
\label{prop:dd1-excess}
For any fixed predictor $h_{a,b}$ and a fresh episode $e' \sim \mathcal{Q}$
(i.e.\ $\gamma_{e'}$ a fresh draw with mean $0$ and variance
$\sigma_\gamma^2$, independent of the training episodes),
\begin{equation}
\label{eq:dd1-excess}
\E_{\gamma_{e'}}\bigl[ L_{e'}(h_{a,b}) \bigr] - L_{e'}(h_{\theta, 0})
= (\theta - a)^2 \;+\; b^2 \bigl( \sigma_\nu^2 + s^2 \sigma_\gamma^2 \bigr)
\;\ge\; 0,
\end{equation}
with equality iff $(a,b) = (\theta, 0)$; note $L_{e'}(h_{\theta,0}) =
\sigma_\varepsilon^2$ deterministically. For the pooled-OLS solution
\eqref{eq:dd1-pooled-limit}, the leading term for small $|\bar\gamma|$ is
\begin{equation}
\label{eq:dd1-excess-rate}
\widehat{b}^{\,2}\bigl( \sigma_\nu^2 + s^2\sigma_\gamma^2 \bigr)
= \frac{\bar{\gamma}^2\, \sigma_\varepsilon^4 \,(\sigma_\nu^2 +
s^2\sigma_\gamma^2)}{D^2},
\qquad
\E\bigl[\bar{\gamma}^2\bigr] = \frac{\sigma_\gamma^2}{E},
\end{equation}
while $(\theta - \widehat a)^2 = O(\bar\gamma^4)$. The expected damage of
pooling thus decays only as $1/E$---negligible for hundreds of episodes,
material for $E = 5$--$7$.
\end{proposition}

\begin{proof}
Take the expectation of \eqref{eq:dd1-Le} over $\gamma_{e'}$ with
$\E\gamma_{e'} = 0$, $\E\gamma_{e'}^2 = \sigma_\gamma^2$:
\[
\E\bigl[(\theta - a - b\gamma_{e'}\theta)^2\bigr]
= (\theta - a)^2 + b^2\theta^2\sigma_\gamma^2,
\qquad
\E\bigl[(1 - b\gamma_{e'})^2\bigr]\sigma_\varepsilon^2
= \sigma_\varepsilon^2 + b^2 \sigma_\varepsilon^2 \sigma_\gamma^2 .
\]
Summing with $b^2\sigma_\nu^2$ and subtracting
$L_{e'}(h_{\theta,0}) = \sigma_\varepsilon^2$ (immediate from
\eqref{eq:dd1-Le} at $(\theta, 0)$) gives \eqref{eq:dd1-excess}, using
$\theta^2 + \sigma_\varepsilon^2 = s^2$. For the rate: $\widehat b =
\bar\gamma \sigma_\varepsilon^2 / D$ from
Proposition~\ref{prop:dd1-pooled-ols}, and $\theta - \widehat a =
\theta \bar\gamma^2 \sigma_\varepsilon^2 / D$ (subtract
\eqref{eq:dd1-pooled-limit} from $\theta$), whose square is $O(\bar\gamma^4)$.
Finally $\E[\bar\gamma^2] = \Var(\bar\gamma) = \sigma_\gamma^2 / E$ for
i.i.d.\ mean-zero $\gamma_{e_j}$.
\end{proof}

\begin{example}[Numbers at niche scale]
\label{ex:dd1-pooled-numbers}
Take $\theta = 1$, $\sigma_\varepsilon^2 = 0.5$ (so $s^2 = 1.5$),
$\sigma_\nu^2 = 1$, and $E = 6$ training episodes with
$\gamma = (0.4, -0.2, 0.3, -0.1, 0.2, 0.0)$---a sample from a mean-zero law
that happens, as small samples do, to lean positive: $\bar\gamma = 0.1$,
$\widehat{m}_2 = 0.34/6 \approx 0.0567$, $\widehat\sigma_\gamma^2 \approx
0.0467$. Then $D = 1 + 1.5 \times 0.0467 + 0.01 \times 0.5 \approx 1.075$ and
\[
\widehat{b} = \frac{0.1 \times 0.5}{1.075} \approx 0.047, \qquad
\widehat{a} \approx 1 \times \frac{1.070}{1.075} \approx 0.995 .
\]
With true $\sigma_\gamma^2 = 0.05$, the expected fresh-episode excess
\eqref{eq:dd1-excess} is approximately $0.047^2 (1 + 1.5 \times 0.05) +
0.005^2 \approx 0.0023$---about $0.5\%$ of the irreducible risk
$\sigma_\varepsilon^2 = 0.5$ per fresh episode, incurred systematically, with
the sign of the error flipping adversarially on episodes where
$\gamma_{e'}$ opposes $\bar\gamma$.
\end{example}

\subsection{The Variance Statistic Detects Exactly the Loading}

The decisive structural fact about the construction is visible in
\eqref{eq:dd1-Le}: the per-episode risk depends on the episode \emph{only
through the products $b\gamma_e$}. A predictor that ignores the spurious
feature has the \emph{same} risk in every episode; a predictor that uses it
has episode-dependent risk. Across-episode variance is therefore a
\emph{detector} of spurious loading.

\begin{proposition}[Variance identifies invariance]
\label{prop:dd1-variance-detects}
Let $\Var_e[\,\cdot\,]$ denote the empirical variance across the training
episodes $e_1, \dots, e_E$, and suppose at least two of the
$\gamma_{e_j}$ are distinct. Then for the model
\eqref{eq:dd1-sem}:
\begin{enumerate}
\item If $b = 0$, then $L_{e_j}(h_{a, 0}) = (\theta - a)^2 +
\sigma_\varepsilon^2$ for every $j$, so $\Var_e\bigl[ L_e(h_{a,0}) \bigr] =
0$.
\item If $b \neq 0$, then $\Var_e\bigl[ L_e(h_{a,b}) \bigr] > 0$, except for
at most finitely many exceptional values of $(a, b)$ (those making the
quadratic $\gamma \mapsto L(\gamma)$ in \eqref{eq:dd1-Le} take equal values
at all the distinct observed $\gamma_{e_j}$, possible only when fewer than
three distinct values are observed and $(a,b)$ is tuned to the pair).
With three or more distinct $\gamma_{e_j}$, $b \neq 0$ implies
$\Var_e[L_e] > 0$ without exception.
\end{enumerate}
Consequently the population of zero-variance predictors is exactly the
$b = 0$ axis (generically), on which the pooled mean risk is minimized at
the invariant predictor $a = \theta$: minimizing
$\Var_e[L_e] + \beta\, \E_e[L_e]$ with $\beta > 0$ and a sufficiently large
variance weight recovers $h_{\theta, 0}$, whereas pooled ERM
($\beta$-only) strictly prefers $\widehat b \neq 0$ by
Lemma~\ref{lem:dd1-erm-descends}.
\end{proposition}

\begin{proof}
(1) Set $b = 0$ in \eqref{eq:dd1-Le}: every $\gamma_e$-dependent term carries
a factor $b$, leaving $(\theta - a)^2 + \sigma_\varepsilon^2$, constant
across episodes; the empirical variance of a constant list is $0$.

(2) Regard \eqref{eq:dd1-Le} as a polynomial in $\gamma$:
\[
L(\gamma) = \underbrace{b^2\theta^2 \,\gamma^2 - 2b\theta(\theta - a)\gamma
+ (\theta - a)^2}_{\text{first term}}
+ \underbrace{\sigma_\varepsilon^2 b^2 \gamma^2 - 2\sigma_\varepsilon^2 b
\gamma + \sigma_\varepsilon^2}_{\text{second term}} + b^2\sigma_\nu^2,
\]
with leading coefficient $b^2(\theta^2 + \sigma_\varepsilon^2) = b^2 s^2 > 0$
when $b \neq 0$: a nondegenerate quadratic. The empirical variance vanishes
iff $L(\gamma_{e_j})$ is the same number for all $j$. A nondegenerate
quadratic takes any given value at most twice, so if three or more distinct
$\gamma_{e_j}$ are observed, equality of all values is impossible and
$\Var_e[L_e] > 0$. With exactly two distinct values $\gamma \neq \gamma'$,
equality forces the quadratic's axis of symmetry to bisect them, i.e.\ pins
one algebraic relation between $a$ and $b$---the finitely-many-exceptions
caveat. The final claims restate Lemma~\ref{lem:dd1-erm-descends} and part
(1) together with the strict positivity in part (2): for large variance
weight, any $b \neq 0$ is dominated by its $b = 0$ projection, and among
$b = 0$ predictors the mean term selects $a = \theta$.
\end{proof}

\begin{keypoint}
The variance statistic is not a generic regularizer; in this construction it
is a \emph{consistent test} for the precise pathology of
Proposition~\ref{prop:dd1-pooled-ols}. Pooled ERM cannot see the difference
between fitting structure and fitting accidents, because both reduce mean
training risk (Lemma~\ref{lem:dd1-erm-descends}); the across-episode variance
sees \emph{only} the accidents, because structure---by definition of being
shared across the niche's episodes---contributes nothing to it
(Proposition~\ref{prop:dd1-variance-detects}). This is also exactly the
quantity the small-sample transfer story of
Section~\ref{sec:dd1-concentration} wants minimized: a smaller
across-episode spread is what would tighten an empirical-Bernstein-style
certificate (Remark~\ref{rem:dd1-small-n}) and what shrinks the variance in
Cantelli's one-sided bound (Theorem~\ref{thm:dd1-cantelli}).
\end{keypoint}

\subsection{Relation to IRM, REx, and Causal Invariance; the \nm{} Objective}

Three strands of the invariance literature meet here.
\textbf{Invariant causal prediction} \cite{peters2016causal} formalizes the
target: under interventions that vary across environments, the conditional
distribution of the response given its \emph{causal parents} is invariant,
and (under conditions) the causal feature set is identifiable as the one
producing invariant residuals. In Definition~\ref{def:dd1-twofeature},
$x_{\mathrm{inv}}$ is a parent and $x_{\mathrm{sp}}$ a child of $y$; the
invariant predictor is the causal one.
\textbf{IRM} \cite{arjovsky2019irm} pursues this target through a bilevel
program---learn a representation such that one classifier on top of it is
simultaneously per-environment optimal---in practice relaxed to a gradient
penalty $\sum_e \|\nabla_{w \mid w = 1} L_e(w \cdot h)\|^2$.
\textbf{REx} \cite{krueger2021rex} observes that a simpler scalarization
often suffices and is better behaved: penalize the spread of per-environment
risks directly, minimizing
$\mathrm{Var}_e[L_e(h)] + \beta\, \E_e[L_e(h)]$-type objectives (V-REx);
robustness to the spread of risks extrapolates to environments whose
``accidents'' are larger than any seen in training.

The \nm{} per-niche objective is the V-REx scalarization applied to a
\emph{decision-focused} episode risk. For niche $r$ with stored episodes
$\mathcal{E}_r$ and head parameters $\theta_r$, define the SPO+ episode risk
$R_e(\theta_r) = \frac{1}{|e|}\sum_{(x, y) \in e}
\ell_{\mathrm{SPO+}}\bigl( \hat y_{\theta_r}(x);\, y \bigr)$ and train by
\begin{equation}
\label{eq:dd1-nm-objective}
J_r(\theta_r)
\;=\; \Var_{e \in \mathcal{E}_r}\bigl[ R_e(\theta_r) \bigr]
\;+\; \beta\; \E_{e \in \mathcal{E}_r}\bigl[ R_e(\theta_r) \bigr],
\qquad \beta > 0 .
\end{equation}
Two substitutions distinguish \eqref{eq:dd1-nm-objective} from textbook REx.
First, the risk inside the variance is the SPO+ \emph{decision} risk, so
``invariance'' is demanded of the quantity that matters
(Section~\ref{sec:dd1-dfl}): a head may have episode-varying \emph{forecast}
error and still be exactly invariant in decision quality, and the objective
correctly declines to punish it. Second, the environments are not given
exogenously but are the episodes routed to the niche---the regime-indexed
environment structure of the invariant predict-and-optimize formulation
\cite{yan2026invpnco}---so the variance is computed within a (putative)
regime, where invariance is actually achievable, rather than across all of
history, where it is not.

\begin{warning}
\textbf{Variance alone ($\beta = 0$) is degenerate.} The zero-variance set
contains far more than invariant-and-good predictors: \emph{any} head whose
episode risks are constant achieves $J_r = 0$ at $\beta = 0$, including
uniformly terrible ones. Concretely, a head whose forecasts induce the same
fixed bad portfolio in all states has $R_e \equiv$ (the same poor value) on
every episode---zero variance, maximal complacency; in the squared-loss model
of Proposition~\ref{prop:dd1-variance-detects}, the entire plane $b = 0$ is
zero-variance regardless of how wrong $a$ is. The mean term is therefore not
a tuning nicety but a load-bearing component: it is the only part of
\eqref{eq:dd1-nm-objective} that distinguishes invariant-good from
invariant-bad. This degeneracy predicts, qualitatively and in advance, the
$\beta = 0$ collapse observed in the companion experiments: post-reactivation
decision regret $0.0356$ for the variance-only objective versus $0.0321$ for
plain retained ERM---the ablation does not merely lose the invariance
benefit, it underperforms doing nothing clever at all, exactly as a
degenerate objective should.
\end{warning}

\begin{remark}[What carries over from squared loss to SPO+]
\label{rem:dd1-carryover}
The exact algebra of
Propositions~\ref{prop:dd1-pooled-ols}--\ref{prop:dd1-variance-detects} is
special to squared loss, where risks are quadratics in $\gamma_e$. What
transfers to the SPO+ risk is the structural trichotomy, none of whose
ingredients used quadraticity: (i) pooled ERM's first-order preference for
episode-specific signal scales with the \emph{empirical mean} of the
accident process and is therefore $\Theta(E^{-1/2})$-sized at niche scale;
(ii) any loading on episode-specific signal makes the per-episode risk
\emph{episode-dependent}, hence visible to the across-episode variance;
(iii) zero variance plus low mean jointly---and only jointly---certify
``invariantly good.'' These three facts, with SPO+ risks in place of squared
risks and with the boundedness $R_e$-analogues supplied by
Proposition~\ref{prop:dd1-regret-range}, are the form in which Part II uses
this section.
\end{remark}

This completes the toolkit. Part II assembles it: the decision layer and
SPO+ machinery of Section~\ref{sec:dd1-dfl} define what each niche head
optimizes \eqref{eq:dd1-nm-objective}; the EG gate of
Section~\ref{sec:dd1-online} allocates episodes and training resources to
niches, deliberately forgoing the adversarial guarantee
(Section~\ref{sec:dd1-online-limits}) in favor of analyzable replicator
dynamics; and the retention and transfer theorems combine the concentration
inequalities of Section~\ref{sec:dd1-concentration} with the Pinsker bridge
of Section~\ref{sec:dd1-distances} to state precisely when a dormant niche's
certified competence survives until its regime returns---the capability gap,
identified in Remark~\ref{rem:dd1-gap}, that the sleeping-experts literature
leaves open and \nm{} is built to close.
% ============================================================================
% deepdive_part23.tex -- Parts II and III of the RISP mathematical
% deep dive. This is a FRAGMENT: it assumes the master preamble defines
% theorem/proposition/lemma/corollary/definition/example/remark (numbered by
% section), the tcolorbox environments intuition/keypoint/warning, and the
% macros \R \E \Var \KL \TV \nm \argmax \argmin.
% All labels are prefixed sec:dd2- / eq:dd2- / tab:dd2-.
% ============================================================================

% ============================================================================
% PART II: THE ENVIRONMENT
% ============================================================================
\part{The Environment: Regimes, Episodes, and Dormancy}
\label{sec:dd2-part-env}

The mechanism documented in Part~III is an answer to a question posed
entirely by the environment: regimes that arrive in persistent blocks
(Section~\ref{sec:dd2-regimes}), episodes as the true statistical unit
(Section~\ref{sec:dd2-episodes}), and the combinatorics of dormancy, during
which a capacity-limited learner is quietly destroyed
(Section~\ref{sec:dd2-dormancy}). Every design choice in \nm{} is a response
to a quantitative fact established in this part.

\section{Regime-Switching Markets}
\label{sec:dd2-regimes}

\subsection{The Regime Process}

\begin{definition}[Block-structured regime process]
\label{def:dd2-regime-process}
Fix a finite regime set $\mathcal{R} = \{1, \ldots, R\}$. A
\textbf{block-structured (semi-Markov) regime process} is a sequence of pairs
\[
(\rho_1, \ell_1),\ (\rho_2, \ell_2),\ (\rho_3, \ell_3),\ \ldots
\qquad \rho_b \in \mathcal{R},\ \ell_b \in \{1, 2, \ldots\},
\]
where $\rho_b$ is the regime of the $b$-th \emph{block} and $\ell_b$ its
duration in days, with $\rho_{b+1} \neq \rho_b$ (blocks are maximal). The
day-level regime label $r_t$ is the piecewise-constant expansion: writing
$T_b = \sum_{b' < b} \ell_{b'}$ for the start of block $b$,
\begin{equation}
\label{eq:dd2-semimarkov}
r_t = \rho_b \qquad \text{for } T_b < t \le T_b + \ell_b .
\end{equation}
The process is semi-Markov when $\rho_{b+1}$ depends on the past only through
$\rho_b$ (a transition kernel on $\mathcal{R}$) while $\ell_b$ is drawn from a
regime-specific duration distribution $G_{\rho_b}$ on $\{1,2,\ldots\}$.
\end{definition}

Hamilton's Markov-switching model \citep{hamilton1989new}---the canonical
formalization of regimes in econometrics---is the special case in which $r_t$
is a first-order Markov chain: block durations are then forced to be
geometric, $\Pr(\ell = k) = P_{rr}^{\,k-1}(1 - P_{rr})$ with mean
$1/(1 - P_{rr})$. The semi-Markov form frees the duration law $G_r$, which
matters because the experimental stream uses bounded, non-geometric blocks
(25--60 days common, 15--30 days crisis) and because real regime durations
are empirically far from geometric \citep{ang2012regime, guidolin2007asset}.
Nothing in Part~III depends on the duration law beyond persistence itself.

The concrete instantiation used throughout: $R = 4$ regimes
$\mathcal{R} = \{\textsf{trend}, \textsf{range}, \textsf{highvol},
\textsf{crisis}\}$; common regimes arrive in blocks of 25--60 days;
\textsf{crisis} arrives in short blocks of 15--30 days and only on every
\emph{fourth} cycle through the common regimes, so consecutive crisis
episodes are separated by roughly 250--450 days of dormancy; horizon
$T \approx 2370$ days ($\approx 9.5$ trading years); each day carries
features $X_t \in \R^{n \times d}$ ($n = 20$ assets, $d = 20$ features per
asset) and realized returns $y_t \in \R^{n}$.

\subsection{Why Block Persistence (Not I.I.D.\ Regime Draws) Is the Whole Game}

It is tempting to model regime-switching as $r_t \stackrel{iid}{\sim} \pi$
for some distribution $\pi$ on $\mathcal{R}$. That model destroys both of the
phenomena this project is about.

\begin{proposition}[Persistence creates local stationarity and dormancy]
\label{prop:dd2-persistence}
Let $r_t$ be a regime process over $\mathcal{R}$.
\begin{enumerate}
  \item \textbf{(I.i.d.\ draws have no usable present.)} If
        $r_t \stackrel{iid}{\sim} \pi$, the conditional law of $r_{t+1}$
        given the entire past equals $\pi$: knowing today's regime says
        nothing about tomorrow's, every day is a fresh mixture draw, and the
        optimal single-day strategy is mixture-optimal, not any
        regime-specialist.
  \item \textbf{(I.i.d.\ draws have no dormancy.)} Under i.i.d.\ draws,
        $\Pr(\text{regime } r \text{ dormant} > D \text{ days})
        = (1 - \pi(r))^{D}$, exponentially small in $D$. Under the block
        process, dormancy is the sum of intervening block durations and is
        typically hundreds of days when $r$ sits on a long cycle---as
        \textsf{crisis} does here ($250$--$450$ days by construction).
\end{enumerate}
\end{proposition}

\begin{proof}
(1) is immediate from independence. (2): with
$\pi(\textsf{crisis}) \approx 1/13$ (one short crisis block per $\approx 12$
common blocks), an i.i.d.\ stream would produce a 250-day crisis dormancy
with probability $(1 - 1/13)^{250} \approx e^{-19.2} \approx 5 \times
10^{-9}$---essentially never---whereas the block schedule produces one every
cycle. The block construction makes long dormancy the rule rather than a
large-deviation event.
\end{proof}

\begin{intuition}[title={Persistence is what makes specialization purchasable}]
Persistence does two jobs at once. \emph{Within} a block, the world is locally
stationary for 25--60 days, long enough for a specialist to be trained,
evaluated, and trusted---this is what makes a regime a \emph{niche} worth
owning. \emph{Between} blocks of the same regime, the world is silent about
that regime for weeks to years---this is what makes memory a problem. Under
i.i.d.\ draws neither holds: there is nothing to specialize in (job one fails)
and nothing is ever dormant (job two never arises). Every result in this
document lives in the gap between those two timescales.
\end{intuition}

\subsection{The Causal Labelers}
\label{sec:dd2-labelers}

All eleven arms of the experiment receive the regime label $\hat r_t$ from a
shared \textbf{causal labeler}: a deterministic function of information
available \emph{strictly before} the decision at day $t$ is made. Two labelers
are used, deliberately spanning a simple/structured spectrum, so that no
conclusion depends on one labeling artifact.

\paragraph{Labeler L1: rolling-percentile volatility band $\times$ trend sign.}
At day $t$, using returns $y_1, \ldots, y_{t-1}$ only:
\begin{enumerate}
  \item \textbf{Trailing volatility.} Realized volatility
        $v_t = \operatorname{sd}\bigl(\bar y_{t-w_v}, \ldots, \bar y_{t-1}\bigr)$
        over a trailing window of $w_v$ days, $\bar y_s$ the cross-sectional
        mean return on day $s$.
  \item \textbf{Rolling percentile thresholds.} The empirical
        $q^{\text{high}}$- and $q^{\text{crisis}}$-quantiles
        $\tau^{\text{high}}_t, \tau^{\text{crisis}}_t$ of the trailing
        history $\{v_s : s \le t-1\}$---quantiles of \emph{past}
        volatilities only, re-estimated daily from the filtration at $t-1$,
        never from the full sample.
  \item \textbf{Trend statistic.} A trailing trend signal $m_t$ (sign and
        magnitude of the mean of $\bar y_{t-w_m}, \ldots, \bar y_{t-1}$, or
        a fast-minus-slow moving-average gap), again from days $\le t-1$.
  \item \textbf{Label.}
        \[
        \hat r_t =
        \begin{cases}
          \textsf{crisis} & v_t \ge \tau^{\text{crisis}}_t, \\
          \textsf{highvol} & \tau^{\text{high}}_t \le v_t < \tau^{\text{crisis}}_t, \\
          \textsf{trend} & v_t < \tau^{\text{high}}_t \text{ and } |m_t| \ge \tau^{m}, \\
          \textsf{range} & v_t < \tau^{\text{high}}_t \text{ and } |m_t| < \tau^{m}.
        \end{cases}
        \]
\end{enumerate}

\paragraph{Labeler L2: forward-only filtered two-state volatility model
$\times$ trend sign.}
L2 replaces the percentile band with a probabilistic state model:
\begin{enumerate}
  \item Maintain a two-state Markov-switching volatility model (low/high-vol
        states with state-dependent return variance), parameters estimated
        on an expanding window through day $t-1$ and refit periodically from
        past data only.
  \item Run the \textbf{Hamilton filter} \citep{hamilton1989new}
        \emph{forward only} for the one-step-ahead predictive probability
        $p_t = \Pr(s_t = \text{high-vol} \mid \bar y_1, \ldots, \bar y_{t-1})$.
        The Kim \emph{smoothed} probability, which conditions on the future,
        is never computed.
  \item Label \textsf{crisis} when $p_t$ exceeds a high threshold and
        trailing volatility is extreme, \textsf{highvol} when $p_t$ exceeds
        a moderate threshold, otherwise \textsf{trend}/\textsf{range} by the
        same trend statistic $m_t$ as in L1.
\end{enumerate}

\paragraph{Leakage analysis.}
Write $\mathcal{F}_{t-1} = \sigma(X_1, y_1, \ldots, X_{t-1}, y_{t-1}, X_t)$
for the information legitimately available when the day-$t$ decision is
placed (today's features are observable; today's returns are not). Both
labelers satisfy the \textbf{causality contract}
\begin{equation}
\label{eq:dd2-causal}
\hat r_t \in \mathcal{F}_{t-1}
\qquad \text{(in fact } \hat r_t \in \sigma(y_1, \ldots, y_{t-1})\text{)},
\end{equation}
because every ingredient---trailing volatility, rolling quantiles, trend
statistics, filter recursions, parameter refits---is a function of returns
through $t-1$. Three subtler channels are also closed:
\emph{(i) contemporaneous leak:} the label for day $t$ never uses $y_t$
itself (a one-day embargo between the last return used and the day labeled);
\emph{(ii) refit leak:} parameters and thresholds are re-estimated only on
data through $t-1$, so no full-sample estimate back-propagates into early
labels;
\emph{(iii) smoothing leak:} L2 uses filtered, never smoothed,
probabilities---the smoother is two-sided, and a single smoothed label can
encode information from arbitrarily far in the future.

\begin{warning}[title={Label leakage: the classic regime-paper sin}]
A large fraction of the regime-switching literature evaluates strategies
against regime labels that \emph{could not have been known at the time}. The
recurring offenders:
\begin{itemize}
  \item \textbf{Smoothed state probabilities:} using Kim-smoothed
        $\Pr(s_t \mid y_{1:T})$ as ``the regime at $t$''---the smoother
        conditions on the future and labels turning points in advance.
  \item \textbf{Full-sample thresholds:} defining \textsf{highvol} as the
        top volatility quintile \emph{of the whole sample}---the 1990
        boundary then depends on 2008.
  \item \textbf{Peak/trough dating:} NBER-style bull/bear datings require
        future data to \emph{confirm} a peak; using the dated label at the
        peak itself is hindsight.
  \item \textbf{Full-sample parameters with filtered probabilities:} the
        leak hides in $\hat\theta$, not in the recursion.
\end{itemize}
Each inflates the apparent value of regime-conditioning for the same reason
published return predictors decay out of sample \citep{mclean2016does}: the
evaluation quietly conditions on information the strategy never had.
Everything downstream---specialist training, competition windows,
pinning---consumes only the causal label
$\hat r_t \in \mathcal{F}_{t-1}$ of \eqref{eq:dd2-causal}; ``reactivation''
means the \emph{causal labeler} says so, with whatever lag and
misclassification that entails.
\end{warning}

\begin{remark}[The label is an input, not a contribution]
\nm{} takes $\hat r_t$ as given and identical across all arms: the
experiment isolates \emph{what to do with} a regime signal (allocation,
training, memory), not how to detect regimes---a separate, well-studied
problem \citep{nystrup2018dynamic, nystrup2020greedy}. A shared causal label
removes regime-detection skill as a confound.
\end{remark}

\section{Episodes as Environments}
\label{sec:dd2-episodes}

\subsection{Episodes}

\begin{definition}[Episode]
\label{def:dd2-episode}
For regime $r \in \mathcal{R}$, the \textbf{$e$-th episode of $r$}, written
$(r, e)$, is the $e$-th maximal block of consecutive days on which
$\hat r_t = r$, i.e.\ the $e$-th element of the subsequence of blocks
$\{b : \rho_b = r\}$ in Definition~\ref{def:dd2-regime-process}. We write
$\mathcal{T}_{r,e} \subset \{1, \ldots, T\}$ for its day set and
$|\mathcal{T}_{r,e}| = \ell_{r,e}$ for its length.
\end{definition}

The central modeling move---inherited from invariant-prediction thinking and
from the decision-focused learning literature
\citep{elmachtoub2022smart, yan2026invpnco}---is to treat each episode as a
distinct \emph{environment}: the regime fixes what is stable, the episode
index fixes what is incidental.

\begin{definition}[Hyper-distribution assumption]
\label{def:dd2-hyper}
Each regime $r$ carries a \textbf{hyper-distribution} $Q_r$ over environment
parameters. Episode $(r,e)$ draws its parameter
$\gamma_{r,e} \sim Q_r$ independently across $e$, and the days within the
episode are then i.i.d.\ given $(r, \gamma_{r,e})$:
\[
\gamma_{r,e} \stackrel{iid}{\sim} Q_r, \qquad
(X_t, y_t) \mid (r, \gamma_{r,e}) \stackrel{iid}{\sim} P_{r, \gamma_{r,e}}
\quad \text{for } t \in \mathcal{T}_{r,e}.
\]
The \textbf{stable component} of regime $r$ is whatever is shared by
$P_{r,\gamma}$ across all $\gamma$ in the support of $Q_r$; the
\textbf{episode-specific component} is the rest.
\end{definition}

\begin{remark}[Two honest idealizations]
\label{rem:dd2-idealizations}
Definition~\ref{def:dd2-hyper} idealizes reality in two ways.
\emph{(i) Macro memory.} Real episodes of the same regime are not independent
draws: the 2008 and 2011 crises share causal machinery and participants who
remember the previous one. Positive dependence across $e$ means the effective
number of independent episodes is \emph{smaller} than the count---scarcity
(Section~\ref{sec:dd2-scarcity}) is, if anything, understated.
\emph{(ii) Secular drift.} $Q_r$ is treated as fixed over the horizon, but
markets drift on decadal scales: predictors decay after publication
\citep{mclean2016does}, and the ``stable'' component of a regime is stable
only relative to the horizon. The synthetic generator below satisfies both
idealizations exactly, so failures of an arm on synthetic data cannot be
blamed on them; transfer to long real panels inherits both caveats.
\end{remark}

\subsection{Episode Scarcity and Its Two Mitigations}
\label{sec:dd2-scarcity}

The statistical currency of Definition~\ref{def:dd2-hyper} is the number of
\emph{episodes} $E_r$, not the number of days: any quantity that varies at the
episode level (an episode-mean loss, an episode-specific coefficient) is
estimated with variance $\propto 1/E_r$, no matter how many days each episode
contains. For common regimes $E_r$ is comfortable (dozens of 25--60-day
blocks over $T \approx 2370$ days). For \textsf{crisis} it is brutal: in
roughly 35 years of equity data one finds on the order of \textbf{5--7}
crisis episodes, and fewer in shorter panels; in the synthetic stream, the
every-fourth-cycle schedule produces a comparable handful over
$T \approx 2370$ days. Two standard mitigations trade scarcity for
dependence, and it is important to state exactly what dependence each one
buys.

\begin{definition}[Sub-episode blocks]
\label{def:dd2-subepisode}
Fix a block length $B$. The \textbf{sub-episode blocks} of episode $(r,e)$
are the contiguous segments obtained by partitioning $\mathcal{T}_{r,e}$ into
pieces of length $B$ (discarding or merging a short remainder). Each block is
treated as a pseudo-episode in episode-level statistics.
\end{definition}

Sub-episode blocks multiply the apparent environment count by
$\approx \ell_{r,e}/B$, but all blocks of one episode \emph{share the same
draw} $\gamma_{r,e}$, so block mean losses are positively correlated through
$\gamma$, and cross-block variance \emph{underestimates} cross-episode
variance: the blocks inform about within-episode noise, not about $Q_r$.
They are a variance mitigation for the stable component only; they cannot
manufacture information about episode heterogeneity the panel does not
contain.

\begin{definition}[Cross-asset episodes]
\label{def:dd2-crossasset}
When the same regime episode is observed simultaneously on $A$ assets or
markets (here $n = 20$ assets share each episode), the per-asset day records
$\{(x_{t,i}, y_{t,i}) : t \in \mathcal{T}_{r,e}\}$ for $i = 1, \ldots, A$ may
be treated as $A$ parallel realizations of the episode environment.
\end{definition}

Cross-asset episodes multiply the day count per episode by $A$, but
introduce \emph{cross-sectional dependence}: assets within a day share
common shocks (market factors, and the episode draw $\gamma_{r,e}$ itself,
common across assets in the generator below), so the effective sample size
per day lies between $1$ and $A$. Again: more data about the within-episode
law, no new draws from $Q_r$.

\begin{keypoint}
Both mitigations enlarge the denominator of \emph{within-episode} estimates
and leave the number of independent draws from $Q_r$ untouched. Whatever
must be learned \emph{about episode heterogeneity}---in particular the
variance-across-episodes term the INV objective penalizes---must come from
$E_r$ genuine episodes, and for \textsf{crisis} $E_r$ is single-digit. This
one fact sizes the episode library ($\le 6$ retained episodes suffices,
because there are never many more) and explains the crisis-pinning wrinkle
of Section~\ref{sec:dd2-competition}.
\end{keypoint}

\subsection{The Synthetic Generator: Heterogeneity Made Explicit}
\label{sec:dd2-generator}

The synthetic stream instantiates Definition~\ref{def:dd2-hyper} exactly. For
day $t$ in episode $(r,e)$ and asset $i$ with feature vector
$x_{t,i} \in \R^{20}$:
\begin{equation}
\label{eq:dd2-generator}
y_{t,i} \;=\;
\underbrace{\theta_r \cdot (x_{t,i,1},\, x_{t,i,2})}_{\text{stable, regime-level}}
\;+\;
\underbrace{\gamma_{r,e}\, x_{t,i,\mathrm{sp}}}_{\text{episode-specific}}
\;+\; b_r \;+\; \varepsilon_{t,i},
\end{equation}
with the following calibration:
\begin{itemize}
  \item \textbf{Stable signal.} $\theta_r \in \R^2$, $\|\theta_r\| = 0.012$,
        directions well-separated across regimes---so a head trained on one
        regime is wrong (not merely stale) on another.
  \item \textbf{Episode-specific signal.}
        $\gamma_{r,e} \sim \mathcal{N}\bigl(0, (0.012 \cdot \mathrm{het})^2\bigr)$,
        \emph{redrawn independently each episode}, loading on a designated
        spurious feature $x_{\mathrm{sp}}$; $\mathrm{het} \ge 0$ is the
        heterogeneity knob.
  \item \textbf{Distractors.} 16 of the 20 features carry no signal in any
        regime: overfitting bait at $n$--$d$ parity ($n = d = 20$).
  \item \textbf{Noise.} $\varepsilon_{t,i} \sim \mathcal{N}(0, \sigma_r^2)$,
        $\sigma_r = 0.02 \cdot (1,\, 0.75,\, 2,\, 3)$ for
        (\textsf{trend}, \textsf{range}, \textsf{highvol}, \textsf{crisis}):
        the rare regime is also the noisiest, by a factor of three.
\end{itemize}

\begin{example}[What $\mathrm{het} = 0$ versus $\mathrm{het} = 2$ means]
\label{ex:dd2-het}
The standard deviation of the episode-specific coefficient is
$\operatorname{sd}(\gamma_{r,e}) = 0.012 \cdot \mathrm{het}$, against a stable
coefficient magnitude of $\|\theta_r\| = 0.012$. Concretely:
\begin{itemize}
  \item $\mathbf{het = 0}$: $\gamma_{r,e} \equiv 0$. Every episode of regime
        $r$ has the \emph{same} law; the spurious feature is a 17th
        distractor; $Q_r$ is a point mass; plain ERM on pooled regime data
        is right, and a variance-across-episodes penalty has nothing to
        penalize.
  \item $\mathbf{het = 1}$: $\operatorname{sd}(\gamma) = 0.012 = \|\theta_r\|$.
        The spurious coefficient is at least as large as the stable one
        whenever $|\gamma_{r,e}| \ge 0.012$, probability
        $\Pr(|Z| \ge 1) \approx 0.317$ for $Z \sim \mathcal{N}(0,1)$: in
        roughly one episode in three, the most profitable in-episode feature
        is the one that betrays you next episode.
  \item $\mathbf{het = 2}$: $\operatorname{sd}(\gamma) = 0.024 = 2\|\theta_r\|$.
        Now $|\gamma_{r,e}| \ge \|\theta_r\|$ with probability
        $\Pr(|Z| \ge 0.5) \approx 0.617$: in the \emph{majority} of episodes
        the episode-specific signal dominates inside the episode. An
        in-episode ERM fit loads heavily on $x_{\mathrm{sp}}$---correctly,
        for that episode---then carries a coefficient of typical size
        $0.024$ with a freshly randomized sign into the next episode.
\end{itemize}
The probe window (the first 15 days after a reactivation,
Definition~\ref{def:dd2-probe}) is precisely where this bites: the learner
meets a new draw of $\gamma_{r,e}$ before any in-episode data has
accumulated, so whatever $\gamma$-chasing is baked into its head is realized
as pure decision regret.
\end{example}

\begin{intuition}[title={The generator is a forgetting--overfitting double trap}]
Equation~\eqref{eq:dd2-generator} arms two independent failure modes. The
\emph{regime} structure (separated $\theta_r$, long dormancy) punishes arms
that forget; the \emph{episode} structure ($\gamma_{r,e}$ redrawn) punishes
arms that remember too literally---memorize the last crisis episode including
its $\gamma$ and you import a coin-flip coefficient into the next one. In the
eleven arms of Section~\ref{sec:dd2-arms}, the allocation axis addresses the
first trap, the objective axis (ERM vs.\ INV) the second.
\end{intuition}

\section{Dormancy Combinatorics and Memory Models}
\label{sec:dd2-dormancy}

\subsection{Dormancy and the Probe Window}

\begin{definition}[Dormancy]
\label{def:dd2-dormancy}
The \textbf{dormancy} of regime $r$ before episode $e+1$ is the gap
\[
D_r(e) \;=\; \min \mathcal{T}_{r,e+1} \;-\; \max \mathcal{T}_{r,e} \;-\; 1,
\]
the number of days between the end of episode $(r,e)$ and the start of
$(r,e+1)$ during which other regimes are active. In the synthetic stream,
common regimes have dormancies on the order of one cycle (tens of days),
while \textsf{crisis} has $D \approx 250$--$450$ days.
\end{definition}

\begin{definition}[Probe window and forgetting deficit]
\label{def:dd2-probe}
A \textbf{reactivation} of $r$ is the first day of an episode $(r, e+1)$ with
dormancy $D_r(e) \ge 90$ days. The \textbf{probe window} is the first
$15$ days of that episode. For an arm $A$, let
$\bar\rho_{\mathrm{probe}}(A)$ denote its mean per-day decision regret
(Definition~\ref{def:dd2-regret}) over probe-window days. The
\textbf{forgetting deficit} of arm $A$ at dormancy $D$ is
\begin{equation}
\label{eq:dd2-phi}
\Phi(D; A) \;=\;
\E\bigl[\bar\rho_{\mathrm{probe}}(A) \,\big|\, \text{dormancy } D\bigr]
\;-\;
\E\bigl[\bar\rho_{\mathrm{probe}}(A^{\circ})\bigr],
\end{equation}
where $A^{\circ}$ is the \emph{retention counterfactual}: the same arm
serving the probe window with the head (and episode library) it held at the
close of the previous episode of $r$, untouched by any intervening foreign
training. $\Phi \approx 0$ means dormancy cost the arm nothing it had;
$\Phi$ large means the dormancy machinery below did its damage.
\end{definition}

The probe window is 15 days: long enough to average regime noise
($\sigma_{\textsf{crisis}} = 0.06$ per day) into a stable regret estimate,
short enough that within-episode relearning cannot mask a cold start---it
measures what you walked in with. All headline numbers in this document
($0.0303$--$0.0320$) are means of $\bar\rho_{\mathrm{probe}}$ across
reactivations and 20 seeds.

\subsection{The Foreign-Update Process}

During $r$'s dormancy, the stream keeps producing days from the $W_a$
\emph{other} active regimes, and the holder of $r$'s head may be made to
\emph{train on a foreign regime}. We abstract this as the
\textbf{foreign-update process}: each dormancy day is a foreign-update day
with probability $p$, where $p$ encodes the arm: $p = 1$ for a monolith that
trains daily on whatever is active; $p \approx$ (win rate) for an unpinned
\nm{} specialist; $p = 0$ for a pinned owner whose regime is dormant. The
damage per foreign update is what the two memory models specify.

\begin{definition}[Hard and soft memory models]
\label{def:dd2-memory-models}
Each specialist has capacity $K < R$: at most $K$ regime heads (linear SPO+
heads, Section~\ref{sec:dd2-specialist}), each with an episode buffer of up to
$6$ episodes $\times$ $40$ days.
\begin{itemize}
  \item \textbf{HARD (LRU eviction).} Heads occupy discrete slots ordered by
        recency of training. Training on a held regime refreshes its recency.
        Training on an \emph{unheld} regime at full capacity instantiates a
        new head and \textbf{evicts} the least-recently-used head
        \emph{together with its episode buffer}---competence and memory are
        destroyed jointly and completely.
  \item \textbf{SOFT (interference decay).} No head is ever evicted. Instead,
        every foreign update multiplies \emph{all other held heads} by
        $(1 - \rho_{\mathrm{int}})$ with $\rho_{\mathrm{int}} = 0.06$:
        after $N$ foreign updates a dormant head retains strength
        $(1 - \rho_{\mathrm{int}})^{N}$ of its trained weights.
\end{itemize}
\end{definition}

The hard model is the cache-eviction picture; the soft model is the
gradient-interference picture of catastrophic forgetting in neural networks
\citep{french1999catastrophic, kirkpatrick2017overcoming, parisi2019continual},
where training on task $B$ perturbs the very weights encoding task $A$.
Progressive networks \citep{rusu2016progressive} avoid interference by
freezing columns---a structural cousin of what pinning will achieve in
Section~\ref{sec:dd2-competition}.

\subsection{Hard Model: The Coupon-Collector Clock}

Under HARD, a dormant head dies when enough \emph{distinct} foreign regimes
have been learned to push it out of the LRU order. The clean question: a
learner holds the dormant head and has $K$ slots; how long until it has
trained on $K$ distinct foreign regimes (each claiming a slot, the $K$-th
overflowing capacity and evicting the dormant head, which---having received no
training all dormancy---is always the least recently used)?

\begin{proposition}[Expected time to eviction]
\label{prop:dd2-coupon}
Suppose each dormancy day is independently a foreign-update day with
probability $p$, and on an update day the active regime is uniform on the
$W_a$ foreign regimes, independently across update days. Let $T_K$ be the
number of days until $K$ distinct foreign regimes have been trained on
($K \le W_a$). Then
\begin{equation}
\label{eq:dd2-coupon}
\E[T_K] \;=\; \frac{1}{p} \sum_{j=0}^{K-1} \frac{W_a}{W_a - j}.
\end{equation}
\end{proposition}

\begin{proof}
Stage $j$ begins when exactly $j$ distinct foreign regimes have been
collected, $0 \le j \le K-1$. On each day of stage $j$, the event ``collect a
new regime today'' requires an update day (probability $p$) on which the
active regime is one of the $W_a - j$ uncollected ones (probability
$(W_a - j)/W_a$, by uniformity and independence). The stage-$j$ waiting time
is therefore geometric with success probability $p (W_a - j)/W_a$ and
expectation $W_a / \bigl(p (W_a - j)\bigr)$. The stages are traversed in
order, so by linearity of expectation
$\E[T_K] = \sum_{j=0}^{K-1} W_a / \bigl(p(W_a - j)\bigr)$, which is
\eqref{eq:dd2-coupon}.
\end{proof}

\begin{example}[Worked numbers: $W_a = 3$, $K = 2$, $p = \tfrac12$]
\label{ex:dd2-coupon}
A learner holds the dormant \textsf{crisis} head with capacity $K = 2$; the
$W_a = 3$ common regimes are active; it trains on the active regime every
other day on average. By \eqref{eq:dd2-coupon},
\[
\E[T_2]
= \frac{1}{0.5}\left(\frac{3}{3} + \frac{3}{2}\right)
= 2 \times 2.5
= 5 \text{ days}:
\]
$1/p = 2$ expected days for any foreign update to claim the spare slot, plus
$3/(0.5 \times 2) = 3$ for an update from one of the remaining two regimes.
Five days---against a crisis dormancy of $250$--$450$ days: the dormant
head's expected lifetime is roughly $1\%$ of the dormancy it must survive.
\end{example}

How certain is eviction by day $D$? For the worked example we can compute the
survival probability exactly.

\begin{proposition}[Survival probability, exact and bounded]
\label{prop:dd2-survival}
In the setting of Proposition~\ref{prop:dd2-coupon} with $W_a = 3$, $K = 2$:
\begin{equation}
\label{eq:dd2-survival}
\Pr(\text{head survives } D \text{ days})
= 3\Bigl(1 - \tfrac{2p}{3}\Bigr)^{D} - 2(1 - p)^{D}.
\end{equation}
In general ($K \le W_a$), a union bound over which $K-1$ regimes were the
only ones updated gives the binomial-tail bound
\[
\Pr(\text{survive } D)
\;\le\; \binom{W_a}{K-1}
\Bigl(1 - p\,\tfrac{W_a - K + 1}{W_a}\Bigr)^{D}.
\]
\end{proposition}

\begin{proof}
Survival through day $D$ means at most $K - 1 = 1$ distinct regime was
trained on. Condition on $N_D \sim \mathrm{Bin}(D, p)$ update days: given
$N_D = m \ge 1$, all $m$ uniform draws land on a single regime with
probability $3 \cdot (1/3)^m$; given $m = 0$, survival is certain. Hence
\[
\Pr(\text{survive})
= (1-p)^D + \sum_{m \ge 1} \binom{D}{m} p^m (1-p)^{D-m} \cdot 3 \cdot 3^{-m}
= (1-p)^D + 3\Bigl[\bigl(1 - p + \tfrac{p}{3}\bigr)^{D} - (1-p)^D\Bigr],
\]
which is \eqref{eq:dd2-survival}. General bound: survival requires the
updated regimes to lie in some $(K-1)$-subset $S$; for fixed $S$, each day
avoids ``update outside $S$'' with probability $1 - p(W_a - K + 1)/W_a$;
union over the $\binom{W_a}{K-1}$ subsets (this matches the leading term of
\eqref{eq:dd2-survival} at $W_a = 3$, $K = 2$).
\end{proof}

\begin{example}[Eviction probabilities by dormancy length]
With $p = 0.5$, evaluating \eqref{eq:dd2-survival} (note
$1 - 2p/3 = 2/3$, $1 - p = 1/2$):

\begin{center}
\begin{tabular}{rccc}
\toprule
$D$ (days) & $3(2/3)^D$ & $2(1/2)^D$ & $\Pr(\text{evicted by } D)$ \\
\midrule
5  & $0.39506$ & $0.06250$  & $1 - 0.33256 = 0.667$ \\
10 & $0.05202$ & $0.00195$  & $1 - 0.05007 = 0.950$ \\
20 & $0.00090$ & $1.9\times 10^{-6}$ & $1 - 0.00090 = 0.9991$ \\
40 & $2.7\times 10^{-7}$ & $1.8\times 10^{-12}$ & $1 - 2.7\times 10^{-7}$ \\
\bottomrule
\end{tabular}
\end{center}

By $D = 20$ days---half of one common-regime block---eviction is a $99.9\%$
event; at a crisis dormancy of $D = 250$ survival is
$3(2/3)^{250} \approx 3 \times 10^{-44}$. Under the hard model an
unprotected crisis head \emph{does not survive dormancy}, full stop. The
i.i.d.-day idealization is conservative here: the real stream arrives in
persistent blocks, each delivering many updates from one regime, so
distinct-regime collection is at least as fast at the block scale.
\end{example}

\begin{example}[LRU mechanics, step by step]
\label{ex:dd2-lru}
A capacity-$K{=}2$ learner ends a crisis episode holding heads for
$\{\textsf{crisis}, \textsf{trend}\}$, with \textsf{crisis} most recent.
The dormancy stream then activates \textsf{trend}, \textsf{range},
\textsf{highvol}, \ldots\ and the learner trains on whatever is active. Slot
contents (ordered LRU $\to$ MRU; each head carries its episode buffer):

\begin{center}
\begin{tabular}{clll}
\toprule
Event & Active regime & Action & Slots after (LRU $\to$ MRU) \\
\midrule
1 & \textsf{trend}   & hit: refresh recency & $[\textsf{crisis}, \textsf{trend}]$ \\
2 & \textsf{range}   & miss at full capacity: \textbf{evict \textsf{crisis}} + buffer & $[\textsf{trend}, \textsf{range}]$ \\
3 & \textsf{crisis} reactivates & miss: \textbf{cold start}, evict \textsf{trend} & $[\textsf{range}, \textsf{crisis}^{\,\text{new}}]$ \\
\bottomrule
\end{tabular}
\end{center}

At event 2---the \emph{first} foreign block from an unheld regime---the
crisis head and its entire episode library are gone. At event 3 the probe
window is served by a freshly initialized head: full cold-start regret, then
a 15--30-day scramble to relearn what was once known, and the cycle repeats.
\end{example}

\subsection{Soft Model: Geometric Erosion and the Half-Life}

Under SOFT nothing is evicted; instead each foreign update multiplies every
other held head by $(1 - \rho_{\mathrm{int}})$ with
$\rho_{\mathrm{int}} = 0.06$. After $N_D$ foreign updates during a dormancy,
the dormant head retains
\begin{equation}
\label{eq:dd2-soft-retention}
\text{retained strength} \;=\; (1 - \rho_{\mathrm{int}})^{N_D} \;=\; 0.94^{N_D},
\end{equation}
of its trained weights (the head decays toward initialization; it is never
refreshed because its regime emits no training days). The natural summary is
the \textbf{half-life in foreign updates}: the $N$ solving
$(1-\rho)^{N} = \tfrac12$,
\begin{equation}
\label{eq:dd2-halflife}
N_{1/2} \;=\; \frac{\log(1/2)}{\log(1 - \rho_{\mathrm{int}})}
\;=\; \frac{-0.693147}{\log(0.94)}
\;=\; \frac{0.693147}{0.061875}
\;\approx\; 11.2 \text{ foreign updates},
\end{equation}
using $\log(0.94) = -0.061875$. Eleven foreign updates---about two competition
windows' worth of foreign training, or a week and a half of daily monolith
training---halve the head. The full erosion schedule:

\begin{center}
\begin{tabular}{rcccccc}
\toprule
Foreign updates $N$ & 11 & 12 & 25 & 50 & 100 & 250 \\
\midrule
Retained $0.94^{N}$ & $0.506$ & $0.476$ & $0.213$ & $0.045$ & $0.0021$ & $1.9 \times 10^{-7}$ \\
\bottomrule
\end{tabular}
\end{center}

For a monolith that trains every day ($N_D = D$), a crisis dormancy of
$250$--$450$ days leaves $0.94^{250} \approx 2 \times 10^{-7}$ of the head:
a cold start without a discrete eviction event. A learner suffering one
foreign update per 20-day window over the same dormancy takes
$N_D \approx 12$--$22$ updates, retention $0.48$--$0.25$: degraded, not
destroyed. The soft model grades the damage by the \emph{training cadence}
of the allocation rule---exactly the quantity the eleven arms vary.

\begin{keypoint}[title={Two models, one phenomenon, bounded from both ends}]
HARD and SOFT bracket the same constraint---finite capacity under foreign
training---from opposite extremes. HARD is the worst case per event: one
overflow update destroys a head completely, at a stopping time whose law
Propositions~\ref{prop:dd2-coupon}--\ref{prop:dd2-survival} pin down. SOFT is
the gentlest plausible case: no single update matters, yet the geometric
decay \eqref{eq:dd2-soft-retention} reaches the same endpoint at the
dormancies in play ($0.94^{250} \approx 10^{-7}$). A conclusion that holds
under \emph{both} is about bounded capacity itself, not a bookkeeping rule;
and a mechanism that protects a head must both prevent overflow updates
(HARD) and drive the foreign-update count to zero (SOFT). Pinning
(Section~\ref{sec:dd2-competition}) does both at once: after a covering
one-per-regime assignment, each specialist's foreign-update probability is
structurally $p = 0$.
\end{keypoint}

\begin{remark}[Why $\Phi$ is defined against the retention counterfactual]
The forgetting deficit \eqref{eq:dd2-phi} subtracts the regret of the
arm's \emph{own frozen past self}, not of an oracle, isolating dormancy
damage from competence: a mediocre head retained perfectly has $\Phi = 0$
(a skill problem, not a memory problem); an excellent head lost to eviction
has large $\Phi$ even at moderate probe regret. The eleven-arm comparison of
Section~\ref{sec:dd2-arms} needs both coordinates---probe regret for the
bottom line, $\Phi$ for the diagnosis.
\end{remark}

% ============================================================================
% PART III: THE RISP MECHANISM
% ============================================================================
\part{The \nm{} Mechanism}
\label{sec:dd2-part-mechanism}

Part~II established the problem: regimes recur after dormancies that no
capacity-limited, reward-chasing learner survives, while episode
heterogeneity punishes whoever survives by memorizing too literally. This
part builds the mechanism that answers it: specialists
(Section~\ref{sec:dd2-specialist}), windowed competition with multiplicative
affinity updates and an absorbing pin (Section~\ref{sec:dd2-competition}),
and the placement of \nm{} among ten alternatives in a two-axis design space
(Section~\ref{sec:dd2-arms}).

\section{The Specialist}
\label{sec:dd2-specialist}

\subsection{State: Identity, Competence, Memory}

\begin{definition}[Specialist state]
\label{def:dd2-specialist}
The population contains $N = 4$ specialists. Specialist $i$ at day $t$
carries three coupled state components:
\begin{enumerate}
  \item \textbf{Identity: niche affinity} $\alpha_i \in \Delta^{R}$,
        initialized uniform ($\alpha_{i,r} = 1/R = 0.25$): the specialist's
        \emph{claim} on each regime, updated only by winning competition
        windows (Section~\ref{sec:dd2-competition}); it controls the
        incumbency bonus and, ultimately, pinning.
  \item \textbf{Competence: regime heads} $\{w_{i,r} : r \in H_i(t)\}$, one
        linear SPO+ head per \emph{held} regime, $|H_i(t)| \le K < R$,
        mapping features to predicted returns
        $\hat y_{t,i'} = w_{i,r} \cdot x_{t,i'}$ and thence to a decision
        via the optimizer (Definition~\ref{def:dd2-regret}).
  \item \textbf{Memory: episode library}
        $\{\mathcal{B}_{i,r} : r \in H_i(t)\}$, at most $6$ episodes
        $\times$ $40$ days per held regime of $(X_t, y_t)$ records keyed by
        episode---created with the head and \textbf{evicted with the head}
        (hard model); the raw material for both training variants below.
\end{enumerate}
\end{definition}

The three components are deliberately welded together: affinity determines
where the specialist wins; winning triggers training, which builds the head
and writes the buffer; head and buffer determine future window performance,
which feeds back into affinity. Identity, competence, and memory rise
together and (under eviction) fall together.

\begin{definition}[Decision regret]
\label{def:dd2-regret}
Fix a compact feasible decision set $\mathcal{A} \subset \R^{n}$ (a polytope
of admissible portfolios over the $n = 20$ assets), and let
$a^{*}(v) \in \argmax_{a \in \mathcal{A}} v^{\top} a$ with optimal value
$z^{*}(v) = \max_{a \in \mathcal{A}} v^{\top} a$. The \textbf{per-day
decision regret} of a prediction $\hat y_t$ against realized returns $y_t$ is
\begin{equation}
\label{eq:dd2-regret}
\rho(\hat y_t,\, y_t) \;=\; z^{*}(y_t) \;-\; y_t^{\top} a^{*}(\hat y_t)
\;\ge\; 0:
\end{equation}
the return foregone by acting on the prediction instead of on the truth.
Regret depends on $\hat y_t$ only through the \emph{decision} it induces---two
predictions with different numbers but the same induced portfolio have
identical regret. This is the decision-focused viewpoint of
\citet{elmachtoub2022smart}: the loss is downstream decision quality, not
prediction error.
\end{definition}

\subsection{One Day in the Life}

Each day $t$, the causal labeler emits $\hat r_t$ (visible to everyone), and
the system follows one of two paths.

\textbf{Pinned case} ($\hat r_t$ pinned to owner $o(\hat r_t)$,
Definition~\ref{def:dd2-pin}):
\begin{enumerate}
  \item The owner alone serves the live decision
        $a_t = a^{*}\bigl(\hat y^{(o)}_t\bigr)$ using $w_{o, \hat r_t}$.
  \item The day's record $(X_t, y_t)$ is appended to the owner's buffer
        $\mathcal{B}_{o, \hat r_t}$ (current episode, capped at $40$ days;
        oldest of the $\le 6$ episodes rotated out when a new one opens).
  \item On its training cadence, the owner takes a training step
        (Section~\ref{sec:dd2-training}) on $\mathcal{B}_{o, \hat r_t}$.
        \emph{No other specialist decides, trains, or updates anything.}
\end{enumerate}

\textbf{Unpinned case}:
\begin{enumerate}
  \item \textbf{Shadow decisions.} Every specialist $i$ forms
        $\hat y^{(i)}_t$ (with $w_{i, \hat r_t}$ if held; with the
        zero-initialized default head, hence a default decision, if not) and
        its shadow decision $a^{(i)}_t = a^{*}(\hat y^{(i)}_t)$. The live
        decision is served by the current best claimant (highest
        $\alpha_{\cdot, \hat r_t}$).
  \item \textbf{Regret accounting.} Each specialist's realized regret
        $\rho^{(i)}_t = \rho(\hat y^{(i)}_t, y_t)$ is appended to the open
        competition window for regime $\hat r_t$. Windows are regime-gated:
        only days with $\hat r_t = r$ advance $r$'s window, so a window
        spans regime interruptions without contamination.
  \item \textbf{Window close.} When $r$'s window reaches $W_c = 20$
        in-regime days, a winner is declared
        (Section~\ref{sec:dd2-competition}), the winner's affinity takes one
        exponentiated-gradient step, and \emph{the winner alone} trains on
        the window's data (appended to its buffer for $r$; if it holds no
        head for $r$, one is instantiated now---the moment HARD eviction or
        SOFT interference strikes its \emph{other} heads).
  \item \textbf{Pin check.} If $\alpha_{i^{*}, r} \ge 0.95$, regime $r$ is
        pinned to $i^{*}$ permanently.
\end{enumerate}

\begin{intuition}[title={Why losers do not train}]
The most consequential line in the loop is the quiet one: only the window
winner trains. Training is simultaneously the reward (the winner improves at
$r$) and the cost-bearer (training on $r$ is a foreign update against every
other head the trainer holds). A specialist that keeps losing on $r$ never
pays $r$'s interference cost---its own niche's head sits at $p = 0$ in the
foreign-update process of Section~\ref{sec:dd2-dormancy}. Competition is not
just selecting who improves; it is rationing who gets \emph{damaged}.
\end{intuition}

\subsection{The SPO+ Training Step}
\label{sec:dd2-training}

Training must improve \emph{decisions}, so the loss is a decision-aware
surrogate: regret \eqref{eq:dd2-regret} itself is piecewise constant in
$\hat y$ (it changes only when the induced $\argmax$ jumps), hence useless
for gradients. The SPO+ surrogate of \citet{elmachtoub2022smart}, written
here in the return-maximization convention, is
\begin{equation}
\label{eq:dd2-spoloss}
\ell_{\mathrm{SPO+}}(\hat y, y)
\;=\;
\max_{a \in \mathcal{A}} \bigl\{ (2\hat y - y)^{\top} a \bigr\}
\;-\; 2 \hat y^{\top} a^{*}(y)
\;+\; z^{*}(y).
\end{equation}

\begin{lemma}[SPO+ is a convex upper bound on regret]
\label{lem:dd2-spoplus}
$\ell_{\mathrm{SPO+}}(\cdot, y)$ is convex in $\hat y$, satisfies
$\ell_{\mathrm{SPO+}}(y, y) = 0$, and
$\ell_{\mathrm{SPO+}}(\hat y, y) \ge \rho(\hat y, y)$ for all $\hat y$.
A subgradient in $\hat y$ is
\begin{equation}
\label{eq:dd2-spograd}
g(\hat y, y) \;=\; 2\bigl(a^{*}(2\hat y - y) - a^{*}(y)\bigr).
\end{equation}
\end{lemma}

\begin{proof}
Convexity: the first term of \eqref{eq:dd2-spoloss} is a maximum of affine
functions of $\hat y$, the rest linear or constant. At $\hat y = y$ the
three terms are $z^{*}(y) - 2z^{*}(y) + z^{*}(y) = 0$. Upper bound:
evaluating the max at the feasible point $a^{*}(\hat y)$,
\[
\ell_{\mathrm{SPO+}}(\hat y, y)
\;\ge\; 2\hat y^{\top}\bigl(a^{*}(\hat y) - a^{*}(y)\bigr)
+ \bigl(z^{*}(y) - y^{\top} a^{*}(\hat y)\bigr),
\]
where the first bracket is $\ge 0$ because $a^{*}(\hat y)$ maximizes
$\hat y^{\top} a$, and the second is exactly $\rho(\hat y, y)$.
Subgradient: Danskin's theorem on the max term gives
$2\,a^{*}(2\hat y - y)$; the linear term contributes $-2a^{*}(y)$.
\end{proof}

With the linear head $\hat y_{t,i} = w^{\top} x_{t,i}$, the chain rule turns
\eqref{eq:dd2-spograd} into a gradient in $w$:
$\nabla_{w}\, \ell_{\mathrm{SPO+}}(\hat y_t, y_t)
= \sum_{i=1}^{n} g_i(\hat y_t, y_t)\, x_{t,i}$.

\begin{example}[SPO+ arithmetic on two assets]
\label{ex:dd2-spoplus}
Let $\mathcal{A} = \{e_1, e_2\}$, truth $y = (0.01, -0.01)$, prediction
$\hat y = (-0.002, 0.004)$ (ranking inverted). Then $a^{*}(y) = e_1$,
$z^{*}(y) = 0.01$, $a^{*}(\hat y) = e_2$, so
$\rho = 0.01 - (-0.01) = 0.02$. For SPO+:
$2\hat y - y = (-0.014, 0.018)$, maximized by $e_2$ at $0.018$; the linear
term is $-2\hat y^{\top} e_1 = 0.004$; hence
$\ell_{\mathrm{SPO+}} = 0.018 + 0.004 + 0.01 = 0.032 \ge \rho = 0.02$.
The subgradient \eqref{eq:dd2-spograd} is $g = 2(e_2 - e_1) = (-2, 2)$:
descent pushes $\hat y_1$ up and $\hat y_2$ down---directly toward repairing
the ranking that caused the bad decision, not toward shrinking prediction
error per se.
\end{example}

\subsection{ERM versus INV over the Episode Library}
\label{sec:dd2-erm-inv}

A training step consumes the episode buffer
$\mathcal{B}_{i,r} = \{B_1, \ldots, B_E\}$, $E \le 6$ episodes. Write
\[
\bar\ell_e(w) \;=\; \frac{1}{|B_e|} \sum_{t \in B_e}
\ell_{\mathrm{SPO+}}\bigl(\hat y_t(w),\, y_t\bigr),
\qquad
\bar\ell(w) \;=\; \frac{1}{E} \sum_{e=1}^{E} \bar\ell_e(w)
\]
for the per-episode and pooled mean surrogate losses. The two variants:

\paragraph{ERM variant.} Sample one episode $e \sim \mathrm{Unif}\{1..E\}$
and take a subgradient step on $\bar\ell_e(w)$. In expectation this descends
the pooled objective $\bar\ell(w)$: all episode-days are equal, and the head
converges toward the buffer-pooled SPO+ minimizer---including whatever
episode-specific $\gamma$ structure dominates the buffer.

\paragraph{INV variant.} Descend the variance-penalized objective
\begin{equation}
\label{eq:dd2-invobj}
J(w) \;=\; \Var_e\bigl[\bar\ell_e(w)\bigr] \;+\; \beta\, \E_e\bigl[\bar\ell_e(w)\bigr]
\;=\; \frac{1}{E} \sum_{e=1}^{E} \bigl(\bar\ell_e(w) - \bar\ell(w)\bigr)^2
\;+\; \beta\, \bar\ell(w),
\end{equation}
which demands not just low mean loss but \emph{evenly} low loss across
episodes---the invariance principle transplanted into decision-focused
learning \citep{yan2026invpnco}.

\begin{proposition}[INV gradient]
\label{prop:dd2-invgrad}
The gradient of \eqref{eq:dd2-invobj} is
\begin{equation}
\label{eq:dd2-invgrad}
\nabla_w J(w) \;=\; \frac{1}{E} \sum_{e=1}^{E}
\Bigl[\, 2\bigl(\bar\ell_e(w) - \bar\ell(w)\bigr) + \beta \,\Bigr]\,
\nabla_w \bar\ell_e(w).
\end{equation}
\end{proposition}

\begin{proof}
Differentiate the variance term of \eqref{eq:dd2-invobj}:
\[
\nabla_w \Var_e
= \frac{1}{E} \sum_{e} 2\bigl(\bar\ell_e - \bar\ell\bigr)
\bigl(\nabla \bar\ell_e - \nabla \bar\ell\bigr)
= \frac{2}{E} \sum_{e} \bigl(\bar\ell_e - \bar\ell\bigr) \nabla \bar\ell_e
\;-\; \frac{2}{E} \nabla \bar\ell \underbrace{\sum_{e} \bigl(\bar\ell_e - \bar\ell\bigr)}_{=\,0}.
\]
The cross term vanishes identically because the centered residuals
$\bar\ell_e - \bar\ell$ sum to zero by definition of $\bar\ell$---this is the
one step worth seeing explicitly, as it is why the INV gradient is a clean
\emph{reweighting} of per-episode gradients rather than a new object. Adding
$\beta \nabla \bar\ell = (\beta/E) \sum_e \nabla \bar\ell_e$ and collecting
terms gives \eqref{eq:dd2-invgrad}.
\end{proof}

\begin{intuition}[title={Reading the INV weights}]
Equation~\eqref{eq:dd2-invgrad} gives episode $e$ the gradient weight
$2(\bar\ell_e - \bar\ell) + \beta$: an episode at the buffer average gets
weight $\beta$ (plain ERM); a worse-than-average episode gets weight
$> \beta$; and a better-than-average episode with
$\bar\ell_e < \bar\ell - \beta/2$ gets a \emph{negative} weight---gradient
\emph{ascent} on episodes served suspiciously well. That last case is the
anti-overfitting tooth. In the generator \eqref{eq:dd2-generator}, loading
on the stable direction $\theta_r$ lowers all $\bar\ell_e$ roughly equally
(no variance contribution), while loading on $x_{\mathrm{sp}}$ helps the
episodes whose $\gamma_{r,e}$ sign matches and hurts the rest---pumping the
variance term. INV strips $\gamma$-chasing weight while keeping
$\theta$-weight: the objective-level answer to the het trap of
Example~\ref{ex:dd2-het}, complementary to the allocation-level answer
(pinning). At $\mathrm{het} = 0$ the penalty has nothing real to penalize
and INV pays a small price---which is why the experiment carries both A5
(ERM) and A6 (INV) rather than declaring a winner a priori.
\end{intuition}

\section{Competition Dynamics}
\label{sec:dd2-competition}

\subsection{Windows, Winners, and the Exponentiated-Gradient Update}

\begin{definition}[Competition window and winner]
\label{def:dd2-window}
For an unpinned regime $r$, the competition window collects, for each
specialist $i$, its shadow decision regrets over $W_c = 20$ in-regime days.
At window close the \textbf{winner} is
\begin{equation}
\label{eq:dd2-winner}
i^{*} \;=\; \argmin_{i \in \{1..N\}} \;
\Bigl[\, \bar\rho^{(i)}_{\mathrm{win}}
\;-\; \lambda \bigl(\alpha_{i,r} - \tfrac{1}{R}\bigr) \Bigr],
\qquad \lambda = 0.3,
\end{equation}
where $\bar\rho^{(i)}_{\mathrm{win}}$ is $i$'s mean window regret: lowest
regret wins, with an \textbf{incumbency bonus} $\lambda(\alpha_{i,r} - 1/R)$
subtracted for affinity above uniform. The winner then applies one
\textbf{exponentiated-gradient (EG) step} to its own affinity vector,
\begin{equation}
\label{eq:dd2-eg}
\alpha_{i^{*}, r} \;\leftarrow\; \alpha_{i^{*}, r}\, e^{\eta},
\qquad \text{then renormalize } \alpha_{i^{*}} \text{ over } \Delta^{R},
\qquad \eta = 0.6,
\end{equation}
and trains on the window's data. Losers' affinity vectors are untouched.
\end{definition}

Two structural remarks. First, the bonus is centered at $1/R$, so an
uncommitted specialist gets zero bonus, and since
$\sum_r (\alpha_{i,r} - 1/R) = 0$, incumbency advantages across regimes are
\emph{zero-sum}: claiming $r$ harder structurally weakens the claim
everywhere else. Second, \eqref{eq:dd2-eg} belongs to the
multiplicative-update family of exponentiated gradient
\citep{kivinen1997exponentiated}, Hedge \citep{freund1997decision}, and
multiplicative weights \citep{arora2012multiplicative}---inverted: each
\emph{specialist} runs the update over \emph{regimes} (what shall I be?),
rather than a master algorithm weighting experts (whom shall I trust?).

Because only the $r$-coordinate is multiplied and the rest are renormalized
proportionally, the update has a two-point reduction: writing
$\alpha = \alpha_{i^{*}, r}$ before the win,
\begin{equation}
\label{eq:dd2-egclosed}
\alpha' \;=\; \frac{\alpha\, e^{\eta}}{\alpha\, e^{\eta} + (1 - \alpha)},
\qquad \text{equivalently} \qquad
\underbrace{\log \frac{\alpha'}{1 - \alpha'}}_{\text{log-odds}}
= \log \frac{\alpha}{1 - \alpha} + \eta .
\end{equation}
Each win adds exactly $\eta$ to the log-odds of regime $r$ within the
winner's identity. Affinity is a win counter, read through a sigmoid.

\begin{example}[Worked EG trace: uniform to pinned in seven wins]
\label{ex:dd2-egtrace}
Start at $\alpha = 1/R = 0.25$ and let the same specialist win consecutive
windows for $r$ (winning nothing else in between). With $\eta = 0.6$,
$e^{0.6} = 1.82212$. First win, by \eqref{eq:dd2-egclosed}:
\[
\alpha_1 = \frac{0.25 \times 1.82212}{0.25 \times 1.82212 + 0.75}
= \frac{0.45553}{1.20553} = 0.37787.
\]
Second win:
\[
\alpha_2 = \frac{0.37787 \times 1.82212}{0.37787 \times 1.82212 + 0.62213}
= \frac{0.68852}{1.31065} = 0.52533.
\]
Iterating (all values rounded from full-precision iteration):

\begin{center}
\begin{tabular}{lccccccccc}
\toprule
Wins $k$ & 0 & 1 & 2 & 3 & 4 & 5 & 6 & 7 & 8 \\
\midrule
$\alpha_k$ & $0.250$ & $0.378$ & $0.525$ & $0.668$ & $0.786$ & $0.870$ & $0.924$ & $\mathbf{0.957}$ & $0.976$ \\
\bottomrule
\end{tabular}
\end{center}

The pin threshold $0.95$ is crossed at win $7$. The closed form confirms it:
$\alpha_k \ge 0.95$ iff the odds satisfy
$\tfrac{\alpha_0}{1-\alpha_0} e^{\eta k} = \tfrac13 e^{0.6k} \ge 19$, i.e.\
$e^{0.6k} \ge 57$, i.e.\ $k \ge \ln 57 / 0.6 = 4.043/0.6 = 6.74$, so
$k = 7$. Seven clean wins---$7 \times 20 = 140$ in-regime days, roughly three
to four common-regime episodes---separate a cold uniform specialist from
permanent ownership.
\end{example}

\subsection{Rich-Get-Richer: The Symmetry-Breaking Engine}

The EG update alone would make ownership a random walk among near-identical
specialists. What breaks the symmetry is the coupling to training:
\emph{winners train $\Rightarrow$ training improves the head and deepens the
buffer $\Rightarrow$ a better head wins the next window more often
$\Rightarrow$ winners win again.}

Formally, let $q_m$ be the probability that the current leader on regime $r$
wins window $m$. Early on, heads are near-identical and $q_m \approx$ the
tie-breaking rate, but the first (possibly lucky) win hands the winner a
window's worth of training on $r$; under the separated-$\theta_r$ generator
this lowers its expected window regret, so $q_{m+1} > q_m$; and the affinity
bonus in \eqref{eq:dd2-winner} adds a deterministic tilt
$\lambda(\alpha_{i,r} - \alpha_{j,r})$ on top. The win process is a
reinforcing urn: the leader's log-odds climb at rate $\eta$ per win, while
losers---who do not train on $r$---suffer no foreign updates from it. Once
one specialist leads on each regime, the system is in the basin of the
covering one-per-regime assignment.

This is competitive exclusion in the sense of \citet{li2026gause}: one
species per niche at equilibrium, with reward competition as the exclusion
force---and the empirical record matches the mechanism's prediction exactly:
in \textbf{all 20 seeds}, the population converges to a covering
one-specialist-per-regime assignment, with each common regime pinned at
$\alpha \ge 0.95$.

\subsection{The Batching Question: A Pre-Registered Prediction and Its Refutation}
\label{sec:dd2-batching}

How large must $W_c$ be? The pre-registered reasoning was a single-window
concentration argument. Suppose two specialists' true mean per-day regrets
on $r$ differ by $\Delta > 0$, with per-day regrets bounded in $[0, B]$. The
window comparison errs when the worse specialist posts the lower sample
mean; by Hoeffding's inequality on the per-day regret differences (range
$2B$, mean $\Delta$),
\begin{equation}
\label{eq:dd2-hoeffding}
\Pr\bigl(\text{wrong winner}\bigr)
\;\le\; \exp\!\left( - \frac{W_c\, \Delta^2}{2 B^2} \right),
\qquad \text{so} \qquad
W_c \;\ge\; \frac{2 B^2}{\Delta^2} \ln \frac{1}{\delta}
\end{equation}
guarantees error $\le \delta$. Early in competition $\Delta$ is small by
construction (symmetric initialization), so \eqref{eq:dd2-hoeffding} blows
up, and the natural prediction follows: \emph{small windows produce noisy
winners, noisy winners scatter the EG updates, and specialization should
require large $W_c$}---batching looked necessary, with $W_c = 1$ expected to
fail outright.

\begin{keypoint}[title={The empirical refutation, stated plainly}]
The pre-registered prediction was wrong. Sweeping
$W_c \in \{1, 5, 20, 60\}$ over 20 seeds, post-reactivation regret is
\textbf{flat}: $0.0303$--$0.0309$ across the sweep, with $W_c = 1$
\emph{marginally best}. Batching within the window is not necessary for
convergence, and the analysis that said otherwise mislocated where the
averaging happens.
\end{keypoint}

The corrected mechanism story. Inequality \eqref{eq:dd2-hoeffding} bounds the
error of a \emph{single} comparison, but nothing in the architecture consumes
single comparisons: the EG recursion \eqref{eq:dd2-egclosed} sums
$\eta \cdot \mathbf{1}\{\text{win}\}$ in log-odds space across windows, and
reaching the pin requires $\approx \ln 57 / \eta \approx 6.7$ \emph{net}
wins, not one correct vote. If the better specialist wins each window with
probability $q$ above its rivals' rates, its lead in win counts---hence in
log-odds---grows linearly in the number of windows with noise
$O(\sqrt{\text{windows}})$: the law of large numbers operates \emph{across}
windows, in $\alpha$-space, regardless of $W_c$. Shrinking $W_c$ from 20 to 1
makes each vote noisier but delivers $20\times$ more votes per regime-day,
and the two effects approximately cancel; small windows also start the
rich-get-richer training loop earlier (the first win arrives after 1 day, not
20), the plausible source of $W_c = 1$'s marginal edge. The affinity is
already an exponentially-weighted tally of the win history---the EG
accumulation \emph{is} the cross-window averaging---so a second averaging
layer inside the window was redundant.

\begin{warning}[title={Do not quietly upgrade the corrected story to a prediction}]
The flat $W_c$ sweep is an empirical result on this stream (synthetic
generator, $R = 4$, $N = 4$, $\eta = 0.6$), not a theorem. The corrected
account is a post-hoc narrative consistent with the data and with standard
multiplicative-weights analyses
\citep{freund1997decision, arora2012multiplicative}, but the experiment that
would test \emph{it} (varying $\eta$ jointly with $W_c$) was not
pre-registered. One box above records a refuted prediction; this box keeps
the replacement story from enjoying a credibility it has not yet earned.
\end{warning}

\subsection{Pinning}
\label{sec:dd2-pinning}

\begin{definition}[Pin]
\label{def:dd2-pin}
When a window winner's updated affinity satisfies
$\alpha_{i^{*}, r} \ge 0.95$, regime $r$ is \textbf{pinned} to owner
$o(r) = i^{*}$, permanently: thereafter $r$'s live decisions are served by
$o(r)$ alone, training on $r$-days goes to $o(r)$ alone, and the competition
machinery for $r$ (shadow decisions, windows, EG updates) shuts down.
\end{definition}

What changes at the pin is the \emph{modality} of protection. Before it, the
owner-apparent is protected statistically: it keeps winning, so it keeps
being the one who trains, but a bad run could in principle let a rival in.
After it, protection is structural: ownership no longer consults the reward
stream at all---the assignment becomes reward-independent, not as an
assumption but as a phase the dynamics enter and cannot leave. Once all
regimes are pinned one-per-specialist (the observed endpoint in every seed),
each specialist trains only on its own regime: $p = 0$, so under HARD it
never overflows and under SOFT it suffers zero interference. The
coupon-collector clock of Proposition~\ref{prop:dd2-coupon} never starts; the
half-life \eqref{eq:dd2-halflife} never ticks. Dormancy becomes silence, with
the head and the six-episode library untouched---a distributed persistent
memory.

Two ablations calibrate how much of the benefit is the pin itself versus the
competition that produces it:
\begin{itemize}
  \item \textbf{Never-pin:} disabling Definition~\ref{def:dd2-pin} costs
        $0.0320$ versus $0.0307$ pinned post-reactivation regret---the
        residual price of merely statistical protection: perpetual shadow
        competition occasionally hands windows, hence foreign training, to
        non-owners, and serving can flicker at reactivations.
  \item \textbf{No incumbency bonus:} $\lambda = 0$ yields $0.0314$ versus
        $0.0307$. The bonus \emph{accelerates} convergence to ownership but
        does not cause it---the rich-get-richer loop breaks symmetry on its
        own. Causality lives in the win--train coupling; $\lambda$ is a
        catalyst.
\end{itemize}

\paragraph{The crisis wrinkle.}
The rare regime exposes a quantization effect in the pin rule. A crisis
episode lasts $15$--$30$ days and yields at most one complete $W_c = 20$
window---often zero---and crisis episodes are themselves few
(Section~\ref{sec:dd2-scarcity}). Across the 20 seeds, the crisis regime's
top affinity reaches only $\alpha \in [0.67, 0.87]$: by the trace of
Example~\ref{ex:dd2-egtrace}, exactly $3$--$5$ EG wins from uniform
($\alpha_3 = 0.668$, $\alpha_5 = 0.870$)---the affinity arithmetic
independently confirms that the stream does not contain the $\approx 7$
complete crisis windows formal pinning requires. The substantive outcome is
nonetheless clean: in every seed the crisis regime has a \emph{stable unique
owner}, so the covering assignment is achieved; only the formal threshold is
starved. The deployment lesson: pin on \emph{ownership stability} (the same
specialist winning a regime's windows across $k$ consecutive episodes), not
on the raw $\alpha \ge 0.95$ threshold, whose win-count requirement assumes
window-rich regimes. For rare regimes, windows---not days, not
episodes---are the binding currency.

\section{The Eleven Arms as Points in Design Space}
\label{sec:dd2-arms}

The companion experiment runs eleven arms. They are not eleven competing
architectures; they are a factorial probe of a two-axis design space, and
each arm exists to occupy a cell or test an edge.

\paragraph{Axis 1: allocation reward-coupling.} How does the regime
$\to$ learner assignment respond to the \emph{current} reward stream?
Three positions: \textbf{coupled} (chases recent performance forever),
\textbf{transient-then-structural} (\nm{}: coupled during competition,
independent after the pin), and \textbf{independent} (lottery or oracle).
Part~II showed why this axis governs forgetting: coupled allocation keeps
every learner's foreign-update probability $p > 0$ during dormancy, and both
memory models convert $p > 0$ over $250$--$450$ days into a destroyed head.

\paragraph{Axis 2: training objective.} Pooled \textbf{ERM} or the
variance-penalized \textbf{INV} of \eqref{eq:dd2-invobj} on the buffered
episodes. This axis governs the het trap of Example~\ref{ex:dd2-het}: ERM
imports episode-specific $\gamma$-weight into the probe window; INV strips
it.

\begin{center}
\begin{tabular}{llll}
\toprule
Arm & Allocation rule & Coupling & Objective \\
\midrule
A1 & monolith: train on whatever is active & coupled (degenerate) & ERM \\
A2 & MoE learned router & coupled & ERM \\
A3 & recent-performance capital shares & coupled (smooth) & ERM \\
A4 & random fixed niches & independent & ERM \\
A5 & \nm{} competition $\to$ pin & transient $\to$ structural & ERM \\
A6 & \nm{} competition $\to$ pin & transient $\to$ structural & INV \\
A7 & monolith & coupled (degenerate) & INV \\
A8a & Hedge over fixed strategies & coupled (weights only) & none (no training) \\
A8b & Hedge over learning experts & coupled & ERM \\
A9 & oracle assignment from day 1 & independent & ERM \\
A10 & oracle assignment from day 1 & independent & INV \\
\bottomrule
\end{tabular}
\end{center}

\paragraph{A1 (monolith-ERM).} One capacity-$K$ learner; the allocation rule
is the degenerate \emph{always me}: it serves and trains on whatever
$\hat r_t$ is active. Its foreign-update process during any dormancy has
$p = 1$: under HARD it is the learner of Example~\ref{ex:dd2-lru}, evicted
within one foreign block; under SOFT it retains
$0.94^{250} \approx 2 \times 10^{-7}$ over a crisis dormancy.
\emph{Prediction:} catastrophic forgetting under both memory models;
worst-tier probe regret.

\paragraph{A2 (MoE learned router).} $N$ experts plus a gate matrix
$G \in \R^{R \times N}$; on each day the gate serves expert
$j_t = \argmax_j G[\hat r_t, j]$ with probability $0.9$ and a uniformly
random expert with probability $\epsilon = 0.1$ ($\epsilon$-greedy); the
served expert trains on the day; the gate updates
\[
G[\hat r_t, j_t] \;\mathrel{+}=\; \mathrm{lr} \cdot
\bigl(\text{baseline} - \rho_{t}\bigr),
\]
with the baseline a running mean regret---below-baseline performance raises
the routing weight. This is reward-coupled allocation \emph{forever}, and
exploration alone guarantees damage: during a crisis dormancy of $D$ days
the crisis expert is dragged into foreign training at rate
$\ge \epsilon/N$ per day, so
$\E[\text{foreign updates}] \ge \epsilon D / N
= 0.1 \times (250\text{--}450)/4 = 6.25\text{--}11.25$---about one full
half-life \eqref{eq:dd2-halflife} before counting routing churn; under HARD
a single explored foreign block can evict. MoE routing in continual settings
is studied in \citet{li2024moecl}; the failure here is not expressiveness
but the absence of any reward-independent protection phase.
\emph{Prediction:} forgets despite having one expert per regime available.

\paragraph{A3 (recent-performance allocator).} All learners persist; capital
shares are a softmax of recent performance, and each learner's
\emph{training intensity equals its share}. The instructive feature is the
\textbf{softmax floor}: a softmax is strictly positive, so every learner
receives a minimum share $s_{\min} > 0$ of every day's training, foreign
days included. The floor \emph{is} the erosion rate: a dormant regime's
natural owner accrues foreign updates at rate $\ge s_{\min}$ per day,
retention $\approx 0.94^{\,s_{\min} D}$ under SOFT---graded and unavoidable.
\emph{Prediction:} graded (not cliff-like) forgetting; better than A1/A2,
worse than any covering reward-independent arm.

\paragraph{A4 (random fixed niches).} Each specialist is assigned a regime
by lottery at $t = 0$, fixed forever: perfectly reward-independent, zero
foreign updates, $\Phi \approx 0$ on every \emph{covered} regime. The defect
is coverage: with $N = R = 4$ and independent uniform draws, all four
regimes are covered with probability $4!/4^4 = 24/256 \approx 9.4\%$---in
over $90\%$ of lotteries some regime has no owner and is served at
default-decision quality forever. Reward-independence is necessary but not
sufficient: the assignment must also cover.
\emph{Prediction:} bimodal---oracle-grade retention on covered regimes,
worst-case service on uncovered ones.

\paragraph{A5 / A6 (\nm{}-ERM / \nm{}-INV).} The mechanism of
Sections~\ref{sec:dd2-specialist}--\ref{sec:dd2-competition}, differing only
in the buffer objective (ERM step vs.\ INV gradient \eqref{eq:dd2-invgrad}).
\nm{} sits at the distinguished point of axis 1: reward-coupled
\emph{exactly long enough} to discover a covering assignment (solving A4's
coverage failure), then structurally independent (solving A1--A3's
protection failure). \emph{Predictions:} both converge to covering pinned
assignments and retain through dormancy (confirmed: 20/20 seeds, common
regimes pinned $\ge 0.95$, pinned regret $0.0307$); A6 outperforms A5 on
probe windows as $\mathrm{het}$ grows, since the probe meets a fresh
$\gamma_{r,e}$ (Example~\ref{ex:dd2-het}).

\paragraph{A7 (monolith-INV).} A1's allocation with A6's objective, blocking
a confound: if \nm{}-INV beats the monolith, is it the mechanism or just the
better loss? INV reshapes \emph{what} is learned from the buffer, but the
buffer itself is evicted (HARD) or the head eroded (SOFT) by the allocation;
no objective can preserve parameters it does not control.
\emph{Prediction:} A7 $\approx$ A1 on retention---the objective axis cannot
rescue a failure on the allocation axis.

\paragraph{A8a / A8b (Hedge over fixed / learning experts).} A8a runs Hedge
\citep{freund1997decision} over a pool of \emph{fixed} strategies: weights
$w_i \propto \exp(-\eta_H \sum_s \rho^{(i)}_s)$, no training anywhere. The
allocation is reward-coupled, but with nothing trainable there is nothing to
forget: A8a is structurally immune to dormancy, and bounded by the best
fixed strategy rather than the best learned head. A8b keeps the Hedge
allocation but lets the weighted experts train on the days they are trusted
with: trainability returns, and with it the foreign-update process.
\emph{A8a vs.\ A8b isolates whether dormancy damage comes from
reward-coupled allocation per se or from training under it} (the latter).
\emph{Predictions:} A8a immune to forgetting but capped within episodes;
A8b recovers in-episode adaptation and re-imports forgetting.

\paragraph{A9 / A10 (oracle assignment, ERM / INV).} The regime $\to$
specialist map is given correctly at $t = 0$---pinning before day one---and
each owner trains only on its regime, with ERM (A9) or INV (A10). These are
the ceilings: reward-independent, covering, zero discovery cost. The
A5/A6 vs.\ A9/A10 gap prices the \emph{discovery phase} (contested serving,
occasional wrong winners, foreign training before ownership settles); the
A9 vs.\ A10 gap reads off the pure objective effect with allocation held
perfect, mirroring A5 vs.\ A6.

\begin{keypoint}[title={The arms are an experiment about classes, not architectures}]
Each arm represents a \emph{cell} in (coupling $\times$ objective) space,
and every headline conclusion is a contrast between cells with one axis held
fixed: A5 vs.\ A2/A3 tests allocation class at fixed objective; A6 vs.\ A5
and A10 vs.\ A9 test objective class at fixed allocation; A7 shows the
objective axis cannot compensate for the allocation axis; A8a vs.\ A8b shows
the damage mechanism is training-under-coupling, not coupling itself; A4
shows reward-independence without coverage is not enough; A9/A10 price
discovery. Substituting a different router, softmax temperature, or fixed
pool would not move any cell. The claim under test is never ``\nm{} beats
architecture X''; it is that \emph{retention through dormancy requires an
assignment that is both reward-independent and covering, and
competition-then-pinning is a cheap emergent source of exactly that
conjunction}---with the objective axis an orthogonal dial governing
probe-window robustness to episode heterogeneity, not retention.
\end{keypoint}

\begin{remark}[Where the empirical record currently stands]
The companion experiments confirm the \nm{} cell directly: 20/20 seeds reach
a covering one-per-regime assignment; common regimes pin at
$\alpha \ge 0.95$; crisis acquires a stable unique owner at
$\alpha \in [0.67, 0.87]$ without a formal pin; never-pinning costs $0.0320$
vs.\ $0.0307$; $\lambda = 0$ costs $0.0314$ vs.\ $0.0307$; and the $W_c$
sweep is flat at $0.0303$--$0.0309$ with $W_c = 1$ marginally best---the
latter refuting the pre-registered batching prediction of
Section~\ref{sec:dd2-batching}. The refutation is reported with the same
prominence as the confirmations; a design-space map is only as trustworthy
as its falsified cells.
\end{remark}
% ============================================================================
% PART IV: THE THEORY IN FULL (author-written; precision-critical)
% ============================================================================
\part{The Theory in Full}

This part assembles the foundations of Part~I into the four results that
organize the \nm{} account: the episode-transfer lemma controlling the
invariance gap, the eviction-rate lemma controlling the forgetting
deficit, the decomposition proposition that yields the $2{\times}2$
prediction (including its super-additive signature), and the break-even
proposition that prices retention against carrying costs. We close with
the Gaussian construction that exhibits ERM's across-episode failure in
closed form, and with the honest remarks on what the coupled competition
dynamics are \emph{not} proven to do. Throughout, the register is that of
the GAUSE companion theory \cite{li2026gause}: idealized-model results
that explain mechanisms and generate falsifiable orderings, with the
quantitative weight carried by the controlled experiments of Part~V.

\section{Assumptions, Stated and Audited}
\label{sec:dd4-assumptions}

\begin{definition}[Episode risk]
\label{def:dd4-risk}
For regime $r$, episode environment $\xi$ with law $P_\xi$ on
$\mathcal{X} \times \mathcal{Y}$, and head $w$, the episode risk is
$L_\xi(w) = \E_{(x,y) \sim P_\xi}[\rho(Xw;\, y)]$, with $\rho$ the
decision regret of Part~I, $\rho \in [0, B]$,
$B = 2 k\, w_{\max}\, y_{\max}$.
\end{definition}

\begin{definition}[Assumption A1: episode exchangeability]
\label{def:dd4-a1}
For each regime $r$ there is a hyper-distribution $Q_r$ over episode
environments such that the historical episodes
$\xi_{r,1}, \dots, \xi_{r,E_r}$ and the deployed episode $\xi_{r,e'}$
are i.i.d.\ draws from $Q_r$. Write
$\mu_r(w) = \E_{\xi \sim Q_r} L_\xi(w)$ and
$\sigma_r^2(w) = \Var_{\xi \sim Q_r} L_\xi(w)$.
\end{definition}

\begin{warning}
A1 idealizes twice. \emph{Independence} fails when macro memory links
episodes (the 2011 crisis is causally downstream of 2008); \emph{identical
distribution} fails under secular drift (market microstructure in 1998
differs from 2020 in ways no regime label captures). We use A1 exactly as
Inv-PnCO uses its environment assumption \cite{yan2026invpnco}: as the
formal rendering of ``the next crisis is a fresh draw from the crisis
distribution,'' with the experiments designed so that A1 holds by
construction (the synthetic generator) and the limitations section of the
companion paper carrying the real-world caveat. Nothing in this part
silently strengthens A1.
\end{warning}

\begin{definition}[Assumption A2: dormancy and foreign updates]
\label{def:dd4-a2}
While regime $r$ is dormant for $D$ steps, an allocation mechanism $A$
routes learning. $A$ is \emph{reward-driven} if the specialist holding
$r$ is made to learn a currently active regime with per-step probability
at least $p > 0$; $A$ is \emph{reward-independent} if the assignment of
$r$ is measurable with respect to information independent of the
realized reward stream and, at convergence, routes to $r$'s owner only
regimes in its fixed niche set.
\end{definition}

The $p > 0$ clause is not vacuous bookkeeping: it is verified \emph{by
construction} for each reward-driven arm in our inventory (the monolith
learns the active regime with probability $1$; the router's
$\epsilon$-greedy exploration plus reward-chasing gate gives
$p \ge \epsilon / N$; the recent-performance allocator's softmax floor
gives $p \ge v_{\min} > 0$), and it is exactly the quantity a fund can
audit in its own allocator: \emph{what fraction of a losing desk's
capacity is reassigned per period?}

\section{The Invariance Gap}
\label{sec:dd4-invgap}

\begin{lemma}[Episode transfer]
\label{lem:dd4-transfer}
Under A1, for any fixed head $w$ and any $\delta \in (0, 1]$, with
probability at least $1 - \delta$ over $\xi_{e'} \sim Q_r$,
\[
  L_{\xi_{e'}}(w) \;\le\; \mu_r(w) +
  \sigma_r(w)\,\sqrt{\tfrac{1-\delta}{\delta}}
  \;\le\; \mu_r(w) + \frac{\sigma_r(w)}{\sqrt{\delta}} .
\]
\end{lemma}

\begin{proof}
Cantelli's inequality (Part~I) applied to the real random variable
$Z = L_\xi(w)$, $\xi \sim Q_r$, which has mean $\mu_r(w)$ and variance
$\sigma_r^2(w)$: for $\tau > 0$,
$\Pr[Z - \mu_r \ge \tau] \le \sigma_r^2 / (\sigma_r^2 + \tau^2)$.
Setting the right side to $\delta$ and solving,
$\tau = \sigma_r \sqrt{(1-\delta)/\delta}$, which is at most
$\sigma_r / \sqrt{\delta}$. No boundedness or distributional shape is
used; only the existence of the second moment, guaranteed by
$Z \in [0, B]$.
\end{proof}

\begin{keypoint}
The bound is monotone in \emph{both} $\mu_r(w)$ and $\sigma_r(w)$, and in
nothing else. Any training objective that drives the pair
$(\mu_r, \sigma_r)$ toward the lower-left of its Pareto frontier
optimizes the guarantee; the \nm{} objective
$\Var_e + \beta\,\E_e$ is the Lagrangian scalarization of exactly that
frontier, with $\beta$ the exchange rate. For each tail level $\delta$
there is a $\beta(\delta)$ whose minimizer over any convex model class
attains the frontier point optimizing the Lemma's bound (standard
scalarization; for non-convex classes the scalarization reaches the
convex hull of the frontier, a caveat we inherit silently from every
penalty method in the invariance literature \cite{krueger2021rex}).
\end{keypoint}

\begin{lemma}[Estimation price of few episodes]
\label{lem:dd4-estimation}
Fix $w$. With $E_r$ training episodes drawn i.i.d.\ from $Q_r$ and
$\hat\mu_r(w)$, $\hat\sigma_r(w)$ the empirical mean and standard
deviation of $\{L_{\xi_e}(w)\}_{e=1}^{E_r}$: for any
$\delta \in (0, 1)$, with probability at least $1 - \delta$,
\[
  |\hat\mu_r(w) - \mu_r(w)| \;\le\; B \sqrt{\frac{\log(2/\delta)}{2 E_r}},
  \qquad
  |\hat\sigma_r(w) - \sigma_r(w)| \;\le\; B \sqrt{\frac{2 \log(2/\delta)}{E_r - 1}} .
\]
\end{lemma}

\begin{proof}
The mean bound is Hoeffding's inequality (Part~I) for i.i.d.\
$[0, B]$-valued variables. For the deviation bound, note that the
empirical standard deviation of $[0,B]$-bounded variables is a function
of the sample with bounded differences: replacing one $L_{\xi_e}$ changes
$\hat\sigma_r$ by at most $B / \sqrt{E_r - 1}$ (each summand
$(L_e - \hat\mu)^2$ moves by at most $B^2$ terms damped by the
$1/(E_r{-}1)$ normalization; a one-line calculation, or see
\cite{maurer2009empirical} for the sharper empirical-Bernstein form).
McDiarmid's inequality then gives the stated tail, and
$\E[\hat\sigma_r] \le \sigma_r$ plus the same concentration controls the
bias direction.
\end{proof}

\begin{warning}[title={Important: the plug-in caveat}]
Lemmas~\ref{lem:dd4-transfer}--\ref{lem:dd4-estimation} hold for
\emph{fixed} $w$. The trained head $\hat w_r$ is a function of the
training episodes, so applying them at $w = \hat w_r$ is a plug-in
heuristic, not a uniform-convergence guarantee; making it uniform costs
a complexity term over the head class (linear heads in $d = 20$
dimensions would admit a standard covering-number argument) that we do
not need for the mechanism claims and therefore do not carry. This is
the same register at which Inv-PnCO deploys its Theorem~1, and we flag
rather than hide it. The practical content survives untouched: with
$E_r = 6$ episodes and $B = 0.4$, the mean's estimation radius at
$\delta = 0.1$ is $B\sqrt{\log 20 / 12} \approx 0.20$---half the regret
range---which is the formal statement of \emph{why} crisis-grade
invariance cannot be certified from six episodes alone and why the
sub-episode and cross-asset enlargements of Part~II are design defaults
rather than afterthoughts.
\end{warning}

\begin{corollary}[Bridge to divergence-style guarantees]
\label{cor:dd4-bridge}
If a representation $\Phi$ renders episode laws $\varepsilon$-invariant
in KL---$\KL(P_\xi(\cdot \mid \Phi)\, \|\, \bar P_r(\cdot \mid \Phi))
\le \varepsilon$ for $Q_r$-almost-every $\xi$---then every head $w$
factoring through $\Phi$ satisfies
$\sup_\xi |L_\xi(w) - \bar L_r(w)| \le B\sqrt{\varepsilon/2}$, hence
$\sigma_r(w) \le B\sqrt{\varepsilon/2}$, and
Lemma~\ref{lem:dd4-transfer} applies with that variance.
\end{corollary}

\begin{proof}
Immediate from the divergence-to-regret bridge of Part~I
($|L_P - L_{P'}| \le B\,\TV \le B\sqrt{\KL/2}$, Pinsker) applied
between each $P_\xi$ and the mixture $\bar P_r$; a uniformly
$c$-close-to-mean random variable has standard deviation at most $c$.
\end{proof}

This corollary is the formal connective tissue between the Inv-PnCO
guarantee language (divergences between solution distributions
\cite{yan2026invpnco}) and the regret currency of our decomposition: one
Pinsker step, at the explicit and admittedly loose price $B$.

\section{The Forgetting Deficit}
\label{sec:dd4-forget}

\begin{lemma}[Eviction rate under reward-driven allocation]
\label{lem:dd4-eviction}
Hard memory, capacity $K$; during $r$'s dormancy let $W_a \ge K$ distinct
regimes be active, let each dormant step independently deliver a foreign
learning assignment to $r$'s owner with probability $\ge p$, and
conditional on an assignment let the regime be drawn from the active set
with each regime's probability $\ge q_{\min} > 0$. Let $U_D$ be the
number of \emph{distinct} previously-unheld regimes learned in $D$ steps.
Then $r$'s head (and its episode buffer) is evicted on
$\{U_D \ge K\}$, and
\[
  \Pr[U_D < K] \;\le\;
  \Pr\!\big[\mathrm{Binom}(D,\ p\, q_{\min} (W_a - K + 1)) < K\big]
  \;\le\; \exp\!\Big(\!-2 D \big(p\,q_{\min}(W_a{-}K{+}1) -
  \tfrac{K}{D}\big)_+^2\Big),
\]
so $\Pr[\text{eviction}] \to 1$ exponentially in $D$. Moreover the
expected eviction time is bounded by the coupon-collector sum
\[
  \E[\text{time to } U = K] \;\le\; \frac{1}{p}
  \sum_{j=0}^{K-1} \frac{1}{q_{\min}\,(W_a - j)} .
\]
Under the soft model with per-foreign-update decay $\rho_{\mathrm{int}}$,
the retained strength after $N_D$ foreign updates is
$(1 - \rho_{\mathrm{int}})^{N_D}$, and $N_D \ge \mathrm{Binom}(D, p)$
stochastically, so $\E[\text{strength}] \le
(1 - p\,\rho_{\mathrm{int}})^{D} \to 0$.
\end{lemma}

\begin{proof}
While fewer than $K$ distinct foreign regimes have been acquired, at
least $W_a - K + 1$ of the active regimes are still unheld, so each step
delivers a \emph{new} distinct foreign regime with probability at least
$p\, q_{\min} (W_a - K + 1)$. The count of such successes in $D$ steps
thus stochastically dominates the stated binomial while $U < K$, giving
the first inequality; the second is Hoeffding's bound for the binomial
lower tail. The waiting-time bound is the sum of the $K$ geometric means
with success probabilities $p\, q_{\min}(W_a - j)$, $j = 0, \dots,
K{-}1$ (the coupon-collector computation of Part~II). LRU evicts $r$
precisely when $K$ distinct foreign heads have displaced it, since the
dormant $r$ is by definition the least recently used. For the soft
model, each foreign update multiplies $r$'s strength by
$(1 - \rho_{\mathrm{int}})$; taking expectations over the binomial count
gives the product form via $\E[(1-\rho)^{\mathrm{Binom}(D,p)}] =
(1 - p\rho)^D$.
\end{proof}

\begin{example}[Numbers at our calibration]
$K = 2$, $W_a = 3$ active common regimes, $p \approx 1$ (the monolith
always learns the active regime), $q_{\min} \approx 1/3$ at block scale:
the expected eviction time is on the order of two block lengths
($\approx 50$--$100$ days), and at crisis dormancy $D \approx 280$ days
the eviction probability is numerically indistinguishable from $1$. This
is why the E3 dormancy sweep finds the retention gap already saturated
at $D = 63$--$126$: the eviction event is binary, happens early in the
dormancy, and cannot happen twice.
\end{example}

\begin{proposition}[Dichotomy: reward-independence zeroes the deficit]
\label{prop:dd4-dichotomy}
Define the forgetting deficit
$\Phi(D; A) = \Pr[\text{$r$'s head lost at reactivation under } A]$
(hard) or $1 - \E[\text{retained strength}]$ (soft). Then:
(i) for any reward-driven $A$ satisfying A2,
$\Phi(D; A) \to 1$ as $D \to \infty$ at the rate of
Lemma~\ref{lem:dd4-eviction}; (ii) for any reward-independent covering
$A$, $\Phi(D; A) = 0$ for all $D$ in the hard model, and in the soft
model the owner's strength in $r$ decays only through the regime's own
drift, with \emph{zero} contribution from foreign updates.
\end{proposition}

\begin{proof}
(i) is the lemma. (ii): coverage gives an owner holding $r$ within
capacity; reward-independence means the assignment map does not respond
to the realized reward stream, so during $r$'s dormancy the owner is
assigned only regimes in its fixed niche set. If its niche set is
$\{r\}$ (the converged one-per-regime case) it receives no updates at
all; if its set contains other regimes, those are by construction
within capacity, so no eviction is triggered ($U_D = 0$ counts only
\emph{unheld} acquisitions) and---in the implemented soft model, where
decay attaches to foreign-\emph{regime} updates---no interference is
triggered either. The head and the episode buffer at reactivation are
bit-identical to their state at dormancy onset.
\end{proof}

\begin{keypoint}
The dichotomy is qualitative, not quantitative: the two classes are
different \emph{functions} of $D$ (one $\to 1$, one $\equiv 0$), which
is what licenses the experimental design---arm classes, not parameter
sweeps, separate them, and the E1 table confirms the class boundary at
$p \sim 10^{-6}$ while E3 confirms the saturation shape. Coverage is the
second conjunct and not decorative: random fixed niches are
reward-independent yet retain only what they own, and their doubled seed
variance in E1 is the measured price of leaving coverage to chance.
\end{keypoint}

\section{The Decomposition and the $2{\times}2$}
\label{sec:dd4-decomp}

\begin{definition}[Relearning cost]
\label{def:dd4-crelearn}
$C_{\mathrm{relearn}}$ is the expected excess of a from-prior head's
mean probe-window regret over a converged retained head's, for the same
regime and probe length $N_p$. It is an empirical constant of the
learning problem (model class, optimizer, signal-to-noise), not a
theorem; the E6 audit shows it collapsing as signal strength rises,
which is the honest reason the entire mechanism has a signal-to-noise
boundary.
\end{definition}

\begin{proposition}[Post-reactivation regret decomposition]
\label{prop:dd4-decomp}
Under A1--A2, when regime $r$ reactivates after dormancy $D$ into a
fresh episode, served under allocation $A$ by a head trained on $E_r$
episodes, then at tail level $\delta$, with the convention that the
invariance terms are evaluated at the served retained head $\hat w_r$,
\[
  \E[\rho_{\mathrm{react}}(r)] \;\le\;
  \underbrace{\hat\mu_r(\hat w_r) + \frac{\hat\sigma_r(\hat w_r)}{\sqrt\delta}
  + O\!\Big(B\sqrt{\tfrac{\log(1/\delta)}{E_r}}\Big)}_{\Ginv}
  \;+\;
  \underbrace{\Phi(D; A)\; C_{\mathrm{relearn}}}_{\Gforget},
\]
where the $O(\cdot)$ collects the estimation radii of
Lemma~\ref{lem:dd4-estimation}.
\end{proposition}

\begin{proof}
Condition on the retention event $R$ ($\Pr[R] = 1 - \Phi$). On $R$ the
served head is $\hat w_r$; Lemma~\ref{lem:dd4-transfer} bounds its
fresh-episode risk by $\mu_r + \sigma_r/\sqrt\delta$, and
Lemma~\ref{lem:dd4-estimation} converts population to empirical
quantities at the stated radii. On $R^c$ ($\Pr = \Phi$) the served head
is fresh; by Definition~\ref{def:dd4-crelearn} its probe regret exceeds
the retained benchmark by $C_{\mathrm{relearn}}$ in expectation, and the
benchmark itself is bounded by the same invariance terms. Total
expectation: $\E[\rho_{\mathrm{react}}] \le (1-\Phi)\,\Ginv + \Phi\,
(\Ginv + C_{\mathrm{relearn}}) = \Ginv + \Phi\, C_{\mathrm{relearn}}$.
\end{proof}

\begin{proposition}[Independent controls; super-additivity where it
matters]
\label{prop:dd4-additive}
$\Gforget$ is a function of $(A, D, K, \text{memory model})$ alone and
is zeroed by any reward-independent covering $A$
(Proposition~\ref{prop:dd4-dichotomy}) regardless of the training
objective; $\Ginv$ is a function of $(\text{objective}, Q_r, E_r)$ alone
and is optimized by the variance-penalized objective
(Lemma~\ref{lem:dd4-transfer} discussion) regardless of $A$. The
\emph{bound} is therefore additive in the two design axes. The
\emph{realized} interaction is predicted super-additive in one specific
cell: under reward-driven $A$ with $\Phi \approx 1$, the served head at
reactivation is fresh whatever objective trained the evicted one, so
improving $\Ginv$ moves nothing observable---the invariance axis is
inert in the reward-driven row.
\end{proposition}

\begin{remark}[The measured interaction]
Define $I = (\text{A1} - \text{A5}) + (\text{A1} - \text{A7}) -
(\text{A1} - \text{A6})$ on post-reactivation regret: $I = 0$ is
additivity, $I < 0$ is synergy. The companion experiments measure
$I = -0.0015 \pm 0.0005$ per-seed (95\% CI excluding zero), realized
almost entirely through $\text{A1} - \text{A7} \approx 0$
($p = 0.88$)---the inert cell, exactly where
Proposition~\ref{prop:dd4-additive} puts it. We are careful about
logical direction: the additive \emph{bound} motivates but does not
entail the super-additive \emph{realization}; the prediction that A7
specifically would sit at A1's level was pre-registered from the
mechanism (an evicted head serves nobody) and is the account's sharpest
falsifiable signature---a ``both tricks help'' story without the
two-clock structure would have A7 strictly between A1 and A6.
\end{remark}

\section{Break-Even Dormancy: Pricing Retention}
\label{sec:dd4-breakeven}

\begin{proposition}[The retention window]
\label{prop:dd4-breakeven}
Charge carrying cost $c$ per idle specialist-day and let $V$ convert one
unit of mean probe regret to value over the $N_p$-day probe. Over one
dormancy--reactivation cycle of length $D$, retaining $r$'s specialist
pays iff
\[
  c\, D \;\le\; \Phi(D; A_{\mathrm{rd}})\;
  C_{\mathrm{relearn}}\, N_p\, V .
\]
Since $\Phi$ is increasing and saturates at $1$ while $cD$ grows
linearly, the solution set is an interval $[D_{\min}, D_{\max}]$
(possibly empty): $D_{\min}$ where $\Phi$ ramps (eviction becomes likely
enough to fear), $D_{\max} = C_{\mathrm{relearn}} N_p V / c$ where
cumulative carrying cost overtakes the one-shot saving. The window is
empty iff $c > \max_D \Phi(D) C_{\mathrm{relearn}} N_p V / D$, i.e.,
for carrying costs above a threshold set entirely by measurable
quantities.
\end{proposition}

\begin{example}[At the E3 calibration]
$C_{\mathrm{relearn}} \approx 0.004$ regret per probe day, $N_p = 15$:
the cycle saving is $\approx 0.06\, V$. At a carrying cost of $1\%$ of
the daily probe saving ($c = 4 \times 10^{-5} V$), $D_{\max} \approx
1{,}500$ trading days and $D_{\min} \approx 40$--$60$ (where E3 shows
$\Phi$ ramping); at $10\%$, $D_{\max} \approx 150$ days---tighter than
one crisis cycle, so a fund at that cost level should \emph{not} retain
a dedicated crisis specialist and should plan to relearn; the indifference
threshold sits near $15\%$. The proposition's value is not a universal
verdict but the conversion of an organizational argument
(``keep the crisis desk?'') into three measurable numbers.
\end{example}

\section{What the Competition Dynamics Are Not Proven to Do}
\label{sec:dd4-honesty}

Two honest closings. First, \emph{we prove no convergence theorem for
the coupled winner-take-all dynamics}: specialists' window wins depend on
heads, heads on training, training on wins---an endogenous feedback for
which adversarial-regret guarantees of the EG update do not apply (Part~I)
and for which we know of no clean fixed-point argument covering the full
mechanism. What we have is the GAUSE precedent (the same dynamics
converge to covering one-per-regime assignments across all its domains
and protocols \cite{li2026gause}), our own twenty-for-twenty convergence
record with the crisis-regime affinity wrinkle reported, and the
single-window concentration remark---whose pre-registered implication
(batching necessary) the $W_c$ ablation \emph{refuted}, teaching the
corrected lesson that cross-window EG accumulation is the operative
averager. Second, \emph{the decomposition's two terms are bounds, not
identities}: a mechanism could in principle beat the bound's additivity
(and ours measurably does, in the predicted direction and cell). The
theory's role is to say which knob owns which term and where each term
vanishes---claims the eleven-arm design tests directly---not to predict
regret to the fourth decimal. Where the theory and the experiments
disagreed, this document reports the experiment and revises the story,
twice (batching; the high-SNR prototype that became E6).
% ============================================================================
% PART V: EXPERIMENTAL EVIDENCE IN FULL
% Fragment file: deepdive_part5.tex (no preamble; included from master)
% All numbers in this part are taken verbatim (or rounded) from the archived
% result files e0_structure_diagnostic.json, e1_synth.json,
% e1s_stitched_crypto.json, e2_capacity_sweep.json, e3_dormancy_sweep.json,
% e4_heterogeneity_sweep.json, e5_ablations.json, e6_snr_audit.json.
% ============================================================================

\part{Experimental Evidence in Full}
\label{sec:dd5-part}

The theory of the preceding parts makes a small number of sharp, falsifiable
claims: that a reward-independent niche assignment retains dormant-regime
competence where every reward-coupled allocator relearns; that an invariance
objective inside a niche buys robustness that an invariance objective inside a
monolith cannot; that the two mechanisms interact super-additively; and that
all of this is conditional on the environment actually containing persistent,
re-activating conditional structure. This part subjects each claim to its
experiment. We present the numbers exhaustively---every arm, every metric,
every pre-registered comparison with its raw and corrected $p$-value---and we
present the null results with the same care as the positive ones, because the
protocol's credibility rests on its ability to return ``no'' (Sections
\ref{sec:dd5-e0e1s} and \ref{sec:dd5-e6}).

\section{Methodology and Statistical Apparatus}
\label{sec:dd5-method}

\subsection{The arms}
\label{sec:dd5-arms}

Eleven arms share an identical decision layer (top-$K$ portfolio construction
over $n=20$ assets from predicted returns) and identical data per seed; they
differ only in how predictive capacity is organized and trained.

\begin{table}[h]
\centering
\caption{The eleven arms. ``Assignment'' is how specialists (or capacity
slots) are bound to regimes; ``Objective'' is the within-specialist training
loss. Oracle arms receive the true regime$\to$specialist map and serve as
skylines, not competitors.}
\label{tab:dd5-arms}
\small
\begin{tabular}{llll}
\toprule
Arm & Name & Assignment & Objective \\
\midrule
A1 & monolith-ERM & none (single model) & ERM \\
A2 & MoE router & learned gate, reward-driven & ERM \\
A3 & recent-perf allocator & capital $\propto$ recent performance & ERM \\
A4 & random fixed niches & random, frozen at init & ERM \\
A5 & RISP-ERM & learned, dormancy-pinned & ERM \\
A6 & RISP-INV (full system) & learned, dormancy-pinned & invariance \\
A7 & monolith-INV & none (single model) & invariance \\
A8a & Hedge over fixed experts & exponential weights, frozen experts & none \\
A8b & Hedge over learners & exponential weights, learning experts & ERM \\
A9 & oracle-pinned-ERM & oracle map, pinned & ERM \\
A10 & oracle-pinned-INV & oracle map, pinned & invariance \\
\bottomrule
\end{tabular}
\end{table}

The synthetic market has $n=20$ assets, $d=20$ features, $R=4$ latent regimes
with dormancy/reactivation dynamics, per-specialist capacity $K=2$, and
$20$ seeds per configuration. Evaluation is restricted to the second half of
each stream so that all arms are past their transient.

\subsection{The seeded paired design}
\label{sec:dd5-paired}

Every seed $s \in \{1,\dots,20\}$ fixes the entire data-generating draw: the
regime schedule, the per-regime parameters, the feature noise, and the
realized returns. All eleven arms are then run on the \emph{same} stream for
seed $s$. Consequently the between-seed component of variance---which regime
schedule was drawn, how harsh the reactivations happened to be---is common to
all arms, and arm differences within a seed are estimates of the treatment
effect alone. This is the experimental-design analogue of common random
numbers: it does not change the means but sharply reduces the variance of
\emph{contrasts}.

We nonetheless report unpaired Welch tests on the per-seed means
(Section~\ref{sec:dd5-welch}) rather than paired $t$-tests. This is
deliberately conservative: the paired test would have strictly more power
(the seed-level draws are positively correlated across arms), so any
comparison that survives the unpaired test would also survive the paired one,
while the converse need not hold. Where a paired analysis materially changes
the reading (the $\lambda$ ablation in Section~\ref{sec:dd5-e5}), we say so
explicitly.

\subsection{The three metrics}
\label{sec:dd5-metrics}

All metrics are means of the per-day \emph{decision regret}
\[
r_t \;=\; V^{\star}_t - V_t,
\]
where $V_t$ is the realized value of the arm's top-$K$ portfolio on day $t$
and $V^{\star}_t$ the value of the best top-$K$ portfolio in hindsight for
that day. Regret is measured in the decision space, not the prediction space,
following the smart-predict-then-optimize critique
\citep{elmachtoub2022smart}: prediction error and decision quality can
dissociate, and our research question lives entirely on the decision side.

Let $\mathcal{T}$ be the evaluation window (second half of the stream). A
\emph{reactivation event} is a day on which a regime that has been dormant
for at least $90$ days re-enters the schedule. The \emph{probe window} of an
event is its first $15$ evaluation days. Then:
\begin{align}
\texttt{post\_react} &= \text{mean of } r_t \text{ over the union of probe
windows in } \mathcal{T}, \label{eq:dd5-postreact}\\
\texttt{steady} &= \text{mean of } r_t \text{ over }
\mathcal{T} \setminus \{\text{probe windows}\}, \label{eq:dd5-steady}\\
\texttt{overall} &= \text{mean of } r_t \text{ over all of } \mathcal{T}.
\label{eq:dd5-overall}
\end{align}
\texttt{post\_react} is the primary endpoint: it isolates the days on which
retention can matter at all. \texttt{steady} measures the price paid for the
mechanism when nothing is reactivating, and \texttt{overall} aggregates both.
Each seed contributes one number per metric per arm; all inference is on
these $n=20$ per-seed means, never on the (serially dependent) daily values.

\subsection{Welch's $t$-test}
\label{sec:dd5-welch}

For arms $a$ and $b$ with per-seed means $\bar{x}_a, \bar{x}_b$, sample
variances $s_a^2, s_b^2$, and $n_a = n_b = 20$, the Welch statistic is
\begin{equation}
t \;=\; \frac{\bar{x}_a - \bar{x}_b}
{\sqrt{\dfrac{s_a^2}{n_a} + \dfrac{s_b^2}{n_b}}},
\label{eq:dd5-welch-t}
\end{equation}
referred to a $t$ distribution with the Welch--Satterthwaite degrees of
freedom
\begin{equation}
\nu \;=\;
\frac{\left(\dfrac{s_a^2}{n_a} + \dfrac{s_b^2}{n_b}\right)^{2}}
{\dfrac{(s_a^2/n_a)^2}{n_a - 1} + \dfrac{(s_b^2/n_b)^2}{n_b - 1}}.
\label{eq:dd5-welch-df}
\end{equation}
We use the unequal-variance form rather than the pooled Student $t$ because
variance heterogeneity across arms is not an edge case here---it is a
\emph{finding}. A4's defining signature is inflated between-seed variance
(its post-reactivation 95\% CI half-width is $0.00126$ against A5's
$0.00070$, a factor of $1.8$; Section~\ref{sec:dd5-e1-a4}), and A6's
high-SNR collapse comes with CI half-widths four to six times A1's
(Section~\ref{sec:dd5-e6}). Pooling variances across such pairs would
miscalibrate exactly the comparisons we care most about. With $n_a=n_b=20$
and roughly comparable variances, $\nu$ ranges over roughly $30$--$38$ in our
comparisons; the test is essentially a $z$-test with a mild small-sample
correction.

\subsection{Holm--Bonferroni and a proof that it controls FWER}
\label{sec:dd5-holm}

The E1 analysis involves a pre-registered family of $m = 11$ pairwise
comparisons (Table~\ref{tab:dd5-e1-tests}). With $11$ tests at level $0.05$
each, the chance of at least one false rejection under the global null is
bounded only by $0.55$ by Bonferroni's inequality---useless. We therefore
control the family-wise error rate (FWER) with Holm's step-down procedure,
which dominates plain Bonferroni at no assumption cost.

\begin{definition}[Holm step-down procedure]
\label{def:dd5-holm}
Order the $p$-values $p_{(1)} \le p_{(2)} \le \cdots \le p_{(m)}$ with
corresponding hypotheses $H_{(1)},\dots,H_{(m)}$. Let
\[
k^\ast = \min\Bigl\{ k : p_{(k)} > \frac{\alpha}{m - k + 1} \Bigr\}
\]
(with $k^\ast = m+1$ if no such $k$ exists). Reject
$H_{(1)},\dots,H_{(k^\ast - 1)}$ and retain the rest. Equivalently, report
adjusted $p$-values
$\tilde{p}_{(k)} = \max_{j \le k}\, \min\{(m-j+1)\, p_{(j)},\, 1\}$
and reject where $\tilde{p} \le \alpha$. The \texttt{holm\_p} entries in our
result files are these adjusted values.
\end{definition}

\begin{proposition}[Holm controls FWER at level $\alpha$]
\label{prop:dd5-holm-fwer}
Under arbitrary dependence among the $p$-values, the probability that Holm's
procedure rejects at least one true null hypothesis is at most $\alpha$.
\end{proposition}

\begin{proof}
Let $I_0$ be the index set of true nulls and $m_0 = |I_0|$; assume $m_0 \ge 1$
(otherwise there is nothing to prove). Suppose the procedure rejects at least
one true null, and let $k$ be the \emph{first} step (in the sorted order) at
which a true null is rejected. All hypotheses rejected at steps
$1,\dots,k-1$ are then false nulls, of which there are $m - m_0$ in total;
hence $k - 1 \le m - m_0$, i.e.\ $k \le m - m_0 + 1$, and therefore
\[
\frac{\alpha}{m - k + 1} \;\le\; \frac{\alpha}{m_0}.
\]
Because step $k$ rejected, $p_{(k)} \le \alpha/(m-k+1) \le \alpha/m_0$, and
since the true null rejected at step $k$ has $p$-value $p_{(k)}$,
\[
\min_{i \in I_0} p_i \;\le\; p_{(k)} \;\le\; \frac{\alpha}{m_0}.
\]
So the event ``at least one true null is rejected'' implies the event
$\{\min_{i \in I_0} p_i \le \alpha/m_0\}$. For each true null,
$\Pr(p_i \le \alpha/m_0) \le \alpha/m_0$ (validity of the individual test),
so by the union bound
\[
\Pr\Bigl(\min_{i \in I_0} p_i \le \frac{\alpha}{m_0}\Bigr)
\;\le\; m_0 \cdot \frac{\alpha}{m_0} \;=\; \alpha. \qedhere
\]
\end{proof}

Multiple-testing discipline is not optional in this literature: the data
mining of trading-strategy families is precisely the failure mode documented
by the reality-check and SPA literature
\citep{white2000reality,hansen2005spa}, by the deflated Sharpe ratio
\citep{bailey2014deflated}, and by the backtesting critiques of
\citet{harvey2016backtesting} and \citet{lopez2018advances}. Our family is
small ($11$), fixed in advance, and corrected; that is the minimum standard
those papers demand.

\subsection{Confidence intervals and the choice of 20 seeds}
\label{sec:dd5-ci}

All reported intervals are per-arm 95\% confidence intervals on the mean over
seeds,
\begin{equation}
\bar{x} \;\pm\; t_{19,\,0.975}\, \frac{s}{\sqrt{20}}, \qquad
t_{19,\,0.975} = 2.093,
\label{eq:dd5-ci}
\end{equation}
where $s$ is the between-seed standard deviation. The \texttt{ci95} fields in
the result files are the half-widths of these intervals.

Why $20$ seeds? Back out the per-seed SD from Eq.~\eqref{eq:dd5-ci}: on the
primary endpoint the major arms have $s \approx 0.0014$--$0.0027$ (A1:
$0.0020$; A5: $0.0015$; A6: $0.0014$; A4: $0.0027$). The tightest
pre-registered contrast is the invariance axis A6 vs.\ A5, with observed
difference $0.00142$. The Welch standard error for that pair is
$\sqrt{(0.0015^2 + 0.0014^2)/20} \approx 0.00046$, giving $t \approx 3.1$ and
the observed raw $p = 0.0061$---comfortably detectable, with power
$\approx 0.85$ at $\alpha = 0.05$ for an effect of that size. Halving the
seed count would have left the invariance axis at the edge of detectability
after Holm correction; doubling it would have bought little, since the
binding constraint on the smallest effects we care about (e.g.\ A6 vs.\ A10)
is that they are \emph{predicted to be null}. Twenty seeds is the point where
every effect we predicted to exist is powered and every effect we predicted
not to exist is bounded by a usefully narrow CI.

\subsection{The noise-floor convention}
\label{sec:dd5-floor}

Raw regret levels conflate two things: irreducible per-day noise (even the
oracle-assigned, invariance-trained A10 cannot drive regret to zero in a
stochastic market) and the mechanism-dependent excess above that floor. We
adopt the convention
\begin{equation}
\text{floor} \;=\; \texttt{steady}(\text{A10}) \;=\; 0.0250, \qquad
\text{excess}(\cdot) \;=\; \text{metric}(\cdot) - 0.0250.
\label{eq:dd5-floor}
\end{equation}
A10's steady state is the best any arm in this family achieves with oracle
assignment, oracle-free training, and no reactivation stress; treating it as
the irreducible floor makes relative improvements interpretable. A $10.3\%$
improvement in raw regret (A6 vs.\ A1, post-reactivation) is a $38.1\%$
reduction in excess-over-floor---the latter is the honest measure of how much
of the \emph{addressable} regret the mechanism removes
(Section~\ref{sec:dd5-e1-axes}).

\subsection{Pre-registration discipline}
\label{sec:dd5-prereg}

Before the final 20-seed E1 run, we fixed in writing: (i) the eleven-arm
roster and all hyperparameters; (ii) the three metrics and the primary
endpoint (\texttt{post\_react}); (iii) the family of eleven pairwise
comparisons and the Holm correction over exactly that family; (iv) a
directional prediction for every comparison, including four predicted
\emph{nulls} (A2$\approx$A1, A7$\approx$A1, A4$\approx$A5,
A6$\approx$A10); and (v) the E0 structure gate for real data, with the
commitment to run E1s and report it regardless of the gate's outcome. The
predicted nulls matter as much as the predicted effects: a theory that only
predicts differences can survive almost any data, whereas predicting that
invariance alone does \emph{nothing} post-reactivation (A7$\approx$A1) is a
claim the data could easily have destroyed.

\begin{intuition}
The entire E1 design is a $2\times 2$ factorial hiding inside an eleven-arm
lineup: \{monolith, niche\} $\times$ \{ERM, INV\} is the square
(A1, A5, A7, A6), and everything else is a control that pins down \emph{why}
the square comes out the way it does. A2 and A3 show that reward-coupled
allocation does not reproduce the niche effect; A4 shows assignment quality
matters beyond mere partitioning; A8a/A8b bracket the system between
``memory without adaptation'' and ``adaptation without memory''; A9/A10 show
how close learned pinning gets to the oracle ceiling.
\end{intuition}

\begin{warning}
\textbf{What these statistics do not establish.} Every $p$-value in this
part is a statement about \emph{mean decision regret on a specified
generator or historical sample, under a specified decision layer, relative
to other arms of the same family}. None of the following is claimed, and the
reader should resist inferring any of it: (i) no claim of positive risk-adjusted
return (``alpha'') on any market---regret relative to a daily hindsight
oracle is not profit; (ii) no claim about live trading, transaction costs,
slippage, capacity, or market impact; (iii) no claim that the synthetic
generator is a calibrated model of any asset class; (iv) no claim that the
real-data nulls of Section~\ref{sec:dd5-e0e1s} would persist under other
universes, frequencies, or feature sets---they are scoped to the data
actually tested. The expected post-publication fate of any genuine
market anomaly is decay \citep{mclean2016does}; our claims are deliberately
about \emph{mechanism}, which is the only thing a synthetic-plus-null-real
evidence base can support.
\end{warning}

% ============================================================================
\section{E1: The Headline Dissociation, Number by Number}
\label{sec:dd5-e1}

\subsection{The full table}
\label{sec:dd5-e1-table}

Table~\ref{tab:dd5-e1-main} reports all eleven arms on all three metrics,
with 95\% CIs, from the archived 20-seed run ($K=2$, heterogeneity $1.0$,
hard memory, probe $15$, minimum dormancy $90$ days).

\begin{table}[h]
\centering
\caption{E1 headline results. Mean daily decision regret ($\times 10^{-2}$,
lower is better) $\pm$ 95\% CI half-width over 20 seeds. Excess = post-react
mean minus the noise floor $2.50 \times 10^{-2}$
(Eq.~\eqref{eq:dd5-floor}). Bold: best non-oracle arm per column.}
\label{tab:dd5-e1-main}
\small
\begin{tabular}{lccccc}
\toprule
Arm & \texttt{post\_react} & \texttt{overall} & \texttt{steady} &
Excess & Post/steady \\
\midrule
A1 monolith-ERM      & $3.424 \pm 0.095$ & $2.690 \pm 0.044$ & $2.459 \pm 0.058$ & $0.925$ & $1.39$ \\
A2 router            & $3.408 \pm 0.110$ & $2.695 \pm 0.052$ & $2.481 \pm 0.070$ & $0.909$ & $1.37$ \\
A3 recent-perf       & $3.776 \pm 0.096$ & $2.940 \pm 0.055$ & $2.685 \pm 0.072$ & $1.277$ & $1.41$ \\
A4 random-fixed      & $3.295 \pm 0.126$ & $2.827 \pm 0.114$ & $2.717 \pm 0.124$ & $0.795$ & $1.21$ \\
A5 RISP-ERM      & $3.214 \pm 0.070$ & $2.750 \pm 0.064$ & $2.644 \pm 0.084$ & $0.715$ & $1.22$ \\
A6 RISP-INV      & $\mathbf{3.072 \pm 0.066}$ & $\mathbf{2.602 \pm 0.059}$ & $2.507 \pm 0.074$ & $\mathbf{0.572}$ & $1.23$ \\
A7 monolith-INV      & $3.435 \pm 0.100$ & $2.663 \pm 0.044$ & $\mathbf{2.428 \pm 0.062}$ & $0.935$ & $1.41$ \\
A8a Hedge-fixed      & $4.390 \pm 0.090$ & $3.815 \pm 0.069$ & $3.704 \pm 0.096$ & $1.890$ & $1.19$ \\
A8b Hedge-learners   & $3.407 \pm 0.100$ & $2.672 \pm 0.046$ & $2.443 \pm 0.057$ & $0.907$ & $1.39$ \\
\midrule
A9 oracle-pinned-ERM & $3.210 \pm 0.077$ & $2.736 \pm 0.066$ & $2.624 \pm 0.088$ & $0.710$ & $1.22$ \\
A10 oracle-pinned-INV& $3.053 \pm 0.078$ & $2.590 \pm 0.059$ & $2.499 \pm 0.075$ & $0.553$ & $1.22$ \\
\bottomrule
\end{tabular}
\end{table}

Three structural facts are visible before any test is run. First, the
\texttt{post\_react} column separates the arms far more than the
\texttt{steady} column: retention is a property of reactivation days.
Second, the post/steady ratio splits the arms into two clean families:
everything with a reward-coupled or absent assignment (A1, A2, A3, A7, A8b)
sits at $1.37$--$1.41$, and everything with a fixed or pinned assignment
(A4, A5, A6, A8a, A9, A10) sits at $1.19$--$1.23$. The ratio is a
fingerprint of the mechanism, independent of the level. Third, the
\texttt{steady} column shows the honest cost: pinning is not free. A5's
steady regret ($2.644$) is \emph{worse} than A1's ($2.459$), because
partitioned specialists see less data per niche; A7 actually posts the best
steady value in the table ($2.428$), better even than A10. RISP wins
\texttt{overall} because reactivation days are frequent and expensive enough
that retention pays for the partitioning cost---not because it dominates
everywhere.

\subsection{The pre-registered comparisons}
\label{sec:dd5-e1-tests}

\begin{table}[h]
\centering
\caption{The eleven pre-registered comparisons on \texttt{post\_react}.
$\Delta$ = first arm minus second arm ($\times 10^{-2}$; negative favors the
first arm). Verdicts at FWER $\alpha = 0.05$ (Holm). ``Null (pred.)''
denotes comparisons pre-registered as approximate equalities.}
\label{tab:dd5-e1-tests}
\small
\begin{tabular}{llrccl}
\toprule
\# & Comparison & $\Delta$ & Welch $p$ & Holm $p$ & Verdict \\
\midrule
1 & A6 vs A1  & $-0.353$ & $9.2\times 10^{-7}$ & $1.0\times 10^{-5}$ & rejected: A6 $<$ A1 \\
2 & A6 vs A7  & $-0.363$ & $1.2\times 10^{-6}$ & $1.2\times 10^{-5}$ & rejected: A6 $<$ A7 \\
3 & A6 vs A8b & $-0.335$ & $4.5\times 10^{-6}$ & $4.1\times 10^{-5}$ & rejected: A6 $<$ A8b \\
4 & A5 vs A1  & $-0.210$ & $1.3\times 10^{-3}$ & $1.0\times 10^{-2}$ & rejected: A5 $<$ A1 \\
5 & A6 vs A5  & $-0.142$ & $6.1\times 10^{-3}$ & $4.3\times 10^{-2}$ & rejected: A6 $<$ A5 \\
6 & A5 vs A2  & $-0.194$ & $6.3\times 10^{-3}$ & $4.3\times 10^{-2}$ & rejected: A5 $<$ A2 \\
7 & A6 vs A9  & $-0.138$ & $1.1\times 10^{-2}$ & $5.7\times 10^{-2}$ & \emph{marginal} (fails Holm) \\
8 & A5 vs A4  & $-0.081$ & $0.281$ & $1.0$ & null (pred.) \\
9 & A6 vs A10 & $+0.019$ & $0.719$ & $1.0$ & null (pred.) \\
10 & A2 vs A1 & $-0.016$ & $0.827$ & $1.0$ & null (pred.) \\
11 & A7 vs A1 & $+0.011$ & $0.882$ & $1.0$ & null (pred.) \\
\bottomrule
\end{tabular}
\end{table}

We now walk through each pre-registered prediction.

\paragraph{P1: A6 $<$ A5 $<$ A1 (both axes contribute).}
Confirmed at every link. A5 beats A1 by $0.210 \times 10^{-2}$
(Holm $p = 0.010$): pinned niches alone, with plain ERM inside them, retain
dormant-regime competence. A6 beats A5 by a further $0.142 \times 10^{-2}$
(Holm $p = 0.043$): the invariance objective adds value \emph{on top of}
retention. The chain A6 $<$ A1 is the largest effect in the family
(Holm $p = 1.0 \times 10^{-5}$).

\paragraph{P2: A7 $\approx$ A1 post-reactivation (the dissociation).}
Confirmed as a null: $\Delta = +0.011 \times 10^{-2}$, raw $p = 0.882$. The
invariance objective, trained in a monolith, does \emph{nothing} for
post-reactivation regret---indeed A7's point estimate is fractionally worse
than A1's. This is the single most theory-laden prediction in the family:
invariance can only protect parameters that still exist, and a monolith
overwrites its dormant-regime parameters regardless of how they were
trained. Note the contrast with \texttt{steady}, where A7 is the best arm in
the table: invariance helps the monolith \emph{between} reactivations and is
helpless \emph{at} them.

\paragraph{P3: A2 $\approx$ A1 (a reward-driven router is a monolith in
disguise).}
Confirmed as a null: $\Delta = -0.016 \times 10^{-2}$, raw $p = 0.827$. The
learned gate receives no gradient from a dormant regime and reallocates its
experts' capacity to live regimes; on reactivation it relearns like the
monolith. The companion test A5 $<$ A2 (Holm $p = 0.043$) confirms the
positive side: dormancy-pinning beats the router by about as much as it
beats the monolith.

\paragraph{P4: A3 is the worst learner.}
Confirmed. A3's post-reactivation regret ($3.776 \times 10^{-2}$) is the
worst of any learning arm---only the frozen A8a is worse---and uniquely
among the learners its damage is \emph{not} confined to probe windows: its
\texttt{overall} ($2.940$) and \texttt{steady} ($2.685$) are also the worst
in the learner family. The failure mode is continuous erosion: the
recent-performance allocator drives capital toward whatever just worked,
and the capital floor (the minimum allocation that keeps a specialist
training at all) is exactly where a temporarily-unlucky specialist's data
diet collapses. Every regime transition triggers a starve-and-relearn cycle,
so A3 pays the relearning tax perpetually, not just after long dormancies.

\paragraph{P5: A8a retains but cannot adapt; A8b adapts but cannot retain.}
Confirmed on both jaws of the bracket. A8a (Hedge over frozen experts) has
the \emph{lowest} post/steady ratio in the table ($1.19$): its relative
degradation at reactivation is minimal, because nothing it could forget was
ever learned online---this is retention in its purest form. But its absolute
level is catastrophic ($\texttt{steady} = 3.704$, $51\%$ above A1), because
frozen experts cannot track within-regime drift at all. A8b (Hedge over
learning experts) is the mirror image: its level tracks the monolith almost
exactly ($\texttt{post\_react}$ $3.407$ vs $3.424$; \texttt{steady} $2.443$
vs $2.459$; ratio $1.39$ vs $1.39$), because exponential weights over
experts that all share the reward stream do not protect any expert's
parameters from being retrained. The sleeping-experts story, in other words,
is not about the \emph{weighting} scheme; it is about whether the experts'
parameters are insulated from reward-coupled reuse. A6 beats A8b by
$0.335\times 10^{-2}$ (Holm $p = 4.1\times 10^{-5}$).

\paragraph{P6: A4 $\approx$ A5 in mean, with higher variance.}
\label{sec:dd5-e1-a4}
Confirmed. Random fixed niches land at $3.295 \times 10^{-2}$ against A5's
$3.214$ (raw $p = 0.281$, Holm $1.0$): a random reward-independent
assignment captures most of the retention effect, exactly as the theory
says it should---retention comes from reward-independence, not from
assignment quality. But A4's between-seed SD is $0.0027$ against A5's
$0.0015$ ($1.8\times$), and the signature is even starker on
\texttt{steady} (CI $0.124$ vs $0.084$) and \texttt{overall} (CI $0.114$ vs
$0.064$). Inspection of A4's per-seed values shows why: most seeds are fine,
but the seeds in which the random assignment leaves a regime with no
well-matched specialist (coverage gaps) blow up the mean (per-seed
\texttt{overall} values up to $0.0341$ against a median near $0.0270$).
Learned pinning buys insurance against the coverage tail, not a better
average draw.

\paragraph{P7: A6 $\approx$ A10 (the skyline is reached).}
Confirmed as a null: $\Delta = +0.019 \times 10^{-2}$, raw $p = 0.719$. The
full system is statistically indistinguishable from the same system with an
oracle-given assignment. Learned dormancy-pinning recovers essentially all
of the oracle assignment's value; the corresponding ERM-side check is even
cleaner (A5 $3.214$ vs A9 $3.210$, not in the test family but visibly
identical).

\paragraph{P8 (exploratory edge): A6 vs A9.}
A6's point estimate beats the \emph{oracle-assigned ERM} arm by
$0.138\times 10^{-2}$---learned assignment plus invariance edging out
perfect assignment without it. Raw $p = 0.011$, but the Holm-adjusted value
is $0.057$, which fails the $0.05$ FWER threshold. We report this exactly
as it is: suggestive, consistent with the rest of the pattern, and not
established at the pre-registered standard. Nothing downstream depends on
it.

\begin{keypoint}
Seven of seven predicted differences went the predicted direction; six of
seven survived Holm at $\alpha=0.05$ (the seventh, A6 $<$ A9, is marginal
at $0.057$). All four predicted nulls came back null with raw
$p \ge 0.28$. No comparison in the pre-registered family contradicted a
prediction.
\end{keypoint}

\subsection{The axes arithmetic}
\label{sec:dd5-e1-axes}

Decomposing the headline effect along the factorial square:
\begin{align}
\text{retention axis (A1} \to \text{A5):}\quad
& \frac{3.424 - 3.214}{3.424} = 6.1\%, \label{eq:dd5-axis-ret}\\[2pt]
\text{invariance axis (A5} \to \text{A6):}\quad
& \frac{3.214 - 3.072}{3.214} = 4.4\%, \label{eq:dd5-axis-inv}\\[2pt]
\text{joint (A1} \to \text{A6):}\quad
& \frac{3.424 - 3.072}{3.424} = 10.3\%. \label{eq:dd5-axis-joint}
\end{align}
In excess-over-floor terms (Eq.~\eqref{eq:dd5-floor}) the joint effect is
\begin{equation}
1 - \frac{3.072 - 2.500}{3.424 - 2.500}
= 1 - \frac{0.572}{0.925} = 38.1\%:
\label{eq:dd5-excess}
\end{equation}
the full system removes a bit over a third of the addressable
post-reactivation regret, and lands within $3\%$ of the oracle skyline's
excess ($0.553$ vs $0.572$).

\subsection{The interaction index: super-additivity}
\label{sec:dd5-e1-interaction}

Define, per seed, the factorial interaction
\begin{equation}
I \;=\; (\text{A1} - \text{A5}) + (\text{A1} - \text{A7})
- (\text{A1} - \text{A6})
\;=\; \text{A1} - \text{A5} - \text{A7} + \text{A6}.
\label{eq:dd5-interaction}
\end{equation}
If the two mechanisms contributed additively, $I$ would be zero: the joint
gain would equal the sum of the marginal gains. Computed per seed on the
paired design (here pairing is legitimate and free, since
Eq.~\eqref{eq:dd5-interaction} is a within-seed contrast),
\[
I \;=\; -0.0015 \;\pm\; 0.0005 \quad (\text{mean} \pm \text{95\% CI},\;
n = 20),
\]
which excludes zero decisively. The sign is the interesting part: $I < 0$
means the joint gain $(\text{A1}-\text{A6}) = 0.0035$ \emph{exceeds} the sum
of the marginals $(\text{A1}-\text{A5}) + (\text{A1}-\text{A7}) =
0.0021 + (-0.0001) = 0.0020$. The interaction is super-additive, and the
arithmetic shows why in one line: the invariance marginal in the monolith is
\emph{zero} (A1 $-$ A7 $= -0.0001$), yet inside niches invariance
contributes $0.0014$ (A5 $-$ A6). Invariance is not an independent additive
bonus; it is a multiplier on retained memories. A specialist that still
holds its dormant regime's parameters can hold \emph{invariant} parameters,
which transfer across the within-regime drift accumulated during dormancy;
a monolith has nothing for the invariance penalty to protect at
reactivation time. This is the empirical signature of the
retention$\times$robustness interaction derived in
Part~III, and it echoes the
invariance-for-decision-making rationale of \citet{krueger2021rex} and
\citet{yan2026invpnco}: invariance pays where the deployment distribution
differs from the training distribution---which, for a retained dormant
niche, is exactly the reactivation moment.

% ============================================================================
\section{E0 and E1s: The Real-Data Nulls}
\label{sec:dd5-e0e1s}

A mechanism paper earns the right to a synthetic headline only by trying,
honestly, to find the mechanism's preconditions in real data, and reporting
what it finds. We found that the precondition---persistent, re-activating,
\emph{decision-relevant} conditional structure---is absent in the real
universes we tested. This section documents that null in full.

\subsection{E0: the structure diagnostic}
\label{sec:dd5-e0}

\paragraph{Protocol.} E0 asks the gating question without running any
learning arm at all: \emph{if an oracle could condition the decision policy
on regime labels, would that beat not conditioning?} The components:

\begin{enumerate}
\item \textbf{Universes.} Five crypto pairs (daily bars) and three
commodities (daily). Evaluation spans $562$ days (crypto) and $1{,}270$
days (commodities).
\item \textbf{Causal labelers.} Two regime labelers, L1 (finer) and L2
(coarser), computed from \emph{trailing} data only, so labels are known
before the day's decision and cannot leak the target. The realized regime
distributions are imbalanced, especially under L2 on crypto
(counts $517/517/31/59$), which we note as a power limitation for the rare
regimes.
\item \textbf{Decision metric.} The same top-$K$ portfolio regret as E1
(Section~\ref{sec:dd5-metrics})---a decision-space test, not a
prediction-space test.
\item \textbf{Oracle conditioning.} Fit the best per-regime decision policy
with hindsight model selection per regime (an upper bound on what any
regime-aware learner could extract), and compare its mean regret to
(a) the same machinery run on \emph{block-shuffled} labels (preserving label
marginals and block-length statistics, destroying the label--day
alignment), and (b) a pooled, unconditioned policy.
\item \textbf{Statistic.} Gap $\% = (\text{shuffled} - \text{real}) /
\text{shuffled} \times 100$ (positive iff real labels help), and
$z = (\text{shuffled} - \text{real}) / \mathrm{sd}(\text{shuffled})$ over
the shuffle distribution.
\end{enumerate}

\begin{table}[h]
\centering
\caption{E0 structure diagnostic on real data. Mean daily decision regret
($\times 10^{-2}$) for oracle regime-conditioning on real labels, on
block-shuffled labels (mean $\pm$ SD of the shuffle distribution), and for
the pooled unconditioned policy. Positive gap = real labels help. No cell
shows exploitable structure.}
\label{tab:dd5-e0}
\small
\begin{tabular}{llccrrr}
\toprule
Domain & Labeler & Real & Shuffled & Pooled & Gap (\%) & $z$ \\
\midrule
Crypto ($n{=}562$ d) & L1 & $1.220$ & $1.202 \pm 0.027$ & $1.195$ & $-1.47$ & $-0.66$ \\
Crypto ($n{=}562$ d) & L2 & $1.212$ & $1.204 \pm 0.047$ & $1.195$ & $-0.66$ & $-0.17$ \\
Commodities ($n{=}1270$ d) & L1 & $0.834$ & $0.832 \pm 0.015$ & $0.816$ & $-0.25$ & $-0.14$ \\
Commodities ($n{=}1270$ d) & L2 & $0.836$ & $0.838 \pm 0.009$ & $0.816$ & $+0.16$ & $+0.16$ \\
\bottomrule
\end{tabular}
\end{table}

\paragraph{Reading.} Across both domains and both labelers, the real-label
oracle never beats the shuffled-label oracle by more than $0.16\%$, all
$|z| < 0.7$, and three of four point estimates are \emph{negative}---real
labels do marginally worse than shuffled ones. Most damning, the
\emph{pooled} policy is the best cell in every row: even hindsight-optimal
regime conditioning loses to not conditioning at all, because splitting the
sample by regime costs more in estimation variance than the (apparently
nonexistent) regime structure returns. This is a null with teeth: it bounds
not just RISP but \emph{any} regime-conditioned method in this family,
on this data, at this frequency, with these features.

\paragraph{Honest scoping.} The null is scoped: five crypto pairs and three
commodities, daily bars, one feature set, two specific (if reasonable)
causal labelers, and a top-$K$ decision layer. Intraday frequencies, larger
cross-sections, alternative labelers (e.g.\ macro-indexed), or other
decision layers could in principle revive the gate, and the protocol is
designed to be rerun cheaply on any of them.

\paragraph{Relation to the GAUSE null.} The predecessor project found a
closely related null in \emph{prediction} space: regime-aware prediction
showed no real-vs-shuffled gap on financial domains \citep{li2026gause}.
That did not settle our question, and pretending otherwise would have been
sloppy: the predict-then-optimize literature documents real cases where
prediction-space and decision-space rankings dissociate
\citep{elmachtoub2022smart,yan2026invpnco}, so a decision-metric retest was
a genuine open question, not a formality. E0 answers it: the dissociation
does not rescue financial regime structure here. The prediction-space null
and the decision-space null now agree.

\subsection{E1s: the stitched-crypto experiment}
\label{sec:dd5-e1s}

E0 was the pre-registered gate for real-data claims, and it failed.
Pre-registration committed us to running the full arm lineup on real data
anyway and predicting \emph{null separation}---a falsifiable prediction in
its own right (if arms had separated despite the failed gate, the gate
would be broken). E1s stitches real crypto regime blocks (block counts
$15/13/7/9$ across the four labeled regimes) into schedules with genuine
dormancy (minimum $80$ days, probe $15$ days), $20$ seeds of stitching
randomization, ten arms (A10 requires generator access and is undefined on
real data).

\begin{table}[h]
\centering
\caption{E1s stitched real crypto. Mean daily decision regret
($\times 10^{-2}$) $\pm$ 95\% CI, 20 seeds. All pairwise comparisons in the
pre-registered family: raw Welch $p > 0.21$, all Holm $p = 1.0$.}
\label{tab:dd5-e1s}
\small
\begin{tabular}{lcc}
\toprule
Arm & \texttt{post\_react} & \texttt{overall} \\
\midrule
A1 monolith-ERM      & $1.396 \pm 0.051$ & $1.285 \pm 0.017$ \\
A2 router            & $1.403 \pm 0.050$ & $1.296 \pm 0.020$ \\
A3 recent-perf       & $1.398 \pm 0.075$ & $1.299 \pm 0.025$ \\
A4 random-fixed      & $1.373 \pm 0.057$ & $1.273 \pm 0.027$ \\
A5 RISP-ERM      & $1.381 \pm 0.056$ & $1.276 \pm 0.025$ \\
A6 RISP-INV      & $1.344 \pm 0.064$ & $1.263 \pm 0.025$ \\
A7 monolith-INV      & $1.410 \pm 0.057$ & $1.290 \pm 0.018$ \\
A8a Hedge-fixed      & $1.325 \pm 0.053$ & $1.239 \pm 0.023$ \\
A8b Hedge-learners   & $1.400 \pm 0.059$ & $1.285 \pm 0.017$ \\
A9 oracle-pinned-ERM & $1.342 \pm 0.052$ & $1.268 \pm 0.027$ \\
\bottomrule
\end{tabular}
\end{table}

All ten arms land in a band of $1.32$--$1.41 \times 10^{-2}$
post-reactivation; the six pre-registered comparisons return raw
$p \in \{0.22, 0.23, 0.40, 0.57, 0.71, 0.95\}$ and Holm $p = 1.0$
throughout. Two details confirm that this is structure-absence rather than
mere noise. First, the numerically best post-reactivation arm is
\emph{A8a}---the frozen-expert arm that cannot adapt to anything---which is
exactly what one expects when there is nothing regime-specific to retain or
relearn: not adapting is as good as adapting. Second, A9 (oracle labels!)
is statistically indistinguishable from A6 with learned labels
($p = 0.95$): even perfect knowledge of the regime labels confers nothing,
because the labels do not predict decision-relevant structure (which is
what E0 measured directly).

\begin{keypoint}
\textbf{The protocol is falsifiable, and falsifies.} The identical
machinery---same arms, same metrics, same test family, same Holm
correction---produces Holm $p \approx 10^{-5}$ separations where the
generator contains dormant-but-persistent structure (E1), and a uniform
$p > 0.2$ flatline where a pre-registered diagnostic says the structure is
absent (E0 $\to$ E1s). A pipeline that can only return positives is
untrustworthy on synthetic data too; the E1s flatline is what makes the E1
separation credible.
\end{keypoint}

% ============================================================================
\section{E2--E4: Boundary Sweeps}
\label{sec:dd5-sweeps}

The headline lives at one point in configuration space ($K=2$, dormancy
$\ge 90$, heterogeneity $1.0$). E2--E4 sweep the three axes the theory says
should govern the effect: capacity pressure, dormancy length, and
cross-regime heterogeneity.

\subsection{E2: capacity}
\label{sec:dd5-e2}

The retention argument requires capacity pressure: a monolith forgets a
dormant regime only if serving the live regimes forces reuse of its
capacity. As $K \to R$ the pressure vanishes and the retention gap should
close. E2 sweeps $K \in \{1,2,3,4\}$ ($R = 4$) under two memory models:
\emph{hard} (slot eviction) and \emph{soft} (graded interference decay).

\begin{table}[h]
\centering
\caption{E2 capacity sweep, \texttt{post\_react} ($\times 10^{-2}$) $\pm$
95\% CI, hard memory. The pinned arms (A5, A6, A9) are exactly constant
across $K$: each specialist serves only its own niches, so the capacity
constraint never binds for them. Rightmost column: retention gap
A1 $-$ A5 ($\times 10^{-2}$).}
\label{tab:dd5-e2-hard}
\small
\begin{tabular}{cccccccc}
\toprule
$K$ & A1 & A2 & A4 & A5 & A6 & A9 & A1$-$A5 \\
\midrule
1 & $3.427 \pm 0.102$ & $3.432 \pm 0.100$ & $3.667 \pm 0.163$ & $3.214 \pm 0.070$ & $3.072 \pm 0.066$ & $3.210 \pm 0.077$ & $0.213$ \\
2 & $3.424 \pm 0.095$ & $3.408 \pm 0.110$ & $3.295 \pm 0.126$ & $3.214 \pm 0.070$ & $3.072 \pm 0.066$ & $3.210 \pm 0.077$ & $0.210$ \\
3 & $3.303 \pm 0.088$ & $3.276 \pm 0.064$ & $3.202 \pm 0.072$ & $3.214 \pm 0.070$ & $3.072 \pm 0.066$ & $3.210 \pm 0.077$ & $0.089$ \\
4 & $3.221 \pm 0.082$ & $3.228 \pm 0.068$ & $3.206 \pm 0.078$ & $3.214 \pm 0.070$ & $3.072 \pm 0.066$ & $3.210 \pm 0.077$ & $0.007$ \\
\bottomrule
\end{tabular}
\end{table}

\begin{table}[h]
\centering
\caption{E2 capacity sweep, \texttt{post\_react} ($\times 10^{-2}$) $\pm$
95\% CI, soft memory (no eviction; graded interference with
overflow-scaled decay). Rightmost column: retention gap A1 $-$ A5.}
\label{tab:dd5-e2-soft}
\small
\begin{tabular}{cccccccc}
\toprule
$K$ & A1 & A2 & A4 & A5 & A6 & A9 & A1$-$A5 \\
\midrule
1 & $3.276 \pm 0.079$ & $3.285 \pm 0.067$ & $3.667 \pm 0.163$ & $3.214 \pm 0.070$ & $3.072 \pm 0.066$ & $3.210 \pm 0.077$ & $0.062$ \\
2 & $3.277 \pm 0.079$ & $3.254 \pm 0.077$ & $3.318 \pm 0.117$ & $3.214 \pm 0.070$ & $3.072 \pm 0.066$ & $3.210 \pm 0.077$ & $0.063$ \\
3 & $3.279 \pm 0.079$ & $3.265 \pm 0.056$ & $3.239 \pm 0.071$ & $3.214 \pm 0.070$ & $3.072 \pm 0.066$ & $3.210 \pm 0.077$ & $0.065$ \\
4 & $3.277 \pm 0.081$ & $3.255 \pm 0.064$ & $3.273 \pm 0.079$ & $3.214 \pm 0.070$ & $3.072 \pm 0.066$ & $3.210 \pm 0.077$ & $0.063$ \\
\bottomrule
\end{tabular}
\end{table}

Four observations.

\begin{enumerate}
\item \textbf{Reward-independent arms are exactly flat.} A5, A6, and A9 post
identical values at every $K$ in both memory models. This is not a rounding
artifact but a structural fact: a pinned specialist serves only its assigned
niches, so its memory never overflows and the capacity model is inert for
it. Flatness is the cleanest possible demonstration that the pinned arms'
advantage is \emph{not} extra effective capacity---they use the same budget
and simply never churn it.
\item \textbf{Gap closure at $K = R$ under hard memory.} The retention gap
A1 $-$ A5 falls from $0.213$ at $K=1$ to $0.007$ at $K=4$---zero to within
a tenth of a CI. With enough slots, the monolith never evicts, never
forgets, and pinning buys nothing. The invariance edge survives, however:
A6 ($3.072$) still beats everything at $K=4$, confirming the axes are
distinct (capacity governs retention, not robustness).
\item \textbf{Non-closure under soft memory---at any capacity.} With no
eviction, capacity slots stop binding and the reward-driven arms go
\emph{flat} in $K$ at $\approx 3.28$: gentler than hard eviction at tight
capacity ($3.43$ at $K \le 2$) but never closing the gap to the pinned
arms ($\approx 0.063$ above A5, $0.21$ above A6, at every $K$), because
graded interference from foreign updates erodes dormant heads regardless
of how many slots exist. The hard model predicts closure at $K = R$ and
closes; the soft model predicts a capacity-independent residue and shows
exactly that.
\item \textbf{A4's coverage problem is worst at $K=1$.} Random fixed niches
post $3.667 \pm 0.163$ at $K=1$---the worst learner cell in the
sweep---because with one slot per specialist a single unlucky assignment
leaves a regime entirely uncovered. The cell's CI ($0.163$) is the widest in
the sweep, the variance signature of Section~\ref{sec:dd5-e1-a4} amplified.
\end{enumerate}

\subsection{E3: dormancy length}
\label{sec:dd5-e3}

Retention should matter more the longer a regime sleeps: the monolith's
relearning debt accumulates with dormancy, while a pinned specialist's
retained parameters do not decay (hard memory). E3 sweeps the minimum
dormancy threshold $D \in \{21, 63, 126, 252, 504\}$ days.

\begin{table}[h]
\centering
\caption{E3 dormancy sweep ($\times 10^{-2}$, $\pm$ 95\% CI). Each $D$ is a
\emph{different schedule generator} (shorter dormancies mean more frequent
transitions and different steady-state mixes), so levels are comparable
\emph{within} a row block only, never across $D$. Rightmost columns:
within-$D$ retention$+$invariance gap A1 $-$ A6, absolute and relative.}
\label{tab:dd5-e3}
\small
\begin{tabular}{clccr@{\quad}r}
\toprule
$D$ & Arm & \texttt{post\_react} & \texttt{steady} &
\multicolumn{2}{c}{A1 $-$ A6 (post)} \\
\midrule
\multirow{5}{*}{21}
& A1 & $5.226 \pm 0.148$ & $3.225 \pm 0.225$ & & \\
& A2 & $5.241 \pm 0.151$ & $3.305 \pm 0.213$ & & \\
& A3 & $5.475 \pm 0.169$ & $3.504 \pm 0.263$ & $0.152$ & $2.9\%$ \\
& A5 & $5.198 \pm 0.167$ & $3.416 \pm 0.220$ & & \\
& A6 & $5.074 \pm 0.142$ & $3.293 \pm 0.222$ & & \\
\midrule
\multirow{5}{*}{63}
& A1 & $5.024 \pm 0.202$ & $2.719 \pm 0.214$ & & \\
& A2 & $5.028 \pm 0.208$ & $2.785 \pm 0.226$ & & \\
& A3 & $5.362 \pm 0.173$ & $2.973 \pm 0.218$ & $0.375$ & $7.5\%$ \\
& A5 & $4.818 \pm 0.198$ & $2.876 \pm 0.208$ & & \\
& A6 & $4.648 \pm 0.193$ & $2.776 \pm 0.209$ & & \\
\midrule
\multirow{5}{*}{126}
& A1 & $4.975 \pm 0.160$ & $2.550 \pm 0.157$ & & \\
& A2 & $4.960 \pm 0.150$ & $2.598 \pm 0.156$ & & \\
& A3 & $5.228 \pm 0.147$ & $2.747 \pm 0.159$ & $0.401$ & $8.1\%$ \\
& A5 & $4.787 \pm 0.185$ & $2.762 \pm 0.141$ & & \\
& A6 & $4.574 \pm 0.181$ & $2.649 \pm 0.130$ & & \\
\midrule
\multirow{5}{*}{252}
& A1 & $4.336 \pm 0.098$ & $2.633 \pm 0.100$ & & \\
& A2 & $4.344 \pm 0.103$ & $2.657 \pm 0.096$ & & \\
& A3 & $4.682 \pm 0.103$ & $2.811 \pm 0.103$ & $0.385$ & $8.9\%$ \\
& A5 & $4.095 \pm 0.144$ & $2.777 \pm 0.099$ & & \\
& A6 & $3.951 \pm 0.103$ & $2.652 \pm 0.092$ & & \\
\midrule
\multirow{5}{*}{504}
& A1 & $3.668 \pm 0.059$ & $2.530 \pm 0.080$ & & \\
& A2 & $3.646 \pm 0.059$ & $2.565 \pm 0.080$ & & \\
& A3 & $3.968 \pm 0.049$ & $2.713 \pm 0.079$ & $0.408$ & $11.1\%$ \\
& A5 & $3.454 \pm 0.076$ & $2.668 \pm 0.077$ & & \\
& A6 & $3.260 \pm 0.061$ & $2.533 \pm 0.077$ & & \\
\bottomrule
\end{tabular}
\end{table}

\begin{warning}
Do not compare levels across $D$ blocks. A $D = 21$ generator produces a
chaotic schedule in which \emph{every} arm is perpetually post-transition
(note the elevated steady values, $3.2$--$3.5$, against $2.5$--$2.7$ at
$D = 504$); a $D = 504$ generator produces long calm stretches. The only
licensed reading is the \emph{within-$D$ gap between arms}, which holds the
schedule fixed.
\end{warning}

The within-$D$ gap A1 $-$ A6 grows steeply from $0.152$ at $D = 21$ to
$0.375$ at $D = 63$, then saturates in absolute terms
($0.375 \to 0.401 \to 0.385 \to 0.408$), while the relative gap continues
to drift up ($7.5\% \to 11.1\%$) because the baseline level falls as
schedules calm down. The shape is exactly the theory's: at $D = 21$ the
monolith's parameters have barely drifted from the dormant regime, so there
is little to retain; by $D \approx 63$--$126$ (one to two quarters of
dormancy) the monolith's relearning debt is fully formed and further
dormancy adds little, because the debt is bounded by ``relearn from
scratch.'' The ordering A6 $<$ A5 $<$ A1 $\approx$ A2, and A3 worst, holds
at every $D$ without exception.

\subsection{E4: heterogeneity}
\label{sec:dd5-e4}

The invariance axis should engage only insofar as regimes differ: at
heterogeneity $0$ all regimes share one mechanism and there is nothing to be
invariant \emph{across}. E4 sweeps the cross-regime heterogeneity scale
$h \in \{0, 0.5, 1.0, 1.5, 2.0\}$ for the pinned arms.

\begin{table}[h]
\centering
\caption{E4 heterogeneity sweep, \texttt{post\_react} ($\times 10^{-2}$)
$\pm$ 95\% CI. Rightmost column: ERM$-$INV gap (A5 $-$ A6).}
\label{tab:dd5-e4}
\small
\begin{tabular}{ccccc}
\toprule
$h$ & A5 RISP-ERM & A6 RISP-INV & A9 oracle-pinned-ERM & A5$-$A6 \\
\midrule
0.0 & $2.998 \pm 0.080$ & $2.900 \pm 0.079$ & $2.997 \pm 0.081$ & $0.097$ \\
0.5 & $3.073 \pm 0.081$ & $2.948 \pm 0.080$ & $3.067 \pm 0.073$ & $0.125$ \\
1.0 & $3.214 \pm 0.070$ & $3.072 \pm 0.066$ & $3.210 \pm 0.077$ & $0.142$ \\
1.5 & $3.459 \pm 0.087$ & $3.249 \pm 0.095$ & $3.418 \pm 0.094$ & $0.210$ \\
2.0 & $3.749 \pm 0.138$ & $3.511 \pm 0.124$ & $3.679 \pm 0.165$ & $0.238$ \\
\bottomrule
\end{tabular}
\end{table}

Three readings.

\begin{enumerate}
\item \textbf{Both arms degrade with heterogeneity, ERM faster.} Fitting
endpoint slopes, A5 degrades at $0.376 \times 10^{-2}$ per unit $h$ against
A6's $0.306$---the ERM arm is roughly $1.2\times$ more sensitive to
heterogeneity than the INV arm, and the oracle-pinned ERM arm A9 sits in
between ($0.341$). The cleaner statement is in the gap itself: the
ERM$-$INV gap grows \emph{monotonically} and roughly $2.4$-fold across the
sweep, from $0.0010$ at $h = 0$ to $0.0024$ at $h = 2$
(in raw units, $0.00097 \to 0.00238$). The invariance objective buys more
the more the regimes differ---which is what an invariance objective is for.
\item \textbf{The $h = 0$ residual gap.} At zero heterogeneity the theory's
first-order prediction is gap $= 0$, yet A6 retains a small, CI-overlapping
but consistently signed edge ($0.097 \times 10^{-2}$). The likely mechanism
is mundane and worth naming: with $d = 20$ features of which only a subset
is causal, the cross-environment variance penalty also suppresses reliance
on \emph{spuriously correlated distractor features within} a niche---it
acts as a targeted regularizer even when there is no cross-regime
heterogeneity to be invariant to. We flag this as a confound on the pure
invariance interpretation at low $h$: part of the A5$-$A6 gap at every $h$
is plain regularization, and only the \emph{growth} of the gap with $h$
(the $0.0010 \to 0.0024$ monotone trend) is unambiguously attributable to
invariance proper.
\item \textbf{A5 tracks A9 everywhere.} Learned pinning matches
oracle pinning under ERM at every heterogeneity ($\le 0.041 \times 10^{-2}$
apart, well inside CIs), extending the P7 skyline result along the whole
sweep.
\end{enumerate}

% ============================================================================
\section{E5--E6: Ablations and the Inversion Audit}
\label{sec:dd5-e5e6}

\subsection{E5: ablations of the full system}
\label{sec:dd5-e5}

E5 ablates A6's four design choices one at a time, holding everything else
at the headline configuration; the response variable is
\texttt{post\_react} ($\times 10^{-2}$, 20 seeds).

\paragraph{The invariance weight $\beta$.}

\begin{table}[h]
\centering
\caption{E5 $\beta$ sweep (A6 variant, \texttt{post\_react},
$\times 10^{-2}$). Reference points: A5 retained-ERM $= 3.214 \pm 0.070$;
headline A6 ($\beta = 1$) $= 3.072 \pm 0.066$.}
\label{tab:dd5-e5-beta}
\small
\begin{tabular}{cccccc}
\toprule
$\beta$ & $0$ & $0.25$ & $1$ & $4$ & $16$ \\
\midrule
\texttt{post\_react} & $3.555 \pm 0.107$ & $3.090 \pm 0.079$ &
$3.072 \pm 0.066$ & $3.170 \pm 0.076$ & $3.522 \pm 0.065$ \\
\bottomrule
\end{tabular}
\end{table}

The response is a clean U-shape with minimum at $\beta = 1$ and a broad
plateau over $\beta \in [0.25, 4]$ (all within $\sim 3\%$ of the minimum,
CIs overlapping). Both extremes collapse, and the collapses are
informative:
\begin{itemize}
\item At $\beta = 0$ the within-niche objective degenerates to the
cross-environment variance term alone, and performance ($3.555$) falls not
merely to the retained-ERM level but \emph{below} it
(A5 $= 3.214$; the variance-only arm is $10.6\%$ worse). This is precisely
the degeneracy the theory predicts (and that the V-REx analysis warns of
\citep{krueger2021rex}): a variance-only criterion is minimized by any
predictor that \emph{equalizes} risk across environments, including
uniformly bad constant-like predictors. Equal risks are not low risks. The
empirical collapse below the ERM reference is the theory's degenerate
solution showing up on schedule.
\item At $\beta = 16$ the risk term swamps the variance term and
performance ($3.522$) gives back the entire invariance gain and more,
landing $14.6\%$ above the $\beta = 1$ optimum and visibly worse than plain
retained ERM. We read the overshoot beyond A5 cautiously (the heavily
re-weighted objective also rescales effective step sizes), but the
qualitative point stands: the benefit requires \emph{both} terms in
balance, and a factor-of-16 window around the default ($0.25$--$4$) is all
flat---the choice $\beta = 1$ is not a tuned knife-edge.
\end{itemize}

\paragraph{The environment window $W_c$.}

\begin{table}[h]
\centering
\caption{E5 $W_c$ sweep (\texttt{post\_react}, $\times 10^{-2}$).}
\label{tab:dd5-e5-wc}
\small
\begin{tabular}{ccccc}
\toprule
$W_c$ & $1$ & $5$ & $20$ & $60$ \\
\midrule
\texttt{post\_react} & $3.033 \pm 0.076$ & $3.057 \pm 0.069$ &
$3.072 \pm 0.066$ & $3.090 \pm 0.070$ \\
\bottomrule
\end{tabular}
\end{table}

We had entertained a deflationary explanation of the invariance gain: that
chunking each niche's stream into $W_c$-day pseudo-environments effectively
\emph{batches} the updates, and the gain is a batching/smoothing artifact
rather than an invariance effect. The sweep refutes this cleanly, and we
report it as the refutation it is: across a $60\times$ range of $W_c$ the
response is flat to within $0.057 \times 10^{-2}$---less than one CI
half-width---and the numerically best cell is $W_c = 1$, i.e.\ \emph{no}
batching at all. If batching were the active ingredient, $W_c = 1$ would be
the worst cell, not the best. The pseudo-environment construction matters
for defining the variance penalty, not for any smoothing side effect.

\paragraph{The assignment-smoothing parameter $\lambda$.}
With $\lambda$ off ($0$): $3.141 \pm 0.082$; on ($0.3$): $3.072 \pm 0.066$.
The unpaired CIs overlap, so by the conservative standard of
Section~\ref{sec:dd5-paired} the effect is not individually established.
The paired within-seed contrast, which is legitimate here, gives
$+0.069 \pm 0.060$ ($\times 10^{-2}$, mean $\pm$ 95\% CI)---a $2.2\%$
improvement whose interval excludes zero, barely. Honest summary:
$\lambda$ helps a little, with weak evidence; nothing in the headline
depends on it.

\paragraph{Pinning schedule: pinned vs.\ never.}
Removing dormancy-pinning entirely (assignments stay perpetually plastic)
degrades \texttt{post\_react} from $3.072 \pm 0.066$ to $3.197 \pm 0.092$;
the paired contrast is $+0.126 \pm 0.059$ ($\times 10^{-2}$), a $3.9\%$
loss that is clearly resolved. Never-pinning gives back most of the
invariance-axis gain (the never-pin arm sits near the A5/A9 retained-ERM
level of $3.21$), which is consistent with the interaction story of
Section~\ref{sec:dd5-e1-interaction}: a perpetually plastic assignment lets
reactivation-time reassignment churn destroy exactly the retained-invariant
parameters the mechanism exists to protect.

\subsection{E6: the SNR audit and the inversion}
\label{sec:dd5-e6}

\paragraph{Development history, stated plainly.} This experiment exists
because our first prototype failed. The initial generator was, we later
understood, calibrated to an effectively high signal-to-noise ratio, and
the arm lineup showed no positive separation---RISP did not beat the
monolith. The tempting path was to tune the generator until the predicted
effect appeared and report that point. Instead we made SNR an explicit
axis, swept it, and report the whole curve, including the region where our
method \emph{loses}. The audit below is the honest version of what could
otherwise have been a silent hyperparameter search.

\begin{table}[h]
\centering
\caption{E6 SNR audit ($\times 10^{-2}$, $\pm$ 95\% CI, 20 seeds per cell).
$\rho$ = post/steady relearning-penalty ratio. ``A6 adv.'' =
$(\text{A1}-\text{A6})/\text{A1}$ on \texttt{post\_react}; positive favors
RISP.}
\label{tab:dd5-e6}
\scriptsize
\begin{tabular}{clccc r}
\toprule
SNR & Arm & \texttt{post\_react} & \texttt{steady} & $\rho$ & A6/A9 adv. \\
\midrule
\multirow{3}{*}{$0.5\times$}
& A1  & $3.778 \pm 0.077$ & $2.963 \pm 0.057$ & $1.28$ & \\
& A6  & $3.466 \pm 0.073$ & $2.890 \pm 0.066$ & $1.20$ & $+8.3\%$ \\
& A9  & $3.584 \pm 0.068$ & $3.000 \pm 0.068$ & $1.19$ & $+5.1\%$ \\
\midrule
\multirow{3}{*}{$1\times$}
& A1  & $3.424 \pm 0.095$ & $2.459 \pm 0.058$ & $1.39$ & \\
& A6  & $3.072 \pm 0.066$ & $2.507 \pm 0.074$ & $1.23$ & $+10.3\%$ \\
& A9  & $3.210 \pm 0.077$ & $2.624 \pm 0.088$ & $1.22$ & $+6.3\%$ \\
\midrule
\multirow{3}{*}{$2\times$}
& A1  & $2.990 \pm 0.075$ & $1.878 \pm 0.062$ & $1.59$ & \\
& A6  & $2.779 \pm 0.117$ & $2.273 \pm 0.110$ & $1.22$ & $+7.0\%$ \\
& A9  & $2.908 \pm 0.121$ & $2.369 \pm 0.128$ & $1.23$ & $+2.7\%$ \\
\midrule
\multirow{3}{*}{$4\times$}
& A1  & $2.733 \pm 0.070$ & $1.480 \pm 0.091$ & $1.85$ & \\
& A6  & $3.089 \pm 0.265$ & $2.577 \pm 0.217$ & $1.20$ & $-13.1\%$ \\
& A9  & $3.210 \pm 0.283$ & $2.660 \pm 0.250$ & $1.21$ & $-17.5\%$ \\
\midrule
\multirow{3}{*}{$8\times$}
& A1  & $3.286 \pm 0.155$ & $1.594 \pm 0.210$ & $2.06$ & \\
& A6  & $4.905 \pm 0.620$ & $4.117 \pm 0.465$ & $1.19$ & $-49.3\%$ \\
& A9  & $4.939 \pm 0.629$ & $4.194 \pm 0.535$ & $1.18$ & $-50.3\%$ \\
\bottomrule
\end{tabular}
\end{table}

\paragraph{The inversion.} The RISP advantage on the primary endpoint
is $+8.3\%$ at $0.5\times$, peaks at $+10.3\%$ at the headline $1\times$,
narrows to $+7.0\%$ at $2\times$, then \emph{inverts}: $-13.1\%$ at
$4\times$ and $-49.3\%$ at $8\times$. The crossover sits between $2\times$
and $4\times$, i.e.\ around $2$--$3\times$ the headline SNR.

\paragraph{The relearning-penalty ratio isolates the mechanism.} The
post/steady ratio $\rho$ tells the mechanistic story with unusual clarity.
A6's ratio is \emph{flat} across the entire sweep:
$1.20, 1.23, 1.22, 1.20, 1.19$. The retention mechanism works identically
at every SNR---a pinned specialist's reactivation penalty is always the
same modest $\sim\!20\%$ premium over its own steady state. A1's ratio, by
contrast, climbs monotonically: $1.28, 1.39, 1.59, 1.85, 2.06$. The
monolith's \emph{relative} relearning penalty gets worse with SNR, not
better. So why does the monolith win at high SNR? Because its
\emph{absolute levels} collapse: A1's steady regret falls from $2.96$ at
$0.5\times$ to $1.48$ at $4\times$, while A6's steady regret
\emph{deteriorates} ($2.89 \to 2.58 \to 4.12$ from $0.5\times$ through
$8\times$, non-monotonically). At high SNR a pooled learner extracts the
strong signal from all data jointly and relearns a reactivated regime
within days---faster than the probe window can penalize it---whereas each
pinned specialist is confined to its own niche's fraction of the data and
carries a variance penalty calibrated for a noisier world. Pooling beats
partitioning when the signal is strong enough to relearn cheaply. (A1's
own non-monotonicity at $8\times$---post-react $3.29$ vs $2.73$ at
$4\times$, with steady CIs widening---suggests additional dynamics at
extreme SNR, e.g.\ aggressive-fit instability, that we have not dissected;
we note it rather than interpret it.)

\paragraph{The skyline inverts identically.} The decisive control: A9, the
\emph{oracle-assigned} pinned arm, inverts in lockstep with A6 ($-17.5\%$
at $4\times$, $-50.3\%$ at $8\times$; their CIs overlap heavily at both
points). The inversion therefore cannot be an artifact of RISP's
learned assignment, its pinning heuristic, or any estimation failure in
our mechanism---perfect assignment inverts too. The inversion is a property
of the \emph{regime}: partitioned, reward-independent memory is the right
architecture in low-SNR, slow-relearning environments and the wrong one in
high-SNR, fast-relearning environments. The boundary is part of the claim,
not a caveat appended to it.

\begin{keypoint}
\textbf{Deployment-boundary reading.} The audit gives an operational rule:
RISP-style partitioned retention should be deployed only where the
post-reactivation relearning cost of a pooled learner is demonstrably
large---empirically, below roughly $2$--$3\times$ the headline SNR in this
family, i.e.\ in environments where a monolith's post/steady ratio is
materially elevated ($\rho \gtrsim 1.4$). The monolith's measured $\rho$
is itself an observable diagnostic: it can be estimated on any stream
\emph{before} committing to the architecture. Financial daily data, where
the E0 diagnostic already found no exploitable conditional structure at
all, sits firmly outside the deployment region; the mechanism's natural
habitat is low-SNR streams with long-dormancy recurring structure.
\end{keypoint}

\begin{intuition}
A useful way to compress E6: \emph{retention is a hedge against expensive
relearning}. The price of the hedge is the partitioning cost (less data per
specialist, a variance penalty, a worse steady state). The value of the
hedge is the relearning debt it cancels. Low SNR makes relearning slow and
expensive, so the hedge is cheap relative to its payout; high SNR makes
relearning nearly free, so the hedge is all premium and no payout. The
post/steady ratios are precisely the two sides of this ledger, read off
the data.
\end{intuition}

% ============================================================================
\section{Threats to Validity}
\label{sec:dd5-threats}

We enumerate the limits of the evidence, in roughly decreasing order of
importance, with what each would take to overturn.

\begin{enumerate}
\item \textbf{The headline is synthetic, and the real-data result is null.}
This is the central limitation and we state it without hedging: every
positive number in Sections \ref{sec:dd5-e1}--\ref{sec:dd5-e5e6} comes from
a generator we wrote. E0/E1s show that the mechanism's precondition is
absent in the real universes we tested, so the contribution is a
\emph{mechanism, a boundary map, and a falsifiable deployment gate}---not a
demonstrated real-market improvement. A reader who needs the latter should
treat this paper as a protocol to run on their own data, starting with E0.
\item \textbf{Small real universes.} The E0/E1s null rests on five crypto
pairs and three commodities at daily frequency with one feature set and two
labelers. The null is well-measured on that data ($|z| < 0.7$ across all
four cells) but could fail to generalize to intraday frequencies, equities
cross-sections, or richer conditioning information. The gate is cheap to
rerun; we make no claim beyond the cells in Table~\ref{tab:dd5-e0}.
\item \textbf{Single decision-layer family.} All arms share one decision
layer (top-$K$ portfolio from predicted returns). The decision-space
dissociation literature \citep{elmachtoub2022smart,yan2026invpnco} cuts
both ways: just as prediction nulls did not settle the decision question,
our decision-layer-specific results do not settle other decision layers
(e.g.\ mean-variance with estimated covariances, or execution-aware
objectives). The within-family rankings are clean; the family is one
point in a larger space.
\item \textbf{Linear specialists.} Specialists are linear in the $d = 20$
features. Retention-vs-relearning economics could shift with
high-capacity nonlinear specialists, in either direction (faster relearning
hurts retention's value, as E6 shows; worse interference helps it). The
GAUSE neural-expert results \citep{li2026gause} suggest the
reward-independence principle transfers to gradient-trained experts, but
we have not verified that here.
\item \textbf{Task-incremental labels are given.} Arms receive the regime
label stream (synthetic: true labels; E1s: causal labeler output). The
harder class-incremental problem---detecting reactivations without
labels---is not addressed, and the GAUSE real-data experience (label-free
recovery collapsing under feature overlap) warns against assuming it is
easy.
\item \textbf{Hyperparameter choices.} Each load-bearing choice is ablated
where it matters: $\beta$ has a $16\times$ flat plateau
(Table~\ref{tab:dd5-e5-beta}), $W_c$ is flat over $60\times$
(Table~\ref{tab:dd5-e5-wc}), $\lambda$ is nearly irrelevant, pinning
matters and is the mechanism itself. The probe length ($15$ days) and
dormancy threshold ($90$ days) were fixed before the final run and the E3
sweep shows the effect is not knife-edged in the dormancy axis; we did not
sweep probe length, and the post/steady split depends on it mechanically.
\item \textbf{Seed count and inference.} $n = 20$ seeds powers the
pre-registered effects (Section~\ref{sec:dd5-ci}) but leaves the marginal
comparison A6 vs A9 unresolved (Holm $p = 0.057$) and gives only moderate
resolution on variance signatures (A4). The conservative unpaired tests
mean reported significances are, if anything, understated.
\item \textbf{What would change our conclusions.} Concretely: (i) an E0 run
on some real universe showing a robust positive gap, followed by an E1s-style
run in which arms \emph{fail} to separate, would break the protocol's
central link; (ii) a demonstration that A2-style routers with explicit
idle-expert protection match A6 would collapse the
reward-independence claim to an implementation detail; (iii) an E6-style
audit on a low-SNR regime showing the monolith's $\rho$ elevated but no
RISP advantage would falsify the deployment gate. All three are
runnable with the released code, and the third is the one we would run
first.
\end{enumerate}

\begin{keypoint}
\textbf{Summary of Part V.} On a generator with the hypothesized structure:
retention is worth $6.1\%$, invariance $4.4\%$, jointly $10.3\%$
($38.1\%$ of excess-over-floor), with a super-additive interaction
($I = -0.0015 \pm 0.0005$) and a reached skyline (A6 $\approx$ A10,
$p = 0.72$); the effect requires capacity pressure (gap $\to 0.007$ at
$K = R$, hard memory), grows then saturates with dormancy, scales with
heterogeneity (gap $0.0010 \to 0.0024$), and survives every ablation with
flat plateaus. On real daily crypto/commodities the precondition is absent
(all E0 $|z| < 0.7$) and the arms correctly fail to separate (all
E1s $p > 0.2$). At high SNR the advantage inverts ($-49.3\%$ at $8\times$),
the oracle skyline inverts with it, and the monolith's own post/steady
ratio provides a measurable deployment gate. The mechanism is real,
bounded, and honest about where it lives.
\end{keypoint}
% ============================================================================
% PART VI: IMPLEMENTATION WALKTHROUGH
% LaTeX fragment -- no preamble; relies on master preamble (lstset python
% style, tcolorboxes intuition/keypoint/warning, theorem family, \R \E \Var \nm)
% ============================================================================

\part{Implementation Walkthrough}
\label{sec:dd6-part}

The preceding parts developed the theory at the blackboard. This part walks
the actual code. The entire system---decision layer, synthetic market,
schedules, capacity-bounded specialists, all nine allocation arms, the run
loop, and the statistics---is roughly 900 lines of seeded \texttt{numpy}
across three files: \texttt{risp.py} (decision layer, market, schedules,
\texttt{Specialist}, arms \texttt{A1}--\texttt{A10}, \texttt{run\_arm},
statistics), \texttt{realdata.py} (loaders, causal labelers, features, block
library, stitched market), and \texttt{run\_experiments.py} (drivers
\texttt{e0}--\texttt{e6}).
There is no GPU, no deep-learning framework, and no external optimizer. The
full experimental battery reruns from scratch in under an hour on a laptop,
and every number quoted in the paper traces to a specific JSON file written
by a specific driver function (Section~\ref{sec:dd6-repro}).

\begin{keypoint}
Smallness is a feature, not a budget constraint. Because the predictors are
linear, the decision oracle is a sort, and the surrogate is SPO+, every
mechanism claimed in the theory---interference under capacity, dormancy
starvation, episode-spurious overfitting, reward-decoupled niche
assignment---is visible in the code as a handful of lines, and nothing else
is present that could secretly produce the effects.
\end{keypoint}


% ============================================================================
\section{The Decision Layer in Code}
\label{sec:dd6-decision}

The decision set is the cardinality-constrained long-only portfolio
$\mathcal{Z} = \{ z \in [0, w_{\max}]^n : \mathbf{1}^\top z \le W,\;
\nm{z}_0 \le k \}$ with linear utility $F(z, y) = y^\top z$. With the budget
non-binding ($W \ge k\, w_{\max}$), the exact maximizer puts $w_{\max}$ on
the $k$ largest \emph{positive} coordinates of the score vector. The whole
decision layer is three functions:

\begin{lstlisting}[caption={The decision layer: exact oracle, regret, and the SPO+ subgradient (\texttt{risp.py}).},label={lst:dd6-decision}]
def solve_topk(y_hat: np.ndarray, k: int, w_max: float) -> np.ndarray:
    """Exact solution of max_z y_hat^T z over Z. O(n log n)."""
    n = y_hat.shape[0]
    z = np.zeros(n)
    idx = np.argsort(-y_hat)[:k]
    pos = idx[y_hat[idx] > 0.0]
    z[pos] = w_max
    return z


def regret(y_hat: np.ndarray, y: np.ndarray, k: int, w_max: float) -> float:
    """Decision regret rho(y_hat; y) = F(z*(y), y) - F(z(y_hat), y) >= 0."""
    z_star = solve_topk(y, k, w_max)
    z_hat = solve_topk(y_hat, k, w_max)
    return float(y @ z_star - y @ z_hat)


def spo_plus_loss_grad(w: np.ndarray, X: np.ndarray, y: np.ndarray,
                       k: int, w_max: float):
    """SPO+ loss and subgradient for the linear predictor y_hat = X w.

    Cast as minimization with cost c = -y, c_hat = -X w.
    l_SPO+(c_hat, c) = max_z (c - 2 c_hat)^T z + 2 c_hat^T z*(c) - c^T z*(c)
    d l / d c_hat = 2 (z*(c) - z*(2 c_hat - c));  d c_hat / d w = -X.
    """
    y_hat = X @ w
    z_star = solve_topk(y, k, w_max)            # argmin c^T z = argmax y^T z
    # z*(2 c_hat - c) = argmin (2 c_hat - c)^T z = argmax (2 y_hat - y)^T z
    z_alt = solve_topk(2.0 * y_hat - y, k, w_max)
    # loss value (for monitoring / variance penalty)
    loss = float((2.0 * y_hat - y) @ z_alt - 2.0 * y_hat @ z_star + y @ z_star)
    grad_chat = 2.0 * (z_star - z_alt)
    grad_w = -X.T @ grad_chat
    return loss, grad_w
\end{lstlisting}

\paragraph{Line-by-line.}
\begin{itemize}
  \item \texttt{solve\_topk}: \texttt{np.argsort(-y\_hat)[:k]} selects the
        $k$ largest scores; the filter \texttt{y\_hat[idx] > 0.0} drops any
        of them that are negative (a long-only book never holds an asset it
        expects to lose money on); the survivors receive weight
        $w_{\max}$. This is the exact $\arg\max_{z \in \mathcal{Z}}
        \hat{y}^\top z$, in $O(n \log n)$.
  \item \texttt{regret}: two oracle calls---one on the realized $y$ (the
        hindsight-optimal book $z^*(y)$), one on the prediction $\hat y$ (the
        served book)---and the difference of realized utilities. This is the
        SPO decision regret $\rho(\hat y; y) \ge 0$ used as the metric
        \emph{everywhere}: it is the per-day loss in \texttt{run\_arm}, the
        gate signal in the router, the competition score in RISP, and
        the weight signal in both Hedge arms.
  \item \texttt{spo\_plus\_loss\_grad}: the minimization-form bookkeeping.
        SPO+ is stated for cost \emph{minimization}, so we set $c = -y$ and
        $\hat c = -Xw$. The loss
        \[
          \ell_{\mathrm{SPO+}}(\hat c, c)
          = \max_{z \in \mathcal{Z}} (c - 2\hat c)^\top z
            + 2 \hat c^\top z^*(c) - c^\top z^*(c)
        \]
        translates, under the sign flip, into exactly the line computing
        \texttt{loss}: the inner max over $(c - 2\hat c)$ becomes a max over
        $(2\hat y - y)$ (the comment on \texttt{z\_alt} spells this out), and
        the remaining two terms become $-2\hat y^\top z^* + y^\top z^*$.
  \item The subgradient requires \emph{two oracle calls}: one at the true
        cost ($z^*(c)$, i.e.\ \texttt{z\_star}) and one at the reflected
        cost $2\hat c - c$ (i.e.\ \texttt{z\_alt}). Then
        $\partial \ell / \partial \hat c = 2(z^*(c) - z^*(2\hat c - c))$
        is \texttt{grad\_chat}, and the chain rule through
        $\hat c = -Xw$ contributes the $-X^\top$ in \texttt{grad\_w}. Note
        the two sign flips cancel into readable code: a head that
        under-ranks an asset the hindsight book holds gets pushed
        \emph{toward} it.
\end{itemize}

\begin{keypoint}
\textbf{No Gurobi needed.} Because the feasible set is a cardinality
polytope with a non-binding budget and the objective is linear, the oracle
is exact and costs one sort. Both the loss and the subgradient of SPO+ only
ever touch the optimization problem through the oracle
$z^*(\cdot)$---\texttt{solve\_topk} is called twice per gradient, never
differentiated through. If the decision problem were upgraded (transaction
costs, sector constraints, integer lots), the \emph{only} change is to swap
\texttt{solve\_topk} for a general solver call; \texttt{regret} and
\texttt{spo\_plus\_loss\_grad} are already written against the oracle
interface and would not change by a character.
\end{keypoint}


% ============================================================================
\section{Markets and Schedules}
\label{sec:dd6-markets}

\subsection{The synthetic generator}
\label{sec:dd6-synth}

The synthetic market implements the structural model of Part~II: per-regime
invariant loadings $\theta_r$, per-episode spurious loadings
$\gamma_{r,e}$, distractor features, and regime-dependent noise.

\begin{lstlisting}[caption={\texttt{SynthConfig} and the day generator (\texttt{risp.py}).},label={lst:dd6-synth}]
@dataclass
class SynthConfig:
    n_assets: int = 20
    d: int = 20                     # x1,x2 invariant; x3 spurious; 16 distractors; const
    k: int = 4
    w_max: float = 0.25
    R: int = R_DEFAULT
    noise: float = 0.02
    het: float = 1.0                # episode heterogeneity scale (E4 sweep)
    gamma_scale: float = 0.012
    theta_scale: float = 0.012     # |theta_r| ~ realistic daily alpha scale
    vol_mult: tuple = (1.0, 0.75, 2.0, 3.0)  # per-regime noise multiplier
    snr_mult: float = 1.0           # audit sweep: scales signal vs noise

    # ... inside SyntheticMarket ...
    def day(self, r: int, e: int):
        """One day's (X, y) in regime r, episode e."""
        c = self.cfg
        X = np.ones((c.n_assets, c.d))
        X[:, :c.d - 1] = self.rng.normal(0, 1, (c.n_assets, c.d - 1))
        g = self.gamma_for(r, e)
        mean = X[:, :2] @ self.theta[r] + g * X[:, 2] + self.b[r]
        y = mean + self.rng.normal(0, c.noise * c.vol_mult[r], c.n_assets)
        return X, y
\end{lstlisting}

Each configuration field maps onto a quantity from the theory:
\begin{itemize}
  \item \texttt{theta\_scale} and the constructor loop (not shown) place the
        invariant loadings $\theta_r = \texttt{theta\_scale} \cdot
        (\cos(2\pi r/R + 0.4), \sin(2\pi r/R + 0.4))$ on well-separated
        directions of the unit circle. This is what makes cross-regime
        interference \emph{real}: a head dragged toward regime $q$'s optimum
        genuinely loses regime $r$'s.
  \item \texttt{gamma\_scale} and \texttt{het} control the episode weather
        $\gamma_{r,e} \sim \mathcal{N}(0, (\texttt{het} \cdot
        \texttt{gamma\_scale})^2)$, drawn lazily and memoized in
        \texttt{gamma\_for} so that revisiting episode $e$ of regime $r$
        replays the same weather. The spurious feature is column
        \texttt{X[:, 2]} in the mean: an ERM learner profitably loads on it
        within an episode, then pays at the next episode when $\gamma$ is
        redrawn. \texttt{het} is the dial swept in E4.
  \item Columns $3$ through $d-2$ are pure distractors---they enter $X$ but
        never the mean---so the predictor has plenty of rope to overfit
        noise, and column $d-1$ is the constant absorbing the regime
        intercept \texttt{b[r]}.
  \item \texttt{vol\_mult} multiplies the noise s.d.\ per regime
        ($3{\times}$ in crisis), so the hard regimes are hard both because
        they are rare \emph{and} because they are noisy.
  \item \texttt{snr\_mult} scales $\theta_r$ and $\gamma_{r,e}$ jointly
        against fixed noise; it is the dial for the E6 audit (at what SNR
        does retention stop mattering because relearning is instant?).
\end{itemize}

\subsection{Schedules and reactivation bookkeeping}
\label{sec:dd6-schedules}

A \texttt{Schedule} is four aligned arrays: the regime id per day, the
episode index per day, a \texttt{block\_start} flag, and---only meaningful
on block starts---the \texttt{dormancy} length that preceded the block.

\begin{lstlisting}[caption={Schedule builders and dormancy bookkeeping (\texttt{risp.py}).},label={lst:dd6-sched}]
def make_schedule(rng: np.random.Generator, R: int = R_DEFAULT,
                  n_cycles: int = 18, block_len=(25, 60),
                  crisis_every: int = 4, crisis_len=(15, 30)) -> Schedule:
    """Common regimes (0,1,2) alternate in blocks; crisis (R-1) appears only
    every `crisis_every` cycles -> long dormancy, recurring reactivations."""
    seq = []
    cyc = 0
    common = [0, 1, 2] if R >= 4 else list(range(R - 1))
    while cyc < n_cycles:
        order = rng.permutation(common)
        for r in order:
            L = int(rng.integers(block_len[0], block_len[1] + 1))
            seq.append((int(r), L))
        cyc += 1
        if cyc % crisis_every == 0:
            L = int(rng.integers(crisis_len[0], crisis_len[1] + 1))
            seq.append((R - 1, L))
    return _seq_to_schedule(seq, R)


def make_dormancy_schedule(rng: np.random.Generator, D: int,
                           R: int = R_DEFAULT, n_react: int = 8,
                           block_len=(25, 60), crisis_len: int = 20) -> Schedule:
    """E3: focal regime R-1 reactivates every ~D days of dormancy exactly."""
    seq = []
    common = [0, 1, 2] if R >= 4 else list(range(R - 1))
    # initial training occurrences of focal regime
    for r in [R - 1, 0, R - 1, 1, R - 1, 2]:
        L = crisis_len if r == R - 1 else int(rng.integers(*block_len))
        seq.append((int(r), L))
    for _ in range(n_react):
        filler = 0
        while filler < D:
            r = int(rng.choice(common))
            L = int(rng.integers(block_len[0], block_len[1] + 1))
            seq.append((r, L))
            filler += L
        seq.append((R - 1, crisis_len))
    return _seq_to_schedule(seq, R)
\end{lstlisting}

Both builders emit a list of \texttt{(regime, length)} pairs and hand it to
\texttt{\_seq\_to\_schedule}, which does the bookkeeping that the
reactivation metric depends on: for each block it increments the per-regime
episode counter, sets \texttt{block\_start[t]} on the first day, and records
\texttt{dormancy[t] = t - last\_seen[r]}, where \texttt{last\_seen[r]} is
the \emph{final day} of the regime's previous block. The method
\begin{lstlisting}
    def reactivation_days(self, min_dormancy: int):
        """Days that begin a block whose regime was dormant >= min_dormancy."""
        out = []
        for t in range(self.T):
            if self.block_start[t] and self.dormancy[t] >= min_dormancy:
                out.append(t)
        return out
\end{lstlisting}
is then the single source of truth for which days count as reactivations:
\texttt{run\_arm} masks the first \texttt{probe} (default 15) days after
each such $t$ as the post-reactivation window, and only windows in the
second half of the run ($t \ge T/2$) enter the headline metric, so warm-up
episodes never contaminate it. \texttt{make\_schedule} gives irregular
dormancy by rationing crisis to every fourth cycle;
\texttt{make\_dormancy\_schedule} gives \emph{controlled} dormancy of
$\approx D$ days for E3, with three training occurrences of the focal
regime up front so there is expertise to forget.

\subsection{Causal labels on real data}
\label{sec:dd6-labels}

On real data the regimes are not given; they are computed by labelers that
must not leak day-$t$ information into the day-$t$ label, because the label
selects which head serves the day-$t$ decision. The L1 labeler is worth
quoting in full:

\begin{lstlisting}[caption={The L1 causal labeler (\texttt{realdata.py}).},label={lst:dd6-l1}]
def label_L1(panel: pd.DataFrame, vol_win=20, trend_win=50,
             vol_pct_win=252) -> np.ndarray:
    """Vol band (rolling-percentile of trailing vol) x trend sign. Causal:
    every input is shifted one day so day t uses data through t-1."""
    px = panel.mean(axis=1) if panel.shape[1] > 1 else panel.iloc[:, 0]
    # equal-weight index in log space to avoid scale domination
    px = np.log(panel).mean(axis=1)
    ret = px.diff()
    vol = ret.rolling(vol_win).std().shift(1)
    vol_rank = vol.rolling(vol_pct_win, min_periods=60).rank(pct=True)
    trend = (px.shift(1) - px.shift(trend_win + 1))
    high_vol = (vol_rank > 0.7).astype(int)
    up = (trend > 0).astype(int)
    lab = high_vol * 2 + (1 - up)        # 0 calm-up 1 calm-down 2 vol-up 3 vol-down
    return lab.values
\end{lstlisting}

The leakage prevention is the \texttt{.shift(1)} discipline, and each
occurrence matters:
\begin{itemize}
  \item \texttt{vol = ret.rolling(vol\_win).std().shift(1)}: the rolling
        window ending at $t$ would include the day-$t$ return $r_t$---which
        is the prediction target. The \texttt{shift(1)} moves the window end
        to $t-1$, so the volatility used to label day $t$ is the s.d.\ of
        $r_{t-20}, \dots, r_{t-1}$.
  \item \texttt{vol\_rank}: the percentile rank is computed over the
        already-shifted \texttt{vol} series, so the trailing-year band
        threshold ($>0.7$) also sees nothing past $t-1$. Ranking rather
        than thresholding in absolute units keeps the band meaningful across
        the secular volatility decay of the crypto sample.
  \item \texttt{trend = px.shift(1) - px.shift(trend\_win + 1)}: the 50-day
        trend ends at the close of $t-1$ (note the matching $+1$ on
        \emph{both} shifts: the window is $[t-51, t-1]$, length 50, never
        touching $t$).
\end{itemize}
The same convention governs the features in \texttt{build\_xy}
(\texttt{mom5} $=$ \texttt{lp.shift(1) - lp.shift(6)}, \texttt{vol20} $=$
\texttt{r1.rolling(20).std().shift(1)}, etc.): all features are functions of
information through $t-1$, while the target \texttt{y = r1} is the day-$t$
return. The L2 labeler is causal by construction rather than by shifting: it
is a forward-only two-state Hamilton filter whose update at day $t$ uses
\texttt{ret[t - 1]} (``yesterday's return only,'' as the inline comment
says), crossed with the same shifted trend sign.

\begin{intuition}
The block-shuffle control in E0 (\texttt{block\_shuffle} in
\texttt{run\_experiments.py}) is the converse check: it keeps the block
\emph{structure} of the labels but randomizes block \emph{identities}. If
regime-conditioned models beat block-shuffled ones on decision regret, the
labels carry real conditional structure, not just autocorrelation.
\end{intuition}


% ============================================================================
\section{The Specialist and Memory Models}
\label{sec:dd6-specialist}

A \texttt{Specialist} is a dictionary of at most $K$ linear heads (one per
held regime), an episode-indexed replay buffer per held regime, and a
training rule (ERM or invariance). Memory pressure is implemented in
\texttt{\_touch}:

\begin{lstlisting}[caption={Memory bookkeeping: hard LRU and soft interference (\texttt{risp.py}).},label={lst:dd6-touch}]
    def _touch(self, r: int):
        self.t += 1
        if r not in self.heads:
            if len(self.heads) >= self.K:
                lru = min(self.heads, key=lambda q: self.heads[q].last_used)
                del self.heads[lru]
                self.buf.pop(lru, None)
            self.heads[r] = Head(w=np.zeros(self.d))
            self.buf[r] = {}
        self.heads[r].last_used = self.t
        if self.memory == "soft":
            for q, h in self.heads.items():
                if q != r:
                    h.strength *= (1.0 - self.rho_int)
                    h.w *= (1.0 - self.rho_int)
            self.heads[r].strength = 1.0
\end{lstlisting}

\paragraph{Hard model.} When a new regime arrives at a full specialist
($\texttt{len(self.heads)} \ge K$), the least-recently-used head is deleted
\emph{together with its buffer}: \texttt{self.buf.pop(lru, None)}. This is
the line that makes forgetting total in the hard model---knowledge lost is
data lost, so the evicted regime must be relearned from scratch at its next
reactivation, with no replay shortcut. The fresh head starts at
$w = 0$, which under \texttt{solve\_topk} means a near-random book (the
\texttt{predict} fallback for unheld regimes is an explicit
$10^{-6}$-scale random tiebreaker for the same reason: zero scores would
deterministically pick no assets).

\paragraph{Soft model.} Under \texttt{memory == "soft"}, every training day
on regime $r$ multiplies every \emph{other} held head's weights and strength
by $(1 - \rho_{\mathrm{int}})$, with $\rho_{\mathrm{int}} = 0.06$ by
default. Untouched expertise decays geometrically toward zero at a rate set
by how often its owner trains on something else---the continuous analogue of
LRU eviction, matching the interference operator of Part~II. The active
head's strength snaps back to $1$.

Training itself lives in \texttt{learn}:

\begin{lstlisting}[caption={ERM vs.\ invariance training over the episode buffer (\texttt{risp.py}).},label={lst:dd6-learn}]
    def learn(self, X, y, r: int, e: int, n_steps: int = 2):
        """One observation: store, touch memory, take SGD steps on the
        regime-r objective over the episode buffer."""
        self._touch(r)
        self.store(X, y, r, e)
        h = self.heads[r]
        eps = self.buf[r]
        for _ in range(n_steps):
            if self.mode == "erm" or len(eps) < 2:
                # plain ERM on pooled buffer
                keys = list(eps.keys())
                ekey = keys[int(self.rng.integers(len(keys)))]
                Xb, yb = ep_sample(eps[ekey], self.rng)
                _, g = spo_plus_loss_grad(h.w, Xb, yb, self.k, self.w_max)
                h.w -= self.lr * g / max(1, Xb.shape[0])
            else:
                # invariance: Var_e[lbar_e] + beta * E_e[lbar_e]
                losses, grads = [], []
                for ekey in eps:
                    Xb, yb = ep_sample(eps[ekey], self.rng)
                    l, g = spo_plus_loss_grad(h.w, Xb, yb, self.k, self.w_max)
                    losses.append(l / max(1, Xb.shape[0]))
                    grads.append(g / max(1, Xb.shape[0]))
                E = len(losses)
                lbar = float(np.mean(losses))
                g_tot = np.zeros(self.d)
                for l, g in zip(losses, grads):
                    g_tot += (2.0 * (l - lbar) + self.beta) * g / E
                h.w -= self.lr * g_tot
\end{lstlisting}

\paragraph{Line-by-line.}
\begin{itemize}
  \item The buffer \texttt{self.buf[r]} maps episode index $e$ to a list of
        stored $(X, y)$ days, capped at \texttt{buf\_days}~$=40$ days per
        episode and \texttt{buf\_episodes}~$=6$ episodes per regime
        (\texttt{store} evicts oldest-first on both axes). Episode identity
        is the unit the invariance objective averages over.
  \item \emph{ERM branch}: sample one episode uniformly at random, draw a
        minibatch of up to 8 stored days from it (\texttt{ep\_sample}), take
        one SPO+ step. Over many steps this pools episodes---exactly the
        estimator that cannot distinguish $\theta_r$ from the episode
        weather $\gamma_{r,e}$.
  \item \emph{INV branch}: evaluate the per-episode mean losses $\bar\ell_e$
        and mean gradients $g_e$ on a fresh minibatch from \emph{every}
        buffered episode, then assemble the gradient of
        $\Var_e[\bar\ell_e] + \beta\, \E_e[\bar\ell_e]$. The identity used
        is
        \[
          \nabla_w \Var_e[\bar\ell_e]
          = \frac{2}{E} \sum_e (\bar\ell_e - \bar\ell)\,
            \big(g_e - \bar g\big)
          = \frac{2}{E} \sum_e (\bar\ell_e - \bar\ell)\, g_e,
        \]
        where the $\bar g$ term vanishes because
        $\sum_e (\bar\ell_e - \bar\ell) = 0$; adding
        $\nabla_w \beta \E_e[\bar\ell_e] = \frac{\beta}{E}\sum_e g_e$ gives
        \[
          \nabla_w \big[\Var_e + \beta \E_e\big]
          = \sum_e \frac{\big(2(\bar\ell_e - \bar\ell) + \beta\big)\, g_e}{E},
        \]
        which is verbatim the accumulation line
        \texttt{g\_tot += (2.0 * (l - lbar) + self.beta) * g / E}. Episodes
        whose loss is above the cross-episode mean get their gradients
        up-weighted; episodes below the mean can have \emph{negative}
        effective weight when $\bar\ell - \bar\ell_e > \beta/2$, actively
        pushing the head away from solutions that win on one episode's
        weather at others' expense.
  \item The guard \texttt{len(eps) < 2} falls back to ERM until a second
        episode exists---a one-episode variance is identically zero, so the
        INV objective is undefined before then.
\end{itemize}

\begin{keypoint}
Eviction travels with the buffer. A head that loses the LRU competition
loses its replay data in the same statement, so no arm can cheat the
capacity constraint by quietly relearning from a retained buffer. Capacity
$K$ bounds \emph{retained regimes}, in both parameters and data.
\end{keypoint}


% ============================================================================
\section{The Arms}
\label{sec:dd6-arms}

All arms share one interface: \texttt{decide(X, r) -> y\_hat} produces the
prediction that is converted into the served book, and
\texttt{observe(X, y, r, e, served\_regret)} performs all internal updates.
Signature-level summary:

\begin{itemize}
  \item \texttt{MonolithArm} (A1 ERM / A7 INV): one capacity-$K$ specialist
        that always learns the active regime; \texttt{observe} is a single
        \texttt{self.s.learn(X, y, r, e)}. This is the canonical
        reward-driven allocation---capacity follows activity, so dormant
        regimes are first in the LRU queue.
  \item \texttt{RecentPerfArm} (A3): the industry baseline. Capital shares
        are a softmax (temperature $0.02$) of trailing 60-day negative
        regret per specialist; the served decision is the top-capital
        specialist's; in \texttt{observe}, \emph{every} specialist trains on
        the active regime with probability equal to its capital share, so
        starved specialists drift into the active regime at the
        capital-floor rate.
  \item \texttt{FixedNicheArm} (A4 random / A9, A10 oracle): a
        reward-independent owner map \texttt{self.owner[r]} fixed at
        construction (random claims, or hand-pinned $r \mapsto r \bmod N$
        for the oracle skyline); \texttt{observe} trains the owner only.
  \item \texttt{HedgeFixedArm} (A8a): exponential weights over $2(d-1)$
        fixed $\pm$ unit-direction strategies, weights decayed by realized
        regret; no learning. \texttt{HedgeLearnersArm} (A8b): Hedge over $N$
        learning specialists, where the current top-weight learner trains on
        the active regime---training follows the weight, so it is
        reward-coupled like the monolith, just with a softer selector.
  \item \texttt{RouterArm} (A2) and \texttt{RISPArm} (A5 ERM / A6 INV):
        quoted in full below; they are the two poles of the paper's central
        contrast.
\end{itemize}

\subsection{RouterArm: where reward-coupling enters}
\label{sec:dd6-router}

\begin{lstlisting}[caption={\texttt{RouterArm.observe}: the reward-driven gate (\texttt{risp.py}).},label={lst:dd6-router}]
    def observe(self, X, y, r, e, served_regret):
        j = self._route(r, explore=True)
        ex = self.experts[j]
        rg = regret(ex.predict(X, r), y, ex.k, ex.w_max)
        b = self.base.get(r, rg)
        self.base[r] = 0.95 * b + 0.05 * rg
        self.G[r, j] += self.gate_lr * (b - rg)      # reward-driven gate
        ex.learn(X, y, r, e)
\end{lstlisting}

Each day, the gate $\varepsilon$-greedily routes the active regime to expert
$j$ (10\% exploration), scores that expert's regret, and updates a running
per-regime baseline $b_r$ by EWMA. The annotated line is where
reward-coupling \emph{enters} the allocation: the gate logit
\texttt{G[r, j]} moves by \texttt{gate\_lr * (b - rg)}---an advantage
update, positive when the routed expert beat the regime's recent baseline.
Crucially, the update is indexed by the \emph{active} regime $r$: a dormant
regime emits no gate gradient at all (Observation~1 of Part~III), so its
routing column \texttt{G[r, :]} is frozen while the experts it points to
are being retrained---and possibly LRU-evicted---by the regimes that are
active. The routed expert then trains on the day
(\texttt{ex.learn}), which is the second coupling: training mass also
follows activity.

\subsection{RISPArm: where reward-coupling exits}
\label{sec:dd6-risp}

\begin{lstlisting}[caption={\texttt{RISPArm.observe}: window accumulation and the pinning short-circuit (\texttt{risp.py}).},label={lst:dd6-nmobserve}]
    def observe(self, X, y, r, e, served_regret):
        if r in self.pinned:
            self.specs[self.pinned[r]].learn(X, y, r, e)
            return
        if self.win_regime is not None and r != self.win_regime:
            self._close_window()
        self.win_regime = r
        for i, sp in enumerate(self.specs):
            self.win_quality[i] += regret(sp.predict(X, r), y,
                                          self.k, self.w_max)
        self.win_data.append((X, y, e))
        self.win_days += 1
        if self.win_days >= self.Wc:
            self._close_window()
\end{lstlisting}

The first two lines are where reward-coupling \emph{exits}: once a regime is
pinned, \texttt{observe} routes the day's data to the pinned owner and
\texttt{return}s. No competition, no affinity update, no gate
gradient---the owner map for that regime is permanently
reward-independent from that day on, exactly the fixed-assignment regime of
the theory. Until pinning, days accumulate into a competition window: every
specialist is scored on the day (counterfactually---only the owner's
prediction was \emph{served}), the day's data is parked in
\texttt{win\_data}, and the window closes after $W_c = 20$ days or at a
regime switch. Window closing does the actual allocation work:

\begin{lstlisting}[caption={\texttt{RISPArm.\_close\_window}: EG update, winner-trains, pinning (\texttt{risp.py}).},label={lst:dd6-nmclose}]
    def _close_window(self):
        if self.win_days == 0 or self.win_regime is None:
            return
        r = self.win_regime
        score = -self.win_quality / self.win_days \
            + self.lam * (self.alpha[:, r] - 1.0 / self.R)
        i_star = int(np.argmax(score))
        # EG / Hedge update on the winner's affinity row
        a = self.alpha[i_star].copy()
        a[r] *= np.exp(self.eta)
        self.alpha[i_star] = a / a.sum()
        # winner trains on the window's data
        for (X, y, e) in self.win_data:
            self.specs[i_star].learn(X, y, r, e, n_steps=1)
        # pinning
        if self.pin_enabled and r not in self.pinned \
                and self.alpha[i_star, r] >= self.pin_thresh:
            self.pinned[r] = i_star
        self.win_quality = np.zeros(self.N)
        self.win_days = 0
        self.win_data = []
\end{lstlisting}

The winner is the specialist with the best (lowest) mean window regret plus
a niche bonus $\lambda(\alpha_{i,r} - 1/R)$ that rewards already-declared
affinity---the incumbency term that prevents ownership thrash. The winner's
affinity row gets a multiplicative EG/Hedge update
($\alpha_{i^*, r} \mathrel{*}= e^{\eta}$, then renormalize over regimes;
since rows are simplexes over regimes, boosting $r$ implicitly taxes
$i^*$'s claim on every other regime---winning here means withdrawing
elsewhere). \emph{Only the winner} trains, on the whole window's parked
data, with \texttt{n\_steps=1} per day to keep total gradient work
comparable to the streaming arms. When the winner's affinity crosses
\texttt{pin\_thresh}~$=0.95$, the regime is pinned and the
\texttt{observe} short-circuit takes over forever.

\begin{intuition}
Read the two \texttt{observe} methods side by side. The router updates its
allocation \emph{every day a regime is active and never otherwise}---it is
reward-coupled by construction, and dormancy starves it. RISP uses
realized regret only as a \emph{transient} signal to discover who should own
what; the moment ownership is settled, the signal is disconnected. The
allocation converges \emph{to} reward-independence. The E5(d) ablation
(\texttt{pinned} vs.\ \texttt{never}) isolates exactly the value of that
disconnection.
\end{intuition}


% ============================================================================
\section{Reproducibility Map}
\label{sec:dd6-repro}

Every quantitative claim in the paper traces to a JSON file in
\texttt{results/}, written by one driver in \texttt{run\_experiments.py}.

\begin{table}[h]
\centering
\small
\begin{tabular}{p{4.6cm} l l}
\hline
\textbf{Paper artifact} & \textbf{Driver} & \textbf{Results JSON} \\
\hline
Real-data structure diagnostic (regime labels vs.\ block-shuffled controls,
crypto + commodities) & \texttt{e0()} & \texttt{e0\_structure\_diagnostic.json} \\
Headline $2{\times}2$ dissociation table; Welch/Holm $p$-values; overall /
post-reactivation / steady-state regret, all arms & \texttt{e1()} &
\texttt{e1\_synth.json} \\
Stitched-crypto replication of the headline table & \texttt{e1s()} &
\texttt{e1s\_stitched\_crypto.json} \\
Capacity sweep ($K \in \{1,2,3,4\}$, hard vs.\ soft memory) & \texttt{e2()}
& \texttt{e2\_K\{K\}\_\{mem\}.json}, \texttt{e2\_capacity\_sweep.json} \\
Dormancy sweep ($D \in \{21, 63, 126, 252, 504\}$) & \texttt{e3()} &
\texttt{e3\_dormancy\_sweep.json} \\
Episode-heterogeneity sweep (\texttt{het} $\in \{0, 0.5, 1, 1.5, 2\}$, ERM
vs.\ INV at fixed allocation) & \texttt{e4()} &
\texttt{e4\_het\{h\}.json}, \texttt{e4\_heterogeneity\_sweep.json} \\
Ablations: $\beta$ sweep (incl.\ variance-only $\beta{=}0$), $W_c$ sweep,
$\lambda$ on/off, pinned vs.\ never-pinned & \texttt{e5()} &
\texttt{e5\_ablations.json} \\
SNR audit (break-even relearning speed) & \texttt{e6()} &
\texttt{e6\_snr\_audit.json} \\
\hline
\end{tabular}
\caption{Map from paper artifacts to experiment drivers and output files.}
\label{tab:dd6-repromap}
\end{table}

\paragraph{Seeding discipline.} Within a seed $s$, three independent RNG
streams are used, and they are constructed so that \emph{all arms see
identical data}:
\begin{itemize}
  \item \emph{Schedule stream}: \texttt{np.random.default\_rng(1000 + s)}
        builds the schedule once per seed (E3 uses $3000 + s$, E1s uses
        $2000 + s$), before any arm runs.
  \item \emph{Market stream}: the market is \emph{reconstructed per arm}
        with the same seed---\texttt{SyntheticMarket(cfg, seed=5000 + s)}
        inside the arm loop ($7000 + s$ in E3)---so every arm replays the
        identical sequence of $(X_t, y_t)$ days. Comparisons across arms
        within a seed are therefore paired, not merely matched in
        distribution.
  \item \emph{Arm stream}: each arm receives a fresh
        \texttt{np.random.default\_rng(99 * s + 7)} ($13s+5$ in E3,
        $55s+3$ in E1s). Multi-specialist arms then spawn one sub-RNG per
        specialist via \texttt{np.random.default\_rng(rng.integers(2**31))}
        in their constructors, so each specialist owns a private stream
        derived from the order of construction-time draws.
\end{itemize}

\paragraph{Rerunning.} From the \texttt{code/} directory:
\begin{lstlisting}[language=bash,numbers=none]
python run_experiments.py e1            # headline table, 20 seeds
python run_experiments.py e1 --seeds 5  # quick smoke run
python run_experiments.py e0            # real-data diagnostic (needs GAUSE/data)
for e in e1s e2 e3 e4 e5 e6; do python run_experiments.py $e; done
\end{lstlisting}
The only external dependencies are \texttt{numpy}, \texttt{scipy} (Welch
tests), and \texttt{pandas} (real-data loaders); \texttt{e0}/\texttt{e1s}
additionally require the Bybit and FRED CSVs under \texttt{GAUSE/data/}.
Each driver prints the absolute path of every JSON it saves.

\begin{warning}
\textbf{The one nondeterminism trap.} All randomness flows from explicit
seeds, but the \emph{construction-time draw order} is part of the
specification. Each arm's top-level RNG is freshly seeded
(\texttt{99*s + 7} for all arms in E1), so adding or removing arms from the
run list is safe; what is \emph{not} safe is changing what an arm does with
that RNG at construction. Every sub-specialist seed comes from a sequential
\texttt{rng.integers(2**31)} call in the constructor, so reordering the arm
factories' internals, changing $N$, or inserting an extra draw in any
\texttt{\_\_init\_\_} silently shifts every downstream stream of that
arm---exploration choices, minibatch sampling, tie-breaking
predictions---and the rerun will no longer bit-match an archived JSON even
though nothing statistically meaningful changed. When validating against
archived results, leave \texttt{ARM\_FACTORIES} and all arm constructors
untouched; when changing them deliberately, regenerate \emph{all} JSONs and
re-archive.
\end{warning}

% ============================================================================
% CONCLUSION
% ============================================================================
\part{Conclusion}

\section*{Summary of the Deep Dive}

The mathematics of \nm{} reduces to one decomposition and two
dichotomies. The decomposition: post-reactivation decision regret splits
into an invariance gap---mean plus variance of episode risks, controlled
by Cantelli's inequality and priced for finite episodes by Hoeffding's---
and a forgetting deficit---eviction probability times relearning cost.
The first dichotomy: the forgetting deficit is identically zero for
reward-independent covering assignments and tends exponentially to its
maximum for reward-driven ones, because a dormant regime emits no signal
with which a reward-driven allocator could protect it; this is a property
of \emph{classes} of mechanisms, confirmed experimentally as a class
boundary. The second dichotomy: the invariance gap is governed by the
training objective alone, is provably reduced by variance penalization
across episodes when episodes are heterogeneous draws of a stable
decision structure, and is \emph{inert}---measurably, predictably---when
the allocation axis has already discarded the head it trained. The
super-additivity of the two fixes is therefore not a curiosity but the
signature of the account: each mechanism is the other's precondition.

Around these results sit the honesty obligations this research program
treats as part of the mathematics: the plug-in caveat on every
generalization bound applied to a trained head; the explicit looseness of
the Pinsker bridge constant; the refusal to claim convergence theorems
for coupled winner-take-all dynamics that we can only demonstrate; the
estimation arithmetic showing that six crisis episodes cannot certify
crisis-grade invariance without environment enlargement; and the three
pre-registered expectations our own experiments refuted---reported here
with the same care as the confirmations, because a theory of honest
retention should itself retain what it got wrong.

\bibliographystyle{unsrtnat}
\bibliography{references}

\end{document}
